\pdfinfoomitdate=1
\pdftrailerid{}
\pdfsuppressptexinfo=15
\documentclass[11pt,letterpaper]{article}
\usepackage[margin=1in]{geometry}
\usepackage[T1]{fontenc}
\usepackage{lmodern,amsmath,amssymb,amsthm,mathtools,booktabs,microtype,longtable,array}
\usepackage[colorlinks=true,linkcolor=black,citecolor=black,urlcolor=blue]{hyperref}
\usepackage{xurl}
\usepackage{enumitem}
\setlist{itemsep=3pt,topsep=5pt}
\setlength{\parskip}{5pt}\setlength{\parindent}{0pt}
\setlength{\emergencystretch}{2em}
\allowdisplaybreaks[2]
\newtheorem{theorem}{Theorem}[section]
\newtheorem{lemma}[theorem]{Lemma}
\newtheorem{proposition}[theorem]{Proposition}
\newtheorem{corollary}[theorem]{Corollary}
\theoremstyle{remark}\newtheorem{remark}[theorem]{Remark}
\numberwithin{equation}{section}
% Macros of the three sources; every shared name has the same definition in each source.
\newcommand{\N}{\mathbb N}
\newcommand{\Prob}{\mathbb P}
\newcommand{\Q}{\mathbb Q}
\newcommand{\Cov}{\operatorname{Cov}}
\newcommand{\HH}{\mathcal H}
\newcommand{\dd}{\mathrm d}
\newcommand{\norm}[1]{\left\|#1\right\|}
\newcommand{\ip}[2]{\langle #1,#2\rangle}
% Added in the merge (batch 105).
\newcommand{\file}[1]{\nolinkurl{#1}}
\newenvironment{writenote}[1][]{\par\smallskip\begingroup\small
 \noindent\textbf{[write]}\if\relax\detokenize{#1}\relax\else~\textbf{#1.}\fi\ }%
 {\par\endgroup\smallskip}
\newcommand{\wtag}{\textup{\textbf{[write]}}}
\newenvironment{partabstract}{\begin{quote}\small\noindent\textbf{Abstract of this Part (as delivered).}\ }{\end{quote}}
\newcommand{\partsource}[1]{\par\noindent\begin{minipage}{\textwidth}\small #1\end{minipage}\par\medskip}
\hypersetup{pdftitle={The proportional placement game (OEIS A398540)},pdfauthor={},pdfsubject={Merged report a398540-proportional-placement-game (batch 105): the computable rate and the leading equivalent (Report 144), the pointwise logarithmic corrections (Report 145), every fixed order (Report 146)}}
\title{\textbf{The proportional placement game (OEIS A398540)}\\[7pt]
\large A computable rate, the leading equivalent,\\ the logarithmic corrections and every fixed order}
\author{Three research reports of one external research session\\(bundle Reports 144, 145 and 146), merged into one ProveIt report}
\date{Parts I--III: 3 October 2026\quad merged: 5 October 2026}
\begin{document}
\hypersetup{pageanchor=false}
\maketitle
\begin{abstract}
In the proportional placement game, $n$ independent uniform values on $(0,1)$ arrive one at a time and must be placed at once into $n$ ordered slots so that the placed values increase from left to right; the proportional policy puts a value into the slot of its containing gap that corresponds to its relative position there. Its winning probabilities $W_n$ satisfy an exact quadratic recurrence with beta-binomial weights, derived in Serra's manuscript on the game, and OEIS A398540 records digits of the limit of $W_{n+1}/W_n$. This report prints three manuscripts on $W_n$ as Parts~I--III, in the order in which each depends on the one before. Part~I proves that $r_*=\lim W_n^{1/n}$ exists and is computable, certifies $527/1000<r_*<79/125$ by exact rational arithmetic, and proves $W_n\sim r_*\sqrt{2\pi n}\,r_*^n$ and $W_{n+1}/W_n\to r_*$; the amplitude $r_*\sqrt{2\pi}$ corrects the value $\sqrt{2\pi}$ that Serra's heuristic integral equation gives. Part~II proves $W_n/(A\sqrt n\,r_*^n)=1+\frac1{4n}-\frac7{540}\frac{\log n}{n^2}+\frac{B_2}{n^2}+o(n^{-2})$, with $A=r_*\sqrt{2\pi}$ and $B_2$ given by absolutely convergent series, the eventual strict decrease of the ratios, and that the ratios have no pure inverse-power expansion through order $n^{-3}$. Part~III proves an expansion in powers of $1/n$ whose $j$th coefficient is a polynomial in $\log n$ of degree at most $j-1$, to every fixed order; the explicit third order, at which the squared logarithm cancels; the first logarithmic term $\frac7{270}\frac{\log n}{n^3}$ of $W_{n+1}/(r_*W_n)$; and computable nested enclosures of the rate with a proved contraction. Each Part also inverts its forward result for the first passage $N(\varepsilon)=\min\{n:W_n\le\varepsilon\}$, with a Lambert-$W_{-1}$ centre and a ceiling-safe bracket. The write adds proofs that Serra's prefactor statements and, read through the asymptotic profile, his integral equation~(200) are wrong, and that his Open Problem~2 fails for every $p\ge2$, and it records a numerical observation: the constant $S_A$ of Part~III agrees with $1/12$ to 18 digits (uncertified). Effective constants, all-index monotonicity of the ratios and the exact value of $S_A$ remain open. The report is unrefereed, and nothing in it is formalized.
\end{abstract}
\hypersetup{pageanchor=true}
\setcounter{tocdepth}{1}
{\small\setlength{\parskip}{0pt}\tableofcontents}
\newpage

\section*{Guide to this report}\phantomsection\label{ppg:sec:guide}
\addcontentsline{toc}{section}{Guide to this report}
The three Parts are three manuscripts of the session bundle of Reports 1--243 (arrival commit \texttt{60f54ea06}), placed in this directory by commit \texttt{e85586b7c} (batch~105). They treat one sequence, the winning probabilities $W_n$ of the proportional placement game, through one recurrence:
\begin{center}
\small
\begin{tabular}{@{}l l >{\raggedright\arraybackslash}p{0.47\textwidth} l@{}}
\toprule
Part & Bundle report & Delivered title & Labels\\
\midrule
I & 144 & Asymptotics and a computable exponential rate for a proportional placement game & \texttt{ppg:lead:}\\
II & 145 & Pointwise corrections for the proportional placement game & \texttt{ppg:log:}\\
III & 146 (base) & All fixed orders for the proportional placement game & \texttt{ppg:all:}\\
\bottomrule
\end{tabular}
\end{center}
\emph{Arrangement.} The base of the merge is Report~146, the strongest of the three: its \file{Report146.tex} was staged as this \file{article.tex} by the placement commit. It is printed last, as Part~\ref{ppg:all:part}, because it is not self-contained. It takes the leading equivalent, the kernel bounds, the radius facts and the certified rate bracket from Report~144, and the starting rate $x_n=1-\frac1{4n}+O((1+\log n)/n^2)$, the exact critical inverse and the series definition of $B_2$ from Report~145, without proving them again (its \eqref{ppg:all:eq:base}, \eqref{ppg:all:eq:oldsecond}, Section~\ref{ppg:all:sec:induction} and Section~\ref{ppg:all:sec:rate}). Report~145 in turn rests on Report~144's leading equivalent and kernel expansion. So all three are printed in full, in dependency order, which is also the order of writing. Reports~145 and~146 shipped unchanged copies of their predecessors in a \file{foundation/} directory; those copies are byte-identical to the standalone deliveries and are printed once, as Parts~\ref{ppg:lead:part} and~\ref{ppg:log:part}.

\emph{Numbering.} Sections are numbered continuously through the three Parts, and statements and equations are numbered within sections. Part~\ref{ppg:lead:part} keeps Report~144's section and statement numbers (Sections~1--12); its Appendix~A is Section~\ref{ppg:lead:app:reproduce}. In Part~\ref{ppg:log:part} a section number is Report~145's plus~$13$: its Theorem~1.1 is Theorem~\ref{ppg:log:thm:main}, its Theorem~8.6 and Corollary~8.7 are Theorem~\ref{ppg:log:thm:full} and Corollary~\ref{ppg:log:cor:monotonicity}, its Theorem~9.1 is Theorem~\ref{ppg:log:thm:ratio}, and its Theorem~10.1 is Theorem~\ref{ppg:log:thm:inverse}. In Part~\ref{ppg:all:part} the offset is~$25$: Report~146's Theorem~1.1 is Theorem~\ref{ppg:all:thm:all}, its Theorem~7.1 is Theorem~\ref{ppg:all:thm:third}, its Corollaries~8.1 and~8.2 are Corollaries~\ref{ppg:all:cor:ratio} and~\ref{ppg:all:cor:extrapolation}, its Theorem~9.1 is Theorem~\ref{ppg:all:thm:threshold}, and its Theorem~10.1 is Theorem~\ref{ppg:all:thm:rate}. The manuscripts number their equations consecutively through the whole document; here equations are numbered within sections, so a delivered equation number $(m)$ is not kept, but every delivered equation keeps its label name, with the Part's prefix. The delivered texts cite each other as ``Report~144'' and ``Report~145''; read Part~\ref{ppg:lead:part} and Part~\ref{ppg:log:part}, with the section offsets just stated (Report~144's Section~$k$ is Section~$k$ here). Statements added in the write (Remark~\ref{ppg:lead:rem:serra}, Proposition~\ref{ppg:lead:prop:op3}, Remarks~\ref{ppg:lead:rem:inverse}, \ref{ppg:log:rem:inverse}, \ref{ppg:all:rem:op2} and~\ref{ppg:all:rem:inverse}) are the last numbered statements of their sections, so no delivered statement number moved; Section~\ref{ppg:sec:further}, which closes the report, was added in the write.

Notes added when the manuscripts were merged are set in small type and tagged \textbf{[write]}, with their date. They give cross-references between the Parts, mark statements that a later Part supersedes or answers, and mark results proved twice. Superseded statements are printed as delivered, with a dated note pointing forward, as in the merged reports \file{a279551-inversion-log-deficit} and \file{a336070-weak-ascents} of the same collection. No delivered statement was changed and no delivered symbol was renamed; the write added labels to six unlabelled headings and one proposition. Section~\ref{ppg:sec:further} collects, under Vladimir's standing rule of 4~October 2026, every question the three Parts leave open, with its source, its status and what is missing.

\section*{What is proved, and where}\phantomsection\label{ppg:sec:status}
\addcontentsline{toc}{section}{What is proved, and where}
Here $W_n$ is the winning probability, $r_*=\lim W_n^{1/n}$, $R=1/r_*$, $\lambda=\log R$, $a=(2\pi)^{-1/2}$, $A=r_*\sqrt{2\pi}=1/(aR)$, $x_n=W_nR^n/(A\sqrt{n+1})$, and $N(\varepsilon)=\min\{n\ge0:W_n\le\varepsilon\}$.
\begin{center}
\small
\begin{longtable}{@{}>{\raggedright\arraybackslash}p{0.52\textwidth} >{\raggedright\arraybackslash}p{0.43\textwidth}@{}}
\toprule
Statement & Where\\
\midrule
\endhead
Exact recurrence $W_n=\frac1n\sum_kW_{k-1}W_{n-k}Q_{n,k}$; $W_n$ positive rationals; $W_2,\ldots,W_5$ & Part~I, Proposition~\ref{ppg:lead:prop:recurrence} (Serra's Theorem~3.12, derived independently); restated in Sections~\ref{ppg:log:sec:model} and~\ref{ppg:all:sec:main}\\
$r_*$ exists, $1/16\le r_*\le\sqrt3/2$; lower and upper certificates converging to $r_*$; $r_*$ is computable & Part~I, Theorems~\ref{ppg:lead:thm:rate} and~\ref{ppg:lead:thm:uppercert}\\
$527/1000<r_*<79/125$, by exact rational arithmetic through $n=100$ & Part~I, \eqref{ppg:lead:eq:finiteLower}--\eqref{ppg:lead:eq:finiteUpper}\\
Nested rational rate enclosures, width contraction $313/400$ per step and a quadratic-times-log width bound (no run) & Part~III, Theorem~\ref{ppg:all:thm:rate}\\
Uniform kernel bounds and first expansion; $\alpha_{i,j}>a$ & Part~I, Theorems~\ref{ppg:lead:thm:kernel} and~\ref{ppg:lead:thm:stronglower}\\
Exact boundary factorization, covariance $1-\frac7{360ij}$, boundary-matched kernel; endpoint-inclusive factorization and uniform kernel expansions to every order & Part~II, Lemmas~\ref{ppg:log:lem:factor}, \ref{ppg:log:lem:covariance}, Proposition~\ref{ppg:log:prop:matched}; Part~III, Lemmas~\ref{ppg:all:lem:factor} and~\ref{ppg:all:lem:kernel}\\
Pole $F\sim1/(a(R-z))$ with logarithmic second term; cumulative laws; $G(z)\sim\frac{\pi}{\sqrt2R}(1-z/R)^{-3/2}$ & Part~I, Theorems~\ref{ppg:lead:thm:pole}, \ref{ppg:lead:thm:log}, \ref{ppg:lead:thm:averages}\\
$W_n=r_*^nn^{1/2+o(1)}$ & Part~I, Theorem~\ref{ppg:lead:thm:subpower} (superseded by the next line)\\
$W_n\sim r_*\sqrt{2\pi n}\,r_*^n$, $W_{n+1}/W_n\to r_*$, $W_n$ eventually decreasing (Serra's Conjecture~6.10) & Part~I, Theorem~\ref{ppg:lead:thm:pointwise}\\
Any equivalent $Cn^\gamma r^n$ has $\gamma=\frac12$, $r=r_*$, $C=r_*\sqrt{2\pi}$; Serra's $C=\sqrt{2\pi}$ is wrong & Part~I, Theorem~\ref{ppg:lead:thm:conditional}, Section~\ref{ppg:lead:sec:sourcecorrection}; \wtag\ Remark~\ref{ppg:lead:rem:serra}\\
Serra's integral equation~(200) fails, read through the profile $A\sqrt t\,r_*^t$: $E(n)/W(n)\to1-r_*$ & \wtag\ Proposition~\ref{ppg:lead:prop:op3}\\
$\frac{W_n}{A\sqrt n\,r_*^n}=1+\frac1{4n}-\frac7{540}\frac{\log n}{n^2}+o(\frac{\log n}{n^2})$; then $+\frac{B_2}{n^2}+o(n^{-2})$ with $B_2$ by absolutely convergent series & Part~II, Theorems~\ref{ppg:log:thm:main} and~\ref{ppg:log:thm:full}\\
Every fixed order: $x_n=1+\sum_{j\le M}P_j(\log n)n^{-j}+O_M((1+\log n)^M/n^{M+1})$, $\deg P_j\le j-1$ & Part~III, Theorem~\ref{ppg:all:thm:all}\\
Third order explicit; the $(\log n)^2/n^3$ term vanishes; log coefficient $\frac{22}{2835}+cS_A$ & Part~III, Theorem~\ref{ppg:all:thm:third}\\
Ratios: eventual strict decrease and log-concavity; $\rho_n/r_*=1+\frac1{2n}-\frac3{8n^2}+o(n^{-2})$; no pure inverse-power expansion through $n^{-3}$; the cubic term $\frac{(7/270)\log n+32/135-2B_2}{n^3}$; a dyadic five-point extrapolation & Part~II, Corollary~\ref{ppg:log:cor:monotonicity}, Theorem~\ref{ppg:log:thm:ratio}; Part~III, Corollaries~\ref{ppg:all:cor:ratio}, \ref{ppg:all:cor:extrapolation}; \wtag\ Remark~\ref{ppg:all:rem:op2} (Serra's Open Problem~2)\\
Inverses: the generating-function inverse $z(Y)$; the first passage $N(\varepsilon)$ inside ceilings around a Lambert-$W_{-1}$ centre, with vanishing uncertainty, then uncertainty $o(x^{-2})$, then $O((1+\log x)^M/x^{M+1})$ at every fixed order, with Newton centres & Part~I, Theorems~\ref{ppg:lead:thm:GFinverse}, \ref{ppg:lead:thm:threshold}; Part~II, Theorem~\ref{ppg:log:thm:inverse}; Part~III, Theorem~\ref{ppg:all:thm:threshold}; \wtag\ Remarks~\ref{ppg:lead:rem:inverse}, \ref{ppg:log:rem:inverse}, \ref{ppg:all:rem:inverse}\\
\midrule
\multicolumn{2}{@{}l}{\emph{Open in all three Parts}}\\
\multicolumn{2}{@{}p{0.96\textwidth}@{}}{Effective constants and remainders, certified decimals of $B_2$, $B_3$ and of $r_*$ beyond the bracket; strict decrease of $W_{n+1}/W_n$ at every index, and an onset; exact log degrees of the $P_j$; growth of $P_j$ with $j$, convergence, a transseries; a practical certified computation; other policies and optimality; the value of $S_A$ (numerically $1/12$). See Section~\ref{ppg:sec:further}.}\\
\bottomrule
\end{longtable}
\end{center}
The finite computations of the three Parts (an exact rational certificate in Part~I, exact algebra checks in Parts~II and~III) are regression checks of formulas, indices and the certificate inequalities; no asymptotic statement rests on them, and none certifies a decimal digit beyond the bracket of Part~I. The manuscripts make no universal priority claim; they state their results ``relative to the inspected source''.

\section*{Notation across the three Parts}\phantomsection\label{ppg:sec:notation}
\addcontentsline{toc}{section}{Notation across the three Parts}
Common to all Parts and consistent between them: $W_n$, $W_0=W_1=1$; $Q_{n,k}$ the beta-binomial weight; $r_*$, $R=1/r_*$, $\lambda=\log(1/r_*)$; $a=(2\pi)^{-1/2}$; $A=r_*\sqrt{2\pi}=1/(aR)$; $x_n=W_nR^n/(A\sqrt{n+1})$ (in Part~I from Theorem~\ref{ppg:lead:thm:subpower} on, where it is written $u_nR^n/A$ with $u_n=W_n/\sqrt{n+1}$); $y_n=x_n-1$; $X(z)=\sum x_nz^n$; the normalized kernel $\beta_{ij}=\alpha_{i,j}/a$, so that $(n+1)x_{n+1}=\sum_{i+j=n}\beta_{ij}x_ix_j$; the defect $e_n=\sum_{i+j=n}(\beta_{ij}-1)x_ix_j$; $N(\varepsilon)$ the first passage; $W_{-1}$ the lower real Lambert branch. Part~II's $b_i$, $H_m$, $B$, $k=B_2-\frac1{32}$, $d=-\frac14$, $p=\frac14$, $q=-\frac5{12}$ mean the same as Part~III's $b_i$, $H_n$, $B$, $k$, $d$, $p$, $q$. The symbols below change meaning between the Parts or inside one; each Part's text fixes its meaning there, and each Part opens with a short table of its own readings.
\begingroup\small
\begin{longtable}{@{}>{\raggedright\arraybackslash}p{0.13\textwidth} >{\raggedright\arraybackslash}p{0.83\textwidth}@{}}
\toprule
Symbol & Meanings\\
\midrule
\endhead
$c$, $C$, $\kappa$ & Part~II: $c=-\frac7{360}$ the covariance coefficient \eqref{ppg:log:eq:pqs}, $\kappa=-\frac7{540}$ the logarithmic coefficient, $C_x$ the second constant of $x_n$, $C_{ij}$ the covariance ratio, $C_h$. Part~III: $c=-\frac7{540}$ (Part~II's $\kappa$) and $C=B_2+\frac14$ (Part~II's $C_x$) in \eqref{ppg:all:eq:oldsecond}; $C_1$, $c_{10}$, $c_{20}$, $C_M$; $C(z,U)$ in \eqref{ppg:all:eq:Crate}. \emph{Tempting false reading:} Part~II's $c=-7/360$ is not Part~III's $c=-7/540$; Part~III writes the covariance constant $S_0=\ip UU=-\frac7{360}$.\\
$p$, $q$ & Part~I: $p=h/2-1/m$, $q=1/(12m)-h/12$ in \eqref{ppg:lead:eq:Dexpand}--\eqref{ppg:lead:eq:Sexpand}; $q<1$ a ratio bound in Lemma~\ref{ppg:lead:lem:modulus}. Parts~II and~III: $p=\frac14$, $q=-\frac5{12}$. But in Part~III, Sections~\ref{ppg:all:sec:models}--\ref{ppg:all:sec:scalar}, $p$ and $q$ are also positive \emph{integer} indices of $B_{p,q}=u^{p-1}L^q$ and $s_{p,q}$. \emph{False reading:} in $b_{p,q}(n)$ and \eqref{ppg:all:eq:blead}, $p$ is not $\frac14$.\\
$s$, $f$, $F$ & Part~II: $s=-\frac{11}{180}$, $F_2=\frac{217}{288}$, $f_i(v)=(1+v/i)^ie^{-v}$, $f_N$ a bound sequence, $f_0$, $f$ inverse models. Part~III: $f=h_2=\frac{217}{288}$ (Part~II's $F_2$), $s_{p,q}$ model coefficients, $f_0$, $f_M$ inverse models, $F(v)=c(\gamma-1)-V(v)$ in \eqref{ppg:all:eq:VF}, $F(t)$ in \eqref{ppg:all:eq:EM}, $F(z)$ a forcing in \eqref{ppg:all:eq:Tmodel}, and $F(z)=\sum u_nz^n$ in Section~\ref{ppg:all:sec:rate} (Part~I's $F$). Part~I: $F$, $G$ the generating functions of $u_n$ and $W_n$; $f_N=\log b_N$ in Lemma~\ref{ppg:lead:lem:merge}; $f(t)$ the model in Theorem~\ref{ppg:lead:thm:threshold}.\\
$a$, $a_i$, $\delta_i$ & $a=(2\pi)^{-1/2}$ everywhere. Part~II: $a_i=(b_i-1)x_i$ and $a^{(2)}=-\frac{269}{1440}$. Part~III: $\delta_i=(b_i-1)x_i$ (Part~II's $a_i$); $a_i(v)=g_i(v)x_i-1$ and $\bar a_i=b_ix_i-1$; $a_2=-\frac{269}{1440}$ (Part~II's $a^{(2)}$) and $a_3$ in \eqref{ppg:all:eq:a23}; but $a_2$ is also the $\ell^2$ coefficient of $\mathcal L_3$ in \eqref{ppg:all:eq:R123}, $a_0$ the leading forcing coefficient in \eqref{ppg:all:eq:inverse3}, and $a$ an integer exponent in $u^aL^q$ (Section~\ref{ppg:all:sec:induction}).\\
$S$ & Part~I: $S_n$ the modal mass (Lemma~\ref{ppg:lead:lem:mass}); $S_N$ partial sums in Lemma~\ref{ppg:lead:lem:subpower}. Part~II: $S_y$, $S_a$ (residue sums \eqref{ppg:log:eq:residuesums}). Part~III: $S_D$ (Part~II's $S_a$), $S_A=\sum_i\bar a_i$ (new, \eqref{ppg:all:eq:SA}), $S_M(z)$ models, $S_0=-\frac7{360}$.\\
$b$, $B$ & Part~I: $B=\int_0^1T$ in \eqref{ppg:lead:eq:factor1}, the Catalan majorant $B$ in \eqref{ppg:lead:eq:majorant}, $B_0(t)$, $b_N=W_{N-1}$. Part~II: $b_i=\lim_j\beta_{ij}$, $b^{(2)}$, $B$ the forcing constant of $e_n$ and $h_n$, $B_2$, and $b=\frac14$ in Theorem~\ref{ppg:log:thm:inverse}. Part~III: $b_i$ (as in Part~II), $b_2^*$, $b_3^*$, $b_{p,q}(n)$, $B_{p,q}(z)$, $B$, $B_2$, $B_3$, $b=\frac14$ in Section~\ref{ppg:all:sec:inverse}, and the interval endpoint $B=z+v/a$ and Catalan root $B_N(t)$ in Section~\ref{ppg:all:sec:rate}.\\
$A$ & The amplitude everywhere, but: Part~II's $A_i$ are endpoint factors (the source says so); Part~III's $A_0,A_1,A_2$ are constants \eqref{ppg:all:eq:Aconstants} and $A=z+v$ is an interval endpoint in \eqref{ppg:all:eq:roundedendpoints}.\\
$x_0$, $x$ & The normalized coefficients $x_n$, so $x_0=1/A$. But Part~II's Theorem~\ref{ppg:log:thm:inverse} writes $x_0=x_0(\varepsilon)$ for the Lambert root and $x_2,x_3$ for centres, Part~III's Section~\ref{ppg:all:sec:inverse} writes $x=x_0(\varepsilon)$, and Part~I's Section~\ref{ppg:lead:sec:certificates} writes $x_N$ for the root of $4xP_N(x)=N$ and $x=125/79$. \emph{False reading:} $x_0$ in \eqref{ppg:log:eq:x0} is not $1/A$.\\
$r$, $R$ & $r_*$ the rate. $r_k$: Robbins's (Stirling's) remainder in all Parts; also Part~I's $r=1/(2m)-h/6$ in \eqref{ppg:lead:eq:Bexpand} and $r_n=1-1/n$ in Theorem~\ref{ppg:lead:thm:pointwise}; Part~II's remainder $r_n$ in \eqref{ppg:log:eq:inverseerror}; Part~III's residuals $r_n$, $r_{M,n}(v)$ and Stirling coefficients $r[k]$. $R=1/r_*$; Part~II's $R_{ij}$, $\widetilde R_{ij}$ kernel remainders; Part~III's $\mathcal R_j$ inverse polynomials and $\mathcal R(t)$.\\
$h$, $H$, $\HH$ & Part~I: $h=1/i+1/j$, $H$ the modal probability, $H(r)$ and $H(s)$ auxiliary functions. Part~II: $h_n=e_n+z_n$ the forcing, $H(z)$ its generating function, $H_m$ the kernel factor. Part~III: $h_n$, $H_n$ as in Part~II; $H=F'/F^2$ in Section~\ref{ppg:all:sec:rate}. $\HH_N$ harmonic numbers (Parts~II and~III).\\
$D$, $d$ & Part~I: $D$ a normalization factor (Section~\ref{ppg:lead:sec:kernel}), $D_0=10^{12}$, a constant in Theorem~\ref{ppg:lead:thm:conditional}. Part~II: $d=-\frac14$, $D_k=(k+1)(k+2)(k+3)$, $D_0=\lambda-1/(2x_0)$, $\mathcal D_i$ the boundary derivative. Part~III: $d=-\frac14$, $D$ a constant in \eqref{ppg:all:eq:thirdh}, $D=U-z$ in Section~\ref{ppg:all:sec:rate}.\\
$L$, $\ell$ & Parts~I and~II: $L=\log(1/\varepsilon)$; Part~I also $L_n(q)$, $L(q)$, $L_M$. Part~II: $\ell_k=\log(k+1)$, $L_h$. Part~III: $L=-\log(1-z)$, $L_\varepsilon=\log(1/\varepsilon)$, $L_1$, $L_2$, $\mathcal L_j$, $\ell=\log x$, and interval endpoints $L$, $L_0$, $L_k$ in Section~\ref{ppg:all:sec:rate}.\\
$T$ & Part~I: the integrand $T(v)$ \eqref{ppg:lead:eq:T} and $T_0(z)$. Parts~II and~III: the critical inverse $T$ (\eqref{ppg:log:eq:T}, \eqref{ppg:all:eq:T}); Part~III also $T[F]$ in \eqref{ppg:all:eq:Tmodel}.\\
$u$, $U$, $v$, $V$ & Part~I: $u_n=W_n/\sqrt{n+1}$, $v_n=u_nR^n$, $v=1/d$ in \eqref{ppg:lead:eq:gapODE}. Part~III: $u=1-z$ (Sections~\ref{ppg:all:sec:models}--\ref{ppg:all:sec:third}) but $u_n=W_n/\sqrt{n+1}$ in Section~\ref{ppg:all:sec:rate}; $U(v)=\frac16-\frac{v^2}2$ but $U$ an upper endpoint in Section~\ref{ppg:all:sec:rate}; $V(v)$ in \eqref{ppg:all:eq:VF} but $v=1/F$ in Section~\ref{ppg:all:sec:rate}.\\
$E$, $e$ & $e_n$, $E(z)$ the defect (Parts~I and~II, same object); Part~I's $E_{i,j}$ kernel remainder; Part~III's shift $Ef(n)=f(2n)$, $E_M(x)$, $E_{n,i}$, $\mathcal E_J$.\\
$\rho$ & Part~I: $\rho=r_m-r_i-r_j$ and a radius in Lemma~\ref{ppg:lead:lem:modulus}. Part~II: $\rho_{ij}$ kernel remainder; $\rho_n=W_{n+1}/W_n$ (Parts~II and~III).\\
$\gamma$ & Euler's constant in Parts~II and~III; an exponent in Part~I, \eqref{ppg:lead:eq:conditionalintro} and Theorem~\ref{ppg:lead:thm:conditional}.\\
$K$, $\mathcal Q$ & Part~II: $K=\sum y_i/(i+1)$, $K_1=\sum y_i/i$, $\mathcal Q=\sum\theta_ix_i$. Part~III: $K=-c/2$, $K_x$; $Q_k(\log n)$ forcing polynomials, $\mathcal Q_J(E)$ the extrapolation operator.\\
$z$ & The generating-function variable; Part~II's convolution $z_n=\sum y_iy_{n-i}$ (the source says so); Part~III's centres $z_M$, $z^{(k)}$ and rational points $z$ in Section~\ref{ppg:all:sec:rate}.\\
$W$ & $W_n$ the probabilities; $W_{-1}$ the Lambert branch. Serra writes $W(n)$ and $r(n)=W(n+1)/W(n)$, $r_\infty$ for $r_*$.\\
\bottomrule
\end{longtable}
\endgroup

\section*{Provenance and merge decisions}\phantomsection\label{ppg:sec:provenance}
\addcontentsline{toc}{section}{Provenance and merge decisions}
Three manuscripts of the session bundle of Reports 1--243, all dated 3~October 2026. Their title blocks read ``Report~144'', ``Report~145'' and ``Report~146''; none names an author, a tool or an addressee (each source sets an empty PDF author field), and none carries ``prepared for private review'' wording. None pins a ProveIt commit or names a repository path. They arrived in commit \texttt{60f54ea06} and were placed by \texttt{e85586b7c} (batch~105).
\begin{itemize}
\item \textbf{Part~\ref{ppg:lead:part}}, bundle Report~144, archive \file{A398540_Placement_Game_Leading_Asymptotics_and_Inverse_Source.zip} (1,059,711 bytes, 12 files; \file{Report144.tex}, 979 lines, 21 pages). Contributes the recurrence, the rate and its certificates, the kernel theory, the pole, the cumulative laws, the leading equivalent and ratio limit, the necessity theorem with the correction of Serra's amplitude, and the first two inverse theorems.
\item \textbf{Part~\ref{ppg:log:part}}, bundle Report~145, archive \file{A398540_Placement_Game_Logarithmic_Corrections_and_Inverse_Source.zip} (1,034,994 bytes, 16 files; \file{Report145.tex}, 995 lines, 21 pages). Its \file{foundation/} directory holds Report~144's source and PDF, byte-identical to Report~144's delivery. Contributes the exact critical inverse, the tail contraction, the boundary covariance and boundary-matched kernel, the logarithmic correction, the constant $B_2$, the eventual decrease of the ratios, the obstruction to pure-power ratio expansions, and the refined inverse.
\item \textbf{Part~\ref{ppg:all:part}}, bundle Report~146, archive \file{A398540_All_Orders_Logarithmic_Asymptotics_and_Inverses_Source.zip} (1,541,057 bytes, 19 files; \file{Report146.tex}, 950 lines, 22 pages). The base. Its \file{foundation/} directory holds Reports~144 and~145, source and PDF, byte-identical to their deliveries (SHA-256 \texttt{cf59fa57\ldots}, \texttt{9397da72\ldots}, \texttt{9e14d0c4\ldots}, \texttt{59a7f2a3\ldots}, as its \file{146-allorders-SOURCES.md} also records). Contributes the endpoint-inclusive kernel, the coefficient models and signed convolution, the all-orders induction, the explicit third order, the ratio logarithm and dyadic extrapolation, the all-order inverse, and the rate-enclosure theorem.
\end{itemize}
\emph{Report~146 was corrected before delivery.} The bundle's own README says: ``Where a saved archive was corrected in place, the latest corrected version is included: Reports 146 and 205. Their superseded versions are not repeated.'' The version printed here is therefore the corrected one, and the version before the correction was never received: it is not in the bundle, not in the repository's history, and not on this machine (searched at placement). What the correction changed cannot be determined, because the archive carries no changelog, erratum or version note. The delivered Report~146 is internally consistent: its checksum manifest verifies (18 of 18 members), its receipts' hashes match the shipped files, and a rerun of its checker on a copy reproduced its receipt byte for byte (placement dossier, 5~October 2026). Its embedded copies of Reports~144 and~145 are those of the standalone deliveries, so the correction did not touch them.

Where the merge had to choose:
\begin{enumerate}
\item \emph{Order and base.} Dependency order, I--II--III, with the base last (see the Guide). Later Parts improve earlier statements; the earlier statements stay, with dated notes.
\item \emph{Embedded copies.} The \file{foundation/} copies in Reports~145 and~146 are printed once, as Parts~\ref{ppg:lead:part} and~\ref{ppg:log:part}, and are not shipped.
\item \emph{Repetition.} Each manuscript opens with the model and the recurrence in its own words; Report~144 derives the recurrence in full (Proposition~\ref{ppg:lead:prop:recurrence}), Reports~145 and~146 in three or four sentences. All three openings are printed, with notes pointing to Part~\ref{ppg:lead:part}, because they fix each Part's notation and labels. Report~145's exact boundary factorization (Lemma~\ref{ppg:log:lem:factor}) and Report~146's endpoint-inclusive factorization (Lemma~\ref{ppg:all:lem:factor}) are both printed: the second extends the first to a zero child.
\item \emph{Superseded results.} Report~144's inverse (Theorem~\ref{ppg:lead:thm:threshold}) and Report~145's (Theorem~\ref{ppg:log:thm:inverse}) follow from the cases $M=1$ and $M=2$ of Report~146's (Theorem~\ref{ppg:all:thm:threshold}), after a small adjustment of the centre (Remarks~\ref{ppg:lead:rem:inverse} and~\ref{ppg:log:rem:inverse}); Report~145's expansion theorems follow from the case $M=2$ of Theorem~\ref{ppg:all:thm:all}, which uses them as inputs; Report~145's obstruction (Theorem~\ref{ppg:log:thm:ratio}) follows from Corollary~\ref{ppg:all:cor:ratio}. All are printed, with notes; Report~145's proof of its obstruction sums a hypothetical ratio expansion and is kept as an independent earlier route.
\item \emph{Front matter and bibliography.} The three title blocks and abstracts are printed at the heads of their Parts; one table of contents replaces Report~146's. The three bibliographies are merged into one, keeping every detail of every entry; Report~145's entry ``foundation'' and Report~146's entries for Reports~144 and~145 now point to Parts~\ref{ppg:lead:part} and~\ref{ppg:log:part}.
\item \emph{Write additions.} Remark~\ref{ppg:lead:rem:serra} and Proposition~\ref{ppg:lead:prop:op3} (Serra's prefactor and his integral equation, with proofs), Remarks~\ref{ppg:lead:rem:inverse}, \ref{ppg:log:rem:inverse} and~\ref{ppg:all:rem:inverse} (the inverses as instances of the transseries volume, and the implications between them, with proofs), Remark~\ref{ppg:all:rem:op2} (Serra's Open Problem~2, with proof), Section~\ref{ppg:sec:further}, the dated notes, the per-Part reading tables and this front matter. Everything else is delivered text.
\end{enumerate}
The write read Serra's version~1 PDF (92 pages, Zenodo record 21946055, licensed CC~BY~4.0 according to the record) and checked every page and equation number the three Parts cite: Theorem~3.12 with (47)--(48) on printed page~23, its proof ending on page~24, (140) on page~46, (147) on page~47, Proposition~6.17 on page~48, (151) on page~49, Section~8.1 on pages~67--69, each printed page being the PDF page minus one. All are as the Parts say.

\section*{Relation to neighbouring reports and to the formal project}\phantomsection\label{ppg:sec:neighbours}
\addcontentsline{toc}{section}{Relation to neighbouring reports and to the formal project}
All neighbours are in the collection directory \file{SetTheory/Cardinals/docs/reports/generating-functions-and-asymptotics/}\allowbreak\file{oeis-sequence-asymptotics/}. No other report, and no Lean or Rocq file, treats A398540, the placement game, Serra's manuscript or this recurrence (searched 5~October 2026). The three other new reports of batch~105 (\file{a217057-unique-pattern-occurrences}, \file{a224182-unique-1432-order}, \file{a292692-weighted-dyck-newton-diagonal}) treat other sequences, and none of their delivered texts mentions the placement game or A398540.

\emph{The transseries volume.} The inverse theorems of the three Parts are instances of results of the repository's transseries volume \cite{tsvol}, which the manuscripts neither cite nor need: the Lambert centre is the large root of $aX+b\log X=L$ with $a=\lambda>0$ and $b=-\frac12<0$, which the branch rule of \texttt{p0:thm:lambert-core} assigns to $W_{-1}$; the first passage is, for small $\varepsilon$, the integer staircase of \texttt{p0:def:three-inverses} for the eventually increasing sequence $\log(1/W_n)$, and a bracket whose two ceilings agree determines it, as in the separation condition of \texttt{p0:thm:staircase}; the mean-value step that moves the forward error to the inverse is the estimate of \texttt{p0:cor:forward-to-inverse}. The Parts reach their brackets by comparing the sequence with smooth models directly, without the monotonicity or the admissible interpolation that \texttt{p0:thm:staircase} assumes. Part~\ref{ppg:all:part}'s centres for $M\ge2$ carry polynomials in $\log x$ (\eqref{ppg:all:eq:formalcenter}), because the forward coefficients $\mathcal L_j$ are polynomials in $\log n$; they are therefore not instances of \texttt{p0:thm:lambert-centered} or \texttt{p0:thm:flattening}, whose correction series have constant coefficients. Remarks~\ref{ppg:lead:rem:inverse}, \ref{ppg:log:rem:inverse} and~\ref{ppg:all:rem:inverse} give the details. These are pointers to known mechanics, not new inverse theory.

\emph{Formal status.} No statement of this report is formalized. Placement in the \file{SetTheory/Cardinals} report collection confers no formal status.

\clearpage
\part{Asymptotics and a computable exponential rate for a proportional placement game}\label{ppg:lead:part}
\partsource{Bundle Report~144, archive \file{A398540_Placement_Game_Leading_Asymptotics_and_Inverse_Source.zip}, delivered as \file{Report144.tex} (979 lines, 21 pages), dated 3~October 2026; title block ``Report 144''; no author line. Printed as delivered, with \textbf{[write]} notes; section and statement numbers are those of the delivery, its Appendix~A is Section~\ref{ppg:lead:app:reproduce}, and equations are renumbered within sections. Remarks~\ref{ppg:lead:rem:serra} and~\ref{ppg:lead:rem:inverse} and Proposition~\ref{ppg:lead:prop:op3} are added in the write.}
\begin{partabstract}
For the continuous online placement game under its proportional policy, we prove
$W_n\sim r_*\sqrt{2\pi n}\,r_*^n$ and $W_{n+1}/W_n\to r_*$, where $r_*\in(0,1)$ is a computable real number. Exact rational certificates through $n=100$ give $0.527<r_*<0.632$. Starting from the probability recurrence, balanced merging establishes the exponential rate, and a uniform beta-bin kernel expansion gives a pole with an unconditional logarithmic second term for $F(z)=\sum W_nz^n/\sqrt{n+1}$. A positive Tauberian argument yields cumulative laws, followed by a half-split induction giving subpower coefficient envelopes. The strict kernel bound $\alpha_{i,j}>(2\pi)^{-1/2}$ makes the Riccati defect coefficientwise nonnegative. A uniform complex modulus bound and Cauchy extraction then prove the individual equivalent and the ratio limit, without assumed ratio monotonicity or boundary continuation. We also obtain a logarithmically refined inverse for $F$ and a Lambert-W first-passage inverse with vanishing uncertainty inside integer ceilings. The leading amplitude corrects an index shift in the source heuristic for the model associated with OEIS A398540. Higher pointwise corrections, monotonicity of the ratios, and additional decimal digits are not established here.
\end{partabstract}
{\small\emph{Readings in this Part} (\wtag). $u_n=W_n/\sqrt{n+1}$; $F=\sum u_nz^n$ and $G=\sum W_nz^n$; $R=1/r_*$; $A=1/(aR)=r_*\sqrt{2\pi}$; the kernel $\alpha_{i,j}$ of \eqref{ppg:lead:eq:alpha} with $n=i+j+1$ the \emph{parent} size, and $\beta_{i,j}=\alpha_{i,j}/a$ from Section~\ref{ppg:lead:sec:pointwise} on, the $\beta_{ij}$ of Parts~\ref{ppg:log:part} and~\ref{ppg:all:part}; $x_n=u_nR^n/A$ from Theorem~\ref{ppg:lead:thm:subpower} on; $\gamma$ an exponent, not Euler's constant; $L=\log(1/\varepsilon)$ in Section~\ref{ppg:lead:sec:inverses}; $T(v)$ an integrand, not the critical inverse of the later Parts; $p,q,r$ in \eqref{ppg:lead:eq:Dexpand}--\eqref{ppg:lead:eq:Bexpand} are functions of $i,j$, not the constants $p=\frac14$, $q=-\frac5{12}$ of the later Parts; $B$ is $\int_0^1T$ in Section~\ref{ppg:lead:sec:kernel} and the Catalan majorant in Section~\ref{ppg:lead:sec:certificates}.\par}
\section{The model and the conclusions}\label{ppg:lead:sec:model}
There are $n$ ordered slots and $n$ independent uniform random variables on $(0,1)$, revealed sequentially. Each value must be placed immediately in an empty slot so that all values placed so far remain increasing from left to right. A failed placement loses the game. The \emph{proportional policy} acts within the unique gap $(L,U)$ containing a new value $x$. If that gap has $k$ empty slots, it selects the slot numbered
\[
 \left\lceil k\frac{x-L}{U-L}\right\rceil
\]
from the left within the gap; if $k=0$, it loses. Initially $L=0$, $U=1$, and $k=n$. Ties have probability zero and may be resolved by any consistent convention. This policy is not assumed to be optimal.

Write $W_n$ for its winning probability and set $W_0=1$. The probability is independent of the endpoints of a gap after affine rescaling. The exact recursion, derived in Section~\ref{ppg:lead:sec:recurrence}, is
\begin{equation}\label{ppg:lead:eq:recurrence}
 W_n=\frac1n\sum_{k=1}^n W_{k-1}W_{n-k}Q_{n,k},\qquad W_0=W_1=1,
\end{equation}
where $Q_{n,k}$ is an increment of a regularized beta distribution. The important feature is the child-size sum $n-1$, rather than $n$.

The following notation is used throughout:
\begin{equation}\label{ppg:lead:eq:notation}
 u_n=\frac{W_n}{\sqrt{n+1}},\quad
 F(z)=\sum_{n\ge0}u_nz^n,\quad G(z)=\sum_{n\ge0}W_nz^n,
 \quad a=\frac1{\sqrt{2\pi}},\quad A=\frac1{aR}.
\end{equation}
Here $R=1/r_*$ is justified by the rate theorem below. The summary expansions for $F$ and $G$ below are along the positive real axis. Section~\ref{ppg:lead:sec:pointwise} separately proves the explicit complex estimates needed for coefficient extraction.

\begin{theorem}[Summary of the unconditional conclusions]\label{ppg:lead:thm:summary}
The limit $r_* = \lim_{n\to\infty}W_n^{1/n}$ exists, is computable, and satisfies
\begin{equation}\label{ppg:lead:eq:numeric}
 \frac{527}{1000}<r_*<\frac{79}{125}.
\end{equation}
The individual probabilities and consecutive ratios satisfy
\begin{equation}\label{ppg:lead:eq:summarypointwise}
 W_n\sim A\sqrt n\,r_*^n=r_*\sqrt{2\pi n}\,r_*^n,
 \qquad\frac{W_{n+1}}{W_n}\longrightarrow r_*.
\end{equation}
With the notation \eqref{ppg:lead:eq:notation}, as $z\uparrow R$,
\begin{equation}\label{ppg:lead:eq:summaryF}
 F(z)=\frac1{a(R-z)}-\frac1{4aR}\log\frac1{R-z}+O(1).
\end{equation}
As $N\to\infty$,
\begin{align}
 \sum_{n=0}^N\frac{W_nR^n}{\sqrt{n+1}}&\sim AN,\label{ppg:lead:eq:summaryV}\\
 \sum_{n=0}^N W_nR^n&\sim\frac{2A}{3}N^{3/2},\label{ppg:lead:eq:summaryB}
\end{align}
and
\begin{equation}\label{ppg:lead:eq:summaryG}
 G(z)\sim\frac{\pi}{\sqrt2 R}(1-z/R)^{-3/2}.
\end{equation}
The unique real inverse $F(z(Y))=Y$, $Y\ge1$, satisfies
\begin{equation}\label{ppg:lead:eq:summaryinverse}
 z(Y)=R-\frac1{aY}+\frac1{4a^2R}\frac{\log Y}{Y^2}+O(Y^{-2}).
\end{equation}
For $0<\varepsilon<1$, the first crossing
$N(\varepsilon)=\min\{n\ge0:W_n\le\varepsilon\}$ exists and obeys
\begin{equation}\label{ppg:lead:eq:summarythreshold}
 N(\varepsilon)=\frac{L+\tfrac12\log L+\log A-\tfrac12\log\lambda}{\lambda}+O(1),
 \quad L=\log(1/\varepsilon),\quad\lambda=\log(1/r_*).
\end{equation}
\end{theorem}

\begin{writenote}[5 October 2026]
Every line of Theorem~\ref{ppg:lead:thm:summary} is proved in this Part; the later Parts sharpen three of them. The equivalent \eqref{ppg:lead:eq:summarypointwise} acquires the corrections $1+\frac1{4n}-\frac7{540}\frac{\log n}{n^2}+\frac{B_2}{n^2}$ (Theorems~\ref{ppg:log:thm:main} and~\ref{ppg:log:thm:full}) and then every fixed order (Theorem~\ref{ppg:all:thm:all}); the ratio limit becomes $W_{n+1}/(r_*W_n)=1+\frac1{2n}-\frac3{8n^2}+O(n^{-3}\log n)$, with eventual strict decrease (Corollaries~\ref{ppg:log:cor:monotonicity} and~\ref{ppg:all:cor:ratio}); and the threshold \eqref{ppg:lead:eq:summarythreshold} is refined by Theorems~\ref{ppg:log:thm:inverse} and~\ref{ppg:all:thm:threshold} (Remark~\ref{ppg:lead:rem:inverse}). The bracket \eqref{ppg:lead:eq:numeric} remains the only certified numerical statement about $r_*$ in this report; the 18 digits $0.575472381511152252$ of OEIS A398540 lie inside it, and Theorem~\ref{ppg:all:thm:rate} gives a contracting algorithm that starts from it but was not run. The equivalent \eqref{ppg:lead:eq:summarypointwise} is Serra's Conjecture~6.10 with its constant determined; Remark~\ref{ppg:lead:rem:serra} shows why his predicted constant is wrong.
\end{writenote}
Theorem~\ref{ppg:lead:thm:threshold} strengthens the last estimate to a Lambert-W formula and a vanishing uncertainty inside ceilings. The proofs separate exponential, averaged, and pointwise information: the radial expansion is not itself used as a coefficient transfer theorem. Section~\ref{ppg:lead:sec:pointwise} supplies the needed complex modulus control. Section~\ref{ppg:lead:sec:conditional} also records the independent consistency statement
\begin{equation}\label{ppg:lead:eq:conditionalintro}
 W_n\sim Cn^\gamma r^n\quad(C,r>0,\ \gamma\in\mathbb R)
 \quad\Longrightarrow\quad
 \gamma=\tfrac12,\quad r=r_*,\quad C=r_*\sqrt{2\pi}.
\end{equation}

\subsection{Relationship to the source and OEIS}\label{ppg:lead:sub:source}
Serra's manuscript~\cite{serra}, version~1, contains the proportional-policy model and the exact recurrence in Theorem~3.12, equations~(47)--(48), printed pages~23--24. The manuscript is dated July~22, 2026; the Zenodo version was published August~15, 2026. Its Section~8.1, printed pages~67--69, lists the full equivalent and ratio questions among the open problems. Printed page numbers are one less than PDF page numbers in that 92-page version.

OEIS A398540~\cite{oeis} records decimal digits attached to the conjectural limit of $W_{n+1}/W_n$, approximately $0.5754723815\ldots$; it does not enumerate the rational probabilities $W_n$. We define $r_*$ first as the nth-root limit and subsequently prove convergence of the consecutive ratios to that same value. This establishes existence of the constant described by the entry, while keeping its numerical certification separate. The bracket \eqref{ppg:lead:eq:numeric} is the numerical certification supplied by this article and does not certify the posted expansion.

The amplitude in \eqref{ppg:lead:eq:summarypointwise} also corrects a normalization in the source's heuristic integral, as explained precisely in Section~\ref{ppg:lead:sec:sourcecorrection}. All proofs below start from \eqref{ppg:lead:eq:recurrence}; they do not rely on numerical extrapolation, a presumed ratio expansion, or the source's proposed order for resolving its open problems. These are contributions for this exact recurrence relative to the inspected source, without a universal priority claim.
\begin{writenote}[5 October 2026]
The write checked the citations of this subsection against Serra's version-1 PDF (see the provenance section); they are accurate. The OEIS entry, fetched on 5~October 2026, gives the 18 digits $0.575472381511152252$ and, in a formula line by Kot\v{e}\v{s}ovec (4~August 2026), $W(n)\sim c\,d^n\sqrt n$ with $c=1.4424953427\ldots$; another formula line calls the ratio limit ``conjectured''. That $c$ is $r_*\sqrt{2\pi}$ evaluated at the posted digits, the corrected amplitude of Section~\ref{ppg:lead:sec:sourcecorrection}, not Serra's $\sqrt{2\pi}$; as the text says, no proof here uses those digits. OEIS data are licensed CC~BY-SA~4.0.
\end{writenote}

\section{The exact probability recurrence}\label{ppg:lead:sec:recurrence}
Put
\[
 p_{m,j}(t)=\binom mj t^j(1-t)^{m-j},\qquad 0\le j\le m.
\]
For integers $b,c\ge1$, let
\[
 I(t;b,c)=\frac{\int_0^t s^{b-1}(1-s)^{c-1}\,ds}
                   {\int_0^1 s^{b-1}(1-s)^{c-1}\,ds}
\]
be the regularized incomplete beta function. Define
\begin{equation}\label{ppg:lead:eq:Q}
 Q_{n,k}=I(k/n;k,n-k+1)-I((k-1)/n;k,n-k+1).
\end{equation}
The beta integral for integers gives the exact identity
\begin{equation}\label{ppg:lead:eq:integral}
 \frac{Q_{n,k}}n=\int_{(k-1)/n}^{k/n}p_{n-1,k-1}(t)\,dt.
\end{equation}

\begin{proposition}[First-step decomposition]\label{ppg:lead:prop:recurrence}
The winning probabilities satisfy \eqref{ppg:lead:eq:recurrence} for every $n\ge1$.
\end{proposition}
\begin{proof}
Condition on the first value $t$, which has density one. For
$(k-1)/n<t<k/n$, the proportional policy places it in slot $k$. A win requires exactly $k-1$ of the remaining $n-1$ values to lie below $t$; the probability of that event is $p_{n-1,k-1}(t)$. Conditional on that count, the subsequences on the two sides are independent uniform samples on their respective intervals, in their original relative arrival order. After affine rescaling they are proportional-policy subgames with $k-1$ and $n-k$ slots. Their interleaving does not affect either policy within its own gap. Their joint winning probability is therefore $W_{k-1}W_{n-k}$. Integrate over the interval and sum over $k$, then use \eqref{ppg:lead:eq:integral}. At $n=1$, $Q_{1,1}=1$, giving $W_1=1$.
\end{proof}

The integral shows that every $W_n$ is a positive rational number. Two exact forms useful for reproduction are
\begin{align}
 \frac{Q_{n,k}}n
 &=\binom{n-1}{k-1}\sum_{h=0}^{n-k}(-1)^h\binom{n-k}{h}
 \frac{(k/n)^{k+h}-((k-1)/n)^{k+h}}{k+h},\label{ppg:lead:eq:poly}\\
 I(t;k,n-k+1)&=\sum_{h=k}^n\binom nh t^h(1-t)^{n-h}.\label{ppg:lead:eq:tail}
\end{align}
The second identity follows by differentiating the sum and telescoping, then checking its value at zero. It supplies an independent, non-alternating polynomial representation of each beta increment. In particular,
\begin{equation}\label{ppg:lead:eq:prefix}
 W_2=\frac34,\quad W_3=\frac{83}{162},\quad
 W_4=\frac{55537}{165888},\quad W_5=\frac{11049709}{51840000}.
\end{equation}
\begin{writenote}[5 October 2026]
Serra's Theorem~3.12 (his (47)--(48), printed page~23) states the same recurrence for $n\ge2$ with $W(0)=W(1)=1$, and his abstract gives $W(2)=3/4$, $W(3)=83/162$ and $W(4)=55537/165888$, as in \eqref{ppg:lead:eq:prefix}. Parts~\ref{ppg:log:part} and~\ref{ppg:all:part} restate the derivation in a few sentences (Sections~\ref{ppg:log:sec:model} and~\ref{ppg:all:sec:main}); this proposition is the full proof for all three Parts.
\end{writenote}

\section{A positive and strictly subunit exponential rate}\label{ppg:lead:sec:rate}
\begin{lemma}[Modal split bound]\label{ppg:lead:lem:mode}
For every $n\ge1$ and $1\le k\le n$, $Q_{n,k}\ge1/n$.
\end{lemma}
\begin{proof}
For $0<t<1$, the ratio of consecutive binomial masses is
\[
 \frac{p_{m,j+1}(t)}{p_{m,j}(t)}
 =\frac{(m-j)t}{(j+1)(1-t)}.
\]
Hence $\lfloor(m+1)t\rfloor$ is a mode when $(m+1)t$ is not an integer; at an integer the adjacent modes tie. On the interior of the $k$th interval in \eqref{ppg:lead:eq:integral}, $k-1$ is a mode of $\operatorname{Bin}(n-1,t)$. Among $n$ masses summing to one, a modal mass is at least $1/n$. Integrating over length $1/n$ proves the claim; interval endpoints have measure zero.
\end{proof}

\begin{lemma}[Contracting total mass]\label{ppg:lead:lem:mass}
Let $S_n=n^{-1}\sum_{k=1}^nQ_{n,k}$. Then $S_n$ is nonincreasing and $S_n\le3/4$ for $n\ge2$. Moreover,
\begin{equation}\label{ppg:lead:eq:expupper}
 0<W_n\le(3/4)^{\lfloor n/2\rfloor}\qquad(n\ge0).
\end{equation}
\end{lemma}
\begin{proof}
The modal partition gives
\[
 S_n=\int_0^1\max_{0\le j<n}p_{n-1,j}(t)\,dt.
\]
For fixed $t$, the mass function of $\operatorname{Bin}(m+1,t)$ is the convolution of that of $\operatorname{Bin}(m,t)$ with a Bernoulli mass function. Every new mass is a convex combination of two old masses, treating masses outside the support as zero. Its maximum cannot increase. Integrate, and compute $S_2=\int_0^1\max(t,1-t)\,dt=3/4$. Positivity in \eqref{ppg:lead:eq:expupper} follows from the recurrence. For its upper bound, use induction and, when $i+j=n-1$,
\[
 \lfloor i/2\rfloor+\lfloor j/2\rfloor\ge\lfloor n/2\rfloor-1.
\]
Inserting the inductive estimates in \eqref{ppg:lead:eq:recurrence} and multiplying by $S_n\le3/4$ completes the induction from $n=0,1$.
\end{proof}

\begin{lemma}[Balanced merging]\label{ppg:lead:lem:merge}
Set $b_N=W_{N-1}$ and $f_N=\log b_N$ for $N\ge1$. The finite limit
$\lambda_0=\lim_{N\to\infty}f_N/N$ exists. For every $m\ge1$,
\begin{equation}\label{ppg:lead:eq:lambda}
 \lambda_0\ge\frac{f_m-2\log m-4\log2}{m}.
\end{equation}
\end{lemma}
\begin{proof}
The single term $k=p$ of the recurrence at $n=p+q-1$, together with Lemma~\ref{ppg:lead:lem:mode}, gives
\begin{equation}\label{ppg:lead:eq:almostsuper}
 b_{p+q}\ge\frac{b_pb_q}{(p+q)^2},\qquad
 f_{p+q}\ge f_p+f_q-2\log(p+q).
\end{equation}
The slightly stronger denominator $(p+q-1)^2$ has been weakened for convenience.

Fix $m$. Start with $q$ blocks each of size $m$, pairing the available blocks at each level and carrying any unpaired block. At level $\ell$ there are at most $q/2^\ell+1$ merges, each of total size at most $2^\ell m$. There are at most $\lceil\log_2q\rceil$ levels. Repeated use of \eqref{ppg:lead:eq:almostsuper} yields
\begin{align*}
 f_{qm}&\ge qf_m-2\sum_{\ell=1}^{\lceil\log_2q\rceil}
       (q2^{-\ell}+1)\log(2^\ell m)\\
 &\ge qf_m-2q(\log m+2\log2)-O_m((\log q)^2),
\end{align*}
since $\sum_{\ell\ge1}2^{-\ell}=1$ and $\sum_{\ell\ge1}\ell2^{-\ell}=2$.
For general $N=qm+s$, $0\le s<m$, merge once more with a block of size $s$ if $s>0$. The finite set of $f_s$ is bounded below and this extra defect is $O(\log N)$. Therefore
\begin{equation}\label{ppg:lead:eq:liminf}
 \liminf_{N\to\infty}\frac{f_N}{N}
 \ge\frac{f_m}{m}-\frac{2\log m+4\log2}{m}.
\end{equation}
Taking $m=1$ bounds the liminf below by $-4\log2$; also $f_N\le0$. Choose arbitrarily large $m$ along a subsequence approaching the finite limsup of $f_m/m$ in \eqref{ppg:lead:eq:liminf}. The error tends to zero, proving that liminf and limsup agree and giving \eqref{ppg:lead:eq:lambda}.
\end{proof}

\begin{theorem}[Rate existence and lower certificates]\label{ppg:lead:thm:rate}
There is a number $r_*\in[1/16,\sqrt3/2]$ such that $W_n^{1/n}\to r_*$. Every $m\ge1$ supplies the lower certificate
\begin{equation}\label{ppg:lead:eq:lowercert}
 \ell_m=\left(\frac{W_{m-1}}{16m^2}\right)^{1/m}\le r_*,
 \qquad \ell_m\longrightarrow r_*.
\end{equation}
Both $F$ and $G$ have finite radius of convergence $R=1/r_*$.
\end{theorem}
\begin{proof}
Exponentiate Lemma~\ref{ppg:lead:lem:merge}, observing
$\log W_n/n=((n+1)/n)f_{n+1}/(n+1)$. Equation~\eqref{ppg:lead:eq:expupper} gives the upper bound $\sqrt3/2$; $m=1$ gives $1/16$. Formula~\eqref{ppg:lead:eq:lambda} gives \eqref{ppg:lead:eq:lowercert}; the convergence follows because the penalty $(2\log m+4\log2)/m$ tends to zero. Cauchy--Hadamard applies to $F$ and $G$; the factor $(n+1)^{-1/2}$ does not change the nth-root rate.
\end{proof}

\section{Uniform control of the normalized kernel}\label{ppg:lead:sec:kernel}
For $i,j\ge0$, put $n=i+j+1$ and
\begin{equation}\label{ppg:lead:eq:alpha}
 \alpha_{i,j}=Q_{n,i+1}\sqrt{\frac{(i+1)(j+1)}{n+1}}.
\end{equation}
Then \eqref{ppg:lead:eq:recurrence} becomes exactly
\begin{equation}\label{ppg:lead:eq:normalized}
 nu_n=\sum_{i+j=n-1}\alpha_{i,j}u_iu_j\qquad(n\ge1).
\end{equation}
In particular $\alpha_{0,0}=1/\sqrt2=u_1$; the constant coefficient is retained when this identity is differentiated as a generating function.

\begin{theorem}[Global kernel estimates]\label{ppg:lead:thm:kernel}
For every $i,j\ge0$, with $n=i+j+1$,
\begin{align}
 ae^{-2/3}&\le\alpha_{i,j}\le1,\label{ppg:lead:eq:kbound}\\
 |\alpha_{i,j}-a|&\le\frac1{i+1}+\frac1{j+1},\label{ppg:lead:eq:kerror}\\
 \alpha_{i,j}&=a\left\{1+\frac14\left(\frac1{i+1}+\frac1{j+1}\right)
                   -\frac5{12n}+E_{i,j}\right\},\label{ppg:lead:eq:kexpansion}\\
 |E_{i,j}|&\le10\left(\frac1{(i+1)^2}+\frac1{(j+1)^2}\right).\label{ppg:lead:eq:kremainder}
\end{align}
Thus $\alpha_{i,j}\to a$ uniformly whenever $\min(i,j)\to\infty$, with no positive lower bound on either split fraction required.
\end{theorem}

\subsection{Exact mode factorization}
Assume first $i,j\ge1$ and set $m=i+j$. The modal probability and a normalization factor are
\[
 H=\binom mi(i/m)^i(j/m)^j,\qquad
 D=\sqrt{\frac{m(i+1)(j+1)}{ij(m+2)}}.
\]
Substitute $t=(i+v)/(m+1)$ in \eqref{ppg:lead:eq:integral}. Then
\begin{align}
 Q_{n,i+1}&=HB,\quad B=\int_0^1 T(v)\,dv,\label{ppg:lead:eq:factor1}\\
 T(v)&=\left(\frac m{m+1}\right)^m
       (1+v/i)^i(1+(1-v)/j)^j.\label{ppg:lead:eq:T}
\end{align}
The mode $i/m$ belongs to the integration interval, so $T(v)\le1$. Robbins' classical bounds~\cite{robbins} state
\begin{equation}\label{ppg:lead:eq:robbins}
 k!=\sqrt{2\pi k}(k/e)^ke^{r_k},\qquad
 \frac1{12k+1}<r_k<\frac1{12k}\quad(k\ge1).
\end{equation}
Writing $\rho=r_m-r_i-r_j$ gives the exact factorization
\begin{equation}\label{ppg:lead:eq:factor2}
 \alpha_{i,j}=aDe^\rho B.
\end{equation}
Let $h=1/i+1/j$. The following elementary inequalities will be used repeatedly:
\begin{equation}\label{ppg:lead:eq:basicfactors}
 -h/12\le\rho<0,\quad 1\le D\le\sqrt2,\quad
 0\le\log D\le h/2.
\end{equation}
For the sign of $\rho$, $m\ge i+1$ implies
$r_m<1/(12m)<1/(12i+1)<r_i$. For $D$, its logarithm is
\[
 \tfrac12\{\log(1+1/i)+\log(1+1/j)-\log(1+2/m)\}.
\]
The lower square bound follows from $m(m+1)\ge2ij$; the upper one follows from $ij\ge m-1$ and
$ij(m+4)-m(m+1)\ge2m-4\ge0$.
Finally $\log(1+x)\ge x-x^2/2$ and $\log(1+x)\le x$ for $x\ge0$ give
\begin{equation}\label{ppg:lead:eq:logT}
 0\ge\log T(v)\ge-\tfrac12\{v^2/i+(1-v)^2/j\}
 \ge-\tfrac12\max(1/i,1/j).
\end{equation}
Consequently $e^{-1/2}\le B\le1$ and $-h/2\le\log B\le0$. Combining the factor bounds proves
\[
 ae^{-2/3}\le\alpha_{i,j}\le a\sqrt2=1/\sqrt\pi<1
\]
for every interior split.

\subsection{The global first error and the endpoints}
Equations~\eqref{ppg:lead:eq:factor2}--\eqref{ppg:lead:eq:logT} imply
\[
 -7h/12\le\log(\alpha_{i,j}/a)\le h/2.
\]
Put $s=1/(i+1)+1/(j+1)$. In the interior, $h\le2s$ and $s\le1$. Hence
\[
 -7s/6\le\log(\alpha_{i,j}/a)\le s.
\]
On $0\le s\le1$, convexity gives $e^s-1\le(e-1)s$, whereas
$1-e^{-7s/6}\le7s/6$. Since $a(e-1)<1$ and $7a/6<1$, this proves \eqref{ppg:lead:eq:kerror} for interior indices.

At an endpoint the exact expression is
\begin{equation}\label{ppg:lead:eq:endpoint}
 \alpha_{0,n-1}=\alpha_{n-1,0}
 =\left\{1-(1-1/n)^n\right\}\sqrt{\frac n{n+1}}.
\end{equation}
It also holds at $n=1$. Its value lies between $(1-e^{-1})/\sqrt2$ and~1, so it satisfies \eqref{ppg:lead:eq:kbound}. Its distance from $a$ is less than~1, while $s=1+1/n\ge1$, proving \eqref{ppg:lead:eq:kerror}. Notice that $\alpha_{0,n-1}\to1-e^{-1}$, not $a$. The requirement that both endpoint distances grow is essential to the kernel limit.

\subsection{The first correction with an explicit remainder}
Continue in the interior. Put $v_2=i^{-2}+j^{-2}$; then $h^2\le2v_2$ and $m^{-2}\le v_2/8$. Taylor's theorem for the logarithm and exponential gives
\begin{align}
 D&=1+p+e_D,& p&=h/2-1/m,& |e_D|&\le3v_2/4,\label{ppg:lead:eq:Dexpand}\\
 e^\rho&=1+q+e_S,& q&=1/(12m)-h/12,& |e_S|&\le v_2/64,\label{ppg:lead:eq:Sexpand}\\
 B&=1+r+e_B,& r&=1/(2m)-h/6,& |e_B|&\le5v_2/8.\label{ppg:lead:eq:Bexpand}
\end{align}
Here are details of the error bounds. First,
$|\log D-(h/2-1/m)|\le v_2/4+m^{-2}\le3v_2/8$.
Since $D\le\sqrt2$ and $0\le\log D\le h/2$, the additional exponential error is at most $\sqrt2 h^2/8\le\sqrt2 v_2/4$, and $3/8+\sqrt2/4<3/4$.
Next, \eqref{ppg:lead:eq:robbins} gives
$|r_k-1/(12k)|\le1/(144k^2)$. As $\rho\le0$ and $|\rho|\le h/12$, the error in \eqref{ppg:lead:eq:Sexpand} is at most
\[
 \frac{v_2+m^{-2}}{144}+\frac{h^2}{288}
 \le(1/128+1/144)v_2<v_2/64.
\]
Finally, uniformly in $0\le v\le1$,
\[
 \log T(v)=\frac1{2m}-\frac{v^2}{2i}-\frac{(1-v)^2}{2j}+\epsilon(v),
 \qquad |\epsilon(v)|\le3v_2/8,
\]
by the cubic remainder bound for $\log(1+x)$. The exponential error is at most $h^2/8\le v_2/4$ by \eqref{ppg:lead:eq:logT}. Integration gives \eqref{ppg:lead:eq:Bexpand}.

For fully explicit multiplication, use $|p|\le h/2$, $|q|\le h/12$,
$|1-e^\rho B|\le7h/12$, and $|1-B|\le h/2$. The identity
\begin{align*}
 De^\rho B-(1+p+q+r)
 =e_De^\rho B+p(e^\rho B-1)+e_SB+q(B-1)+e_B
\end{align*}
has absolute value at most
\[
 \left(\frac34+\frac7{12}+\frac1{64}+\frac1{12}+\frac58\right)v_2
 =\frac{395}{192}v_2.
\]
Since $p+q+r=h/4-5/(12m)$, replacing $h$ by
$1/(i+1)+1/(j+1)$ and $m$ by $n=m+1$ adds at most
\[
 v_2/4+5/(12m^2)\le29v_2/96.
\]
The total is $151v_2/64$. As
$v_2\le4\{(i+1)^{-2}+(j+1)^{-2}\}$ and $151/16<10$, this proves \eqref{ppg:lead:eq:kexpansion}--\eqref{ppg:lead:eq:kremainder} in the interior. At an endpoint the displayed bracket without $E_{i,j}$ is $5/4-1/(6n)$. Formula~\eqref{ppg:lead:eq:endpoint}, $0<\alpha_{i,j}\le1$ and $1/a=\sqrt{2\pi}<3$ bound the absolute relative error by~3, while the sum of squared reciprocals is at least~1. The same constant~10 covers the endpoints and completes the proof of Theorem~\ref{ppg:lead:thm:kernel}.

\subsection{A strict lower bound at the bulk constant}
The coarse lower bound in \eqref{ppg:lead:eq:kbound} suffices for earlier estimates, but the following strengthening is crucial for the pointwise theorem.

\begin{theorem}[Positive normalized excess]\label{ppg:lead:thm:stronglower}
For every $i,j\ge0$, $\alpha_{i,j}>a$. For interior indices $i,j\ge1$,
\begin{equation}\label{ppg:lead:eq:stronglower}
 \log(\alpha_{i,j}/a)\ge\frac5{324(i+j)}.
\end{equation}
\end{theorem}
\begin{proof}
For $i,j\ge1$, use the exact factorization \eqref{ppg:lead:eq:factor2}, put $m=i+j$, and apply Jensen's inequality to $B=\int_0^1T(v)\,dv$. Robbins' lower estimate for $r_m$ and upper estimates for $r_i,r_j$ give, after direct integration of $\log T$,
\[
 \log(\alpha_{i,j}/a)\ge \Phi(i)+\Phi(j)-\Psi(m),
\]
where
\begin{align*}
 \Phi(k)&=\{k(k+1)+1/2\}\log(1+1/k)-k-1/(12k),\\
 \Psi(m)&=m\log(1+1/m)+\tfrac12\log(1+2/m)-1/(12m+1).
\end{align*}
For example, the integration uses
$\int_0^1 i\log(1+v/i)\,dv=i(i+1)\log(1+1/i)-i$.

For $x=1/k\in(0,1]$, set $t=x/(2+x)$. The positive series
$\log(1+x)=2(t+t^3/3+t^5/5+\cdots)$ gives
\begin{align*}
 \Phi(k)&\ge\frac{1+x}{2+x}-\frac x{12}
       +\frac{2x(1+x+x^2/2)}{3(2+x)^3}\\
 &\ge\frac12+\frac x6-\frac x{12}+\frac{2x}{81}
  =\frac12+\frac{35x}{324}.
\end{align*}
For $m\ge2$, the inequalities
$\log(1+y)\le y-y^2/2+y^3/3$ and
$\log(1+y)\le y-y^2/[2(1+y)]$ for $y\ge0$ yield
\begin{align*}
 \Psi(m)&\le1+\frac1{2m}+\frac1{3m^2}
                 -\frac1{m(m+2)}-\frac1{12m+1}\\
 &\le1+\frac5{12m}.
\end{align*}
The second inequality is equivalent to
$4(2m-2)(12m+1)\ge m(m+2)$; their difference is
$95m^2-90m-8>0$ for $m\ge2$. Since $1/i+1/j\ge4/m$,
\[
 \log(\alpha_{i,j}/a)\ge\frac1m\left(\frac{35}{81}-\frac5{12}\right)
 =\frac5{324m}>0.
\]
At an endpoint, \eqref{ppg:lead:eq:endpoint} gives
$\alpha_{0,n-1}\ge(1-e^{-1})/\sqrt2>a$.
For instance $1-e^{-1}>3/5$ and $\pi>3$ suffice for the strict comparison. This also covers $n=1$, where $\alpha_{0,0}=1/\sqrt2$.
\end{proof}
\begin{writenote}[5 October 2026]
Part~\ref{ppg:log:part} refines the exact mode factorization \eqref{ppg:lead:eq:factor2} into a product of endpoint profiles, $\beta_{ij}=b_ib_jH_mC_{ij}$ with $\beta_{ij}=\alpha_{i,j}/a$ (Lemma~\ref{ppg:log:lem:factor}), with the mixed covariance $C_{ij}=1-\frac7{360ij}+O(i^{-2}j^{-1}+i^{-1}j^{-2})$ (Lemma~\ref{ppg:log:lem:covariance}); Part~\ref{ppg:all:part} writes it as $\beta_{ij}=H_n\int_0^1g_i(v)g_j(1-v)\,\dd v$, including a zero child (Lemma~\ref{ppg:all:lem:factor}), and expands $g_n$ and $H_n$ to every order (Lemma~\ref{ppg:all:lem:kernel}). The bounds $a<\alpha_{i,j}\le1$ and \eqref{ppg:lead:eq:kerror} of this section are the inputs \eqref{ppg:all:eq:alphainput} of Part~\ref{ppg:all:part}'s rate algorithm.
\end{writenote}

\section{Finite convergent enclosures and exact certification}\label{ppg:lead:sec:certificates}
Let $P_N(z)=\sum_{0\le n<N}u_nz^n$, for $N\ge1$.
\begin{theorem}[Upper certificates and computability]\label{ppg:lead:thm:uppercert}
The equation
\begin{equation}\label{ppg:lead:eq:xN}
 4x_NP_N(x_N)=N
\end{equation}
has a unique positive solution. For every $N\ge1$,
\[
 x_N\le R,\qquad r_*\le1/x_N,\qquad x_N\longrightarrow R.
\]
Consequently the finite bounds
\begin{equation}\label{ppg:lead:eq:nested}
 L_M=\max_{1\le m\le M}\ell_m,\qquad
 U_M=\min\left(\frac{\sqrt3}{2},\min_{1\le N\le M}\frac1{x_N}\right)
\end{equation}
satisfy $L_M\le r_*\le U_M$, are nested, and have $U_M-L_M\to0$.
\end{theorem}
\begin{proof}
The left side of \eqref{ppg:lead:eq:xN} is continuous and strictly increasing from zero to infinity because $P_N$ has positive constant coefficient. Let $B$ be the nonnegative formal solution of
\begin{equation}\label{ppg:lead:eq:majorant}
 B=P_N+\frac zN B^2.
\end{equation}
Coefficient induction gives $[z^n]B\ge u_n$: for $n<N$ the term $P_N$ already supplies $u_n$, while for $n\ge N$ equation~\eqref{ppg:lead:eq:normalized} and $\alpha_{i,j}\le1$ give
\[
 u_n\le\frac1N\sum_{i+j=n-1}u_iu_j.
\]
The factor $z$ in \eqref{ppg:lead:eq:majorant} ensures that only earlier coefficients occur. The Catalan expansion of $B$ converges at every real $0\le x<x_N$, since $4xP_N(x)/N<1$. Positivity and coefficient domination imply that $F(x)$ converges there; hence $R\ge x_N$.

For fixed $0<x<R$, $P_N(x)\le F(x)<\infty$, so $4xP_N(x)<N$ for all sufficiently large $N$. Thus eventually $x_N>x$. Together with $x_N\le R$, this proves $x_N\to R$. Now use $\ell_m\to r_*$ from Theorem~\ref{ppg:lead:thm:rate} to obtain \eqref{ppg:lead:eq:nested} and convergence of its width.
\end{proof}

The theorem is a computability result, with no claim of an efficient convergence rate. Every needed $W_n$ is rational, and $P_N$ has explicitly computable algebraic coefficients. Rational interval arithmetic encloses those coefficients and the positive root in \eqref{ppg:lead:eq:xN} to arbitrary precision. Exact equality at a trial point need not be decided: the derivative of $4xP_N(x)-N$ is at least~4 for $x\ge0$, so an interval enclosure of the function value also provides a certified local root enclosure. Refine the algebraic quantities and increase $M$ until the rational outer endpoints have the requested gap. Since the true gaps tend to zero, this procedure terminates for every prescribed positive tolerance.

\subsection{A finite exact demonstration}
Use $m=N=101$, the rational $x=125/79$, and scale $D_0=10^{12}$. Define
\[
 s_n=\operatorname{isqrt}((n+1)D_0^2),\qquad
 P^{+}_{101}(x)=\sum_{n=0}^{100}\frac{W_nD_0}{s_n}x^n.
\]
The integer-square-root inequalities
\[
 s_n^2\le(n+1)D_0^2<(s_n+1)^2
\]
prove $0<s_n/D_0\le\sqrt{n+1}$, so $P_{101}(x)\le P^+_{101}(x)$. Exact arithmetic on the rational recurrence gives
\begin{align}
 W_{100}&>16\cdot101^2(527/1000)^{101},\label{ppg:lead:eq:finiteLower}\\
 4xP^+_{101}(x)&<95111/1000<101.\label{ppg:lead:eq:finiteUpper}
\end{align}
The first inequality and \eqref{ppg:lead:eq:lowercert} give $r_*>527/1000$. The second gives $x<x_{101}\le R$ and thus $r_*<79/125$. The companion recomputes all $101$ probabilities using both \eqref{ppg:lead:eq:poly} and \eqref{ppg:lead:eq:tail}, checks equality of the entire rational vector, and verifies these strict inequalities. No floating-point value, extrapolation, timing measurement, or numerical tolerance participates in the certificate.
\begin{writenote}[5 October 2026]
The write reran the certificate verifier on a copy of the shipped files (\file{code/144-leading-companion-verify_certificate.py} with \texttt{--fixture} pointing to \file{data/144-leading-companion-fixtures-rate_certificate.json}; Python~3.14 on Windows): it exits with status~0 and reports the strict bounds $527/1000$ and $79/125$. The README of this report gives the route. Theorem~\ref{ppg:all:thm:rate} of Part~\ref{ppg:all:part} takes this bracket as its starting interval \eqref{ppg:all:eq:initialR} and proves a width contraction by the factor $313/400$ per update; no update has been run.
\end{writenote}

\section{Divergence and the real axis singularity}\label{ppg:lead:sec:abelian}
By absolute convergence and the bounded kernel, \eqref{ppg:lead:eq:normalized} gives
\begin{equation}\label{ppg:lead:eq:derivative}
 F'(z)=\sum_{i,j\ge0}\alpha_{i,j}u_iu_jz^{i+j}\qquad(0\le z<R).
\end{equation}

\begin{lemma}[Divergence at the radius]\label{ppg:lead:lem:divergence}
$F(z)\to\infty$ as $z\uparrow R$.
\end{lemma}
\begin{proof}
If the positive boundary sum $F(R)$ were finite, choose $N$ such that $4RF(R)/N<1$. Then $4RP_N(R)/N<1$, so continuity gives $x>R$ with $4xP_N(x)/N<1$. The Catalan majorant \eqref{ppg:lead:eq:majorant} converges at $x$ and dominates $F$, contradicting its radius. This argument does not assume a singularity type or analytic continuation.
\end{proof}

\begin{theorem}[Leading pole]\label{ppg:lead:thm:pole}
As $z\uparrow R$,
\begin{equation}\label{ppg:lead:eq:pole}
 F(z)\sim\frac1{a(R-z)}=\frac A{1-z/R}.
\end{equation}
\end{theorem}
\begin{proof}
Divide \eqref{ppg:lead:eq:derivative} by $F(z)^2$. The quotient is the average of $\alpha_{I,J}$ under two independent distributions with masses $u_iz^i/F(z)$. For every fixed $M$, the probability that $I<M$ tends to zero by Lemma~\ref{ppg:lead:lem:divergence}, and the same holds for $J$. Boundedness of $\alpha$ and its uniform convergence to $a$ as $\min(i,j)\to\infty$ therefore give $F'/F^2\to a$. Since $1/F$ tends to zero,
\[
 \frac1{F(z)}=\int_z^R\frac{F'(t)}{F(t)^2}\,dt\sim a(R-z).
\]
The reciprocal gives \eqref{ppg:lead:eq:pole}.
\end{proof}

\begin{theorem}[Logarithmic second term]\label{ppg:lead:thm:log}
As $z\uparrow R$, \eqref{ppg:lead:eq:summaryF} holds. Equivalently,
\begin{equation}\label{ppg:lead:eq:xlog}
 F(Rx)=\frac A{1-x}+\frac A4\log(1-x)+O(1)\qquad(x\uparrow1).
\end{equation}
\end{theorem}
\begin{proof}
Write $\delta=R-z$ and introduce the positive weighted series
\begin{align*}
 J(z)&=\sum_{n\ge0}\frac{u_nz^n}{n+1}=\frac1z\int_0^zF(t)\,dt,\\
 K(z)&=\sum_{n\ge0}\frac{u_nz^n}{(n+1)^2}=\frac1z\int_0^zJ(t)\,dt,\\
 T_0(z)&=\sum_{i,j\ge0}\frac{u_iu_jz^{i+j}}{i+j+1}
         =\frac1z\int_0^zF(t)^2\,dt.
\end{align*}
At zero the integral formulas use their removable values. By \eqref{ppg:lead:eq:kerror},
\[
 F'=aF^2+O(FJ).
\]
The leading pole implies $J=O(\log(1/\delta))$. Divide by $F^2$, integrate the reciprocal derivative to $R$, and invert to get successively
\begin{align}
 \frac{F'}{F^2}&=a+O(\delta\log(1/\delta)),\notag\\
 \frac1F&=a\delta+O(\delta^2\log(1/\delta)),\notag\\
 F&=\frac1{a\delta}+O(\log(1/\delta)).\label{ppg:lead:eq:roughlog}
\end{align}
The error in \eqref{ppg:lead:eq:roughlog} is integrable near $R$. Therefore
\begin{equation}\label{ppg:lead:eq:Jfine}
 J(z)=\frac1{aR}\log(1/\delta)+O(1),\qquad K(z)=O(1),\qquad T_0(z)=O(F(z)).
\end{equation}
The exact denominator structure in \eqref{ppg:lead:eq:kexpansion} now yields
\begin{equation}\label{ppg:lead:eq:refinedDE}
 F'=aF^2+\frac a2FJ-\frac{5a}{12}T_0+O(FK)
    =aF^2+\frac a2FJ+O(F).
\end{equation}
Indeed the two first-reciprocal terms each produce $FJ$, the common denominator produces $T_0$, and the two squared-reciprocal errors are bounded by $20aFK$.
Equations~\eqref{ppg:lead:eq:roughlog} and \eqref{ppg:lead:eq:Jfine} give
\[
 \frac JF=\frac\delta R\log(1/\delta)+O(\delta),\qquad
 \frac{F'}{F^2}=a+\frac a{2R}\delta\log(1/\delta)+O(\delta).
\]
Integrate to the boundary, using
$\int_0^\delta t\log(1/t)\,dt=(\delta^2/2)\log(1/\delta)+\delta^2/4$:
\[
 \frac1F=a\delta+\frac a{4R}\delta^2\log(1/\delta)+O(\delta^2).
\]
Inversion gives \eqref{ppg:lead:eq:summaryF}; its residual quadratic inversion error is $O(\delta\log^2(1/\delta))=o(1)$. Replacing $\delta$ by $R(1-x)$ changes the logarithm only by an additive constant, proving \eqref{ppg:lead:eq:xlog}.
\end{proof}

\section{Positive averages and the original generating function}\label{ppg:lead:sec:tauberian}
\begin{theorem}[Cumulative asymptotics]\label{ppg:lead:thm:averages}
The laws \eqref{ppg:lead:eq:summaryV}, \eqref{ppg:lead:eq:summaryB}, and \eqref{ppg:lead:eq:summaryG} hold without any monotonicity assumption on the probabilities or normalized coefficients.
\end{theorem}
\begin{proof}
Set $v_n=u_nR^n\ge0$ and $H(s)=\sum_{n\ge0}v_ne^{-ns}=F(Re^{-s})$. By the leading pole, $H(s)\sim A/s$ as $s\downarrow0$. We include a positive-measure proof of the necessary Tauberian step.

For real $x>0$, put a mass $x^{-1}v_ne^{-n/x}$ at the point $y=e^{-n/x}$ of $[0,1]$, obtaining a finite measure $\eta_x$. For every integer $k\ge0$,
\[
 \int_0^1y^k\,d\eta_x(y)=\frac1xH((k+1)/x)\longrightarrow\frac A{k+1}.
\]
These are the moments of $A$ times Lebesgue measure on $[0,1]$. The $k=0$ case bounds total masses. Polynomial approximation of continuous functions proves weak convergence to that measure. The bounded function equal to $1/y$ on $[e^{-1},1]$ and zero elsewhere has only limiting-null discontinuities, so continuous upper and lower approximations give
\[
 \frac1x\sum_{n\le x}v_n
 =\int_{e^{-1}}^1\frac1y\,d\eta_x(y)
 \longrightarrow A\int_{e^{-1}}^1\frac{dy}{y}=A.
\]
This proves \eqref{ppg:lead:eq:summaryV}.

Let $V(t)=\sum_{n\le t}v_n\sim At$. Discrete summation by parts yields
\[
 \sum_{n=0}^Nv_n\sqrt{n+1}
 =\sqrt{N+1}V(N)-\sum_{n=0}^{N-1}V(n)(\sqrt{n+2}-\sqrt{n+1}).
\]
The second term is asymptotic to $(A/3)N^{3/2}$, since its summand is asymptotic to $(A/2)\sqrt n$. This proves \eqref{ppg:lead:eq:summaryB}.

Finally write $B_0(t)=\sum_{n\le t}W_nR^n$. Positivity and termwise integration give the exact identity
\[
 G(Re^{-s})=s\int_0^\infty e^{-st}B_0(t)\,dt.
\]
Since $B_0(t)\sim(2A/3)t^{3/2}$ and $B_0(t)\le C(1+t)^{3/2}$, substitution $y=st$ and dominated convergence give
\[
 G(Re^{-s})\sim\frac{2A}{3}\Gamma(5/2)s^{-3/2}
 =A\Gamma(3/2)s^{-3/2}.
\]
For $0<s\le1$ the global dominating function after multiplying by $s^{3/2}$ is $C(1+y)^{3/2}e^{-y}$, which also controls the atom at $n=0$. Now $1-e^{-s}\sim s$ and $A\Gamma(3/2)=\pi/(\sqrt2 R)$, proving \eqref{ppg:lead:eq:summaryG}.
\end{proof}

The positivity in this proof controls cumulative mass; it does not by itself control individual fluctuations. For a warning matching even the logarithmic term, take $0<d<3/4$, set $c_0=A(1+d)$, and let
\[
 c_n=A\{1+d(-1)^n-1/(4n)\}\quad(n\ge1).
\]
These coefficients are positive and have generating function
\[
 \sum_{n\ge0}c_nx^n=\frac A{1-x}+\frac A4\log(1-x)+\frac{Ad}{1+x},
\]
but do not converge to $A$. This is not a solution of the placement recurrence. It shows why a coefficient-level conclusion must use additional structure, as the next section does for a log-scale law.

\section{An individual logarithmic square root law}\label{ppg:lead:sec:subpower}
The cumulative law can be fed back into a symmetric positive recurrence. The following abstract lemma supplies a coefficient-level result weaker than an equivalent but stronger than an exponential rate.

\begin{lemma}[Subpower envelope from a half split]\label{ppg:lead:lem:subpower}
Suppose $x_n>0$ satisfies
\[
 nx_n=\sum_{i+j=n-1}\beta_{i,j}x_ix_j\qquad(n\ge1),
\]
where $\beta$ is symmetric, nonnegative and uniformly bounded,
$\beta_{i,j}\to1$ uniformly as $\min(i,j)\to\infty$, and
$S_N=\sum_{n=0}^Nx_n\sim N$. For every $p>0$ there exist constants $0<c_p\le C_p<\infty$ such that
\begin{equation}\label{ppg:lead:eq:subpower}
 c_p(n+1)^{-p}\le x_n\le C_p(n+1)^p\qquad(n\ge0).
\end{equation}
In particular $\log x_n/\log n\to0$.
\end{lemma}
\begin{proof}
Let $h_n=\lfloor(n-1)/2\rfloor$ and assign weight $w_{n,i}=2$ for $0\le i\le h_n$, except that the central index $i=(n-1)/2$, when integral, has weight~1. Symmetry gives
\[
 nx_n=\sum_{i=0}^{h_n}w_{n,i}\beta_{i,n-1-i}x_ix_{n-1-i}.
\]
For fixed real $q$, put
\[
 L_n(q)=\frac1n\sum_{i=0}^{h_n}w_{n,i}\beta_{i,n-1-i}x_i(1-i/n)^q.
\]
We first show
\begin{equation}\label{ppg:lead:eq:halfmultipliers}
 L_n(q)\longrightarrow L(q)=2\int_0^{1/2}(1-t)^q\,dt.
\end{equation}
The positive measures $n^{-1}\sum_{i\le n/2}x_i\delta_{i/n}$ converge weakly on $[0,1/2]$ to Lebesgue measure, since their cumulative masses converge to $t$ for each $0<t\le1/2$ by $S_{\lfloor nt\rfloor}/n\to t$, and the mass at zero vanishes. Furthermore
$x_k=S_k-S_{k-1}=o(k)$, so removing a boundary index or replacing the doubled central weight by~1 changes the scaled sum by $o(1)$. The function $(1-t)^q$ is bounded and continuous on this interval.

To replace $\beta$ by~1, fix $M$. The indices $i<M$ contribute $O_M(1/n)$, using boundedness of the kernel and the fixed finite prefix. On the rest of the half sum, the other index is at least $(n-1)/2$; uniform kernel convergence therefore makes $|\beta-1|$ arbitrarily small when $M$ and then $n$ are large. The scaled sum of $x_i$ remains bounded. This proves \eqref{ppg:lead:eq:halfmultipliers}.

For $p>0$, strict inequalities on $0<t\le1/2$ give
$L(p)<1<L(-p)$. Choose $N_0\ge1$ so that $L_n(p)<1$ and $L_n(-p)>1$ for all $n\ge N_0$. Choose $C_p$ to dominate $x_k/(k+1)^p$ on $0\le k<N_0$. If the upper bound holds at earlier indices, then $j=n-1-i<n$ and $j+1=n-i$, hence
\[
 x_n\le C_p n^pL_n(p)\le C_p n^p\le C_p(n+1)^p.
\]
Induction proves the upper bound. Choose $c_p>0$ below all $x_k(k+1)^p$ on the same finite prefix. The parallel induction gives
\[
 x_n\ge c_p n^{-p}L_n(-p)\ge c_p n^{-p}\ge c_p(n+1)^{-p}.
\]
This proves \eqref{ppg:lead:eq:subpower}. Take logarithms, divide by $\log n$, and then let $p\downarrow0$.
\end{proof}

\begin{theorem}[Individual logarithmic exponent]\label{ppg:lead:thm:subpower}
For the placement probabilities,
\begin{equation}\label{ppg:lead:eq:summarysubpower}
 W_n=r_*^n n^{1/2+o(1)},\qquad
 \frac{\log W_n-n\log r_*}{\log n}\longrightarrow\frac12.
\end{equation}
More precisely, with
\[
 x_n=\frac{u_nR^n}{A},
\]
the two-sided bound \eqref{ppg:lead:eq:subpower} holds for every $p>0$.
\end{theorem}
\begin{proof}
The normalized recurrence becomes
$nx_n=\sum_{i+j=n-1}(\alpha_{i,j}/a)x_ix_j$, because $AR=1/a$.
The kernel is symmetric by reflection in the beta integral and satisfies all the kernel hypotheses of Lemma~\ref{ppg:lead:lem:subpower} by Theorem~\ref{ppg:lead:thm:kernel}. The cumulative law gives $\sum_{n\le N}x_n\sim N$. Apply the lemma and use
$W_n=A\sqrt{n+1}\,r_*^nx_n$. The fixed factor $A$ is absorbed in $n^{o(1)}$.
\end{proof}

The subpower theorem alone supplies no constant-amplitude equivalent. The next section uses it to bound a coefficientwise error, then uses stronger kernel positivity to prove $x_n\to1$.


\section{The pointwise equivalent and the ratio limit}\label{ppg:lead:sec:pointwise}
We now use additional positivity of the normalized kernel to obtain the coefficient-level bridge. Set
\[
 X(z)=\sum_{n\ge0}x_nz^n=\frac{F(Rz)}A,\qquad
 \beta_{i,j}=\frac{\alpha_{i,j}}a.
\]
The radius of $X$ is~1. The results already proved give
\begin{equation}\label{ppg:lead:eq:Xreal}
 X(r)=\frac1{1-r}+H(r),\qquad
 |H(r)|\le C\left(1+\log\frac1{1-r}\right)\quad(r\uparrow1).
\end{equation}
No coefficient equivalent has been used to establish this radial estimate.

\begin{lemma}[Positive and small coefficient defect]\label{ppg:lead:lem:defect}
Define
\[
 e_m=\sum_{i+j=m}(\beta_{i,j}-1)x_ix_j,\qquad
 E(z)=\sum_{m\ge0}e_mz^m.
\]
Then $e_m\ge0$, and for each $\varepsilon>0$,
\begin{equation}\label{ppg:lead:eq:smallE}
 e_m=O_\varepsilon((m+1)^\varepsilon),\qquad
 X'(z)=X(z)^2+E(z)\quad(|z|<1).
\end{equation}
In particular $e_{n-1}/n\to0$.
\end{lemma}
\begin{proof}
Theorem~\ref{ppg:lead:thm:stronglower} gives $\beta_{i,j}>1$. The normalized recurrence and absolute convergence give the differential identity. By the inverse-index kernel bound and symmetry,
\[
 0\le e_m\le C\sum_{i+j=m}
 \left(\frac1{i+1}+\frac1{j+1}\right)x_ix_j
 =2C\sum_{i=0}^m\frac{x_i}{i+1}x_{m-i}.
\]
For any fixed $p>0$, the subpower envelope bounds this by
\[
 2CC_p^2(m+1)^p\sum_{i=0}^m(i+1)^{p-1}
 =O_p((m+1)^{2p}).
\]
Taking $p=\varepsilon/2$ proves the claim. This is a coefficientwise error estimate, stronger than a bound for $E$ only on the positive real axis.
\end{proof}

\begin{lemma}[Complex Abelian estimate in a fixed cone]\label{ppg:lead:lem:cone}
As $z\to1$ inside $|z|<1$, uniformly on every region
$|1-z|\le K(1-|z|)$ with fixed $K$,
\begin{equation}\label{ppg:lead:eq:cone}
 X(z)=\frac1{1-z}+o\bigl((1-|z|)^{-1}\bigr).
\end{equation}
\end{lemma}
\begin{proof}
The cumulative law is $S_n=\sum_{k\le n}x_k\sim n$. Abel summation gives
\[
 X(z)-\frac1{1-z}=(1-z)\sum_{n\ge0}\{S_n-(n+1)\}z^n.
\]
For any $\eta>0$, the summand in braces has magnitude at most $\eta(n+1)$ past a fixed index. The corresponding tail is bounded by
$\eta|1-z|/(1-|z|)^2\le\eta K/(1-|z|)$.
After multiplication by $1-|z|$, the finite prefix tends to zero uniformly. Let $\eta\downarrow0$.
\end{proof}

\begin{lemma}[Uniform modulus bound]\label{ppg:lead:lem:modulus}
There is a constant $C$ such that for every $r$ sufficiently close to~1 and every $-\pi\le\theta\le\pi$,
\begin{equation}\label{ppg:lead:eq:modulus}
 |X(re^{i\theta})|\le\frac C{1-r+|\theta|}
       +C\left(1+\log\frac1{1-r}\right)^2.
\end{equation}
\end{lemma}
\begin{proof}
We first obtain an initial angular gap at a radius depending on the angle. Put $t=|\theta|$. For $t\downarrow0$, take $\rho=1-t$ in Lemma~\ref{ppg:lead:lem:cone}. Since
$1-\rho e^{i\theta}=t-i\theta+O(t^2)$ and $tX(\rho)\to1$,
\[
 \frac{|X(\rho e^{i\theta})|}{X(\rho)}\longrightarrow\frac1{\sqrt2}
\]
for either sign of $\theta$. Choose $0<t_0<1/2$ small enough that this ratio is at most some fixed $q<1$ for $0<t\le t_0$, and that \eqref{ppg:lead:eq:Xreal} holds for $r\ge1-t_0$. For $t_0\le|\theta|\le\pi$, keep $\rho=1-t_0$. Strict triangle inequality holds because $x_0,x_1>0$: equality would require the constant and linear terms to have the same argument. Compactness gives a uniform strict gap on this angle set. Thus, putting $\tau=\min(|\theta|,t_0)$ and $\rho=1-\tau$, there are constants independent of $\theta\ne0$ such that
\begin{equation}\label{ppg:lead:eq:initialgap}
 \frac c\tau\le X(\rho)-|X(\rho e^{i\theta})|\le\frac C\tau.
\end{equation}

Fix $\theta\ne0$ and let $h(r)=|X(re^{i\theta})|$ for $\rho\le r<1$. The nonnegative coefficients of $E$ give
$|E(re^{i\theta})|\le E(r)$. On compact intervals $h$ is absolutely continuous, and the differential identity implies
\[
 h'(r)\le h(r)^2+E(r)
\]
almost everywhere. At a zero of $X(re^{i\theta})$, its upper right Dini derivative is also at most $E(r)$, so zeros cause no exception to scalar comparison. Let $y$ solve
\[
 y'=y^2+E(r),\qquad y(\rho)=h(\rho).
\]
Local existence and uniqueness apply below~1. The real function $X(r)$ solves the same equation and has a larger starting value. Uniqueness prevents crossing, and $0\le y<X$ extends the solution to every $r<1$, since $X$ is bounded on compact subintervals. Scalar comparison, or Gronwall's inequality for $(h-y)_+$, gives $h\le y$.

The positive gap $d=X-y$ and its reciprocal $v=1/d$ satisfy the exact equations
\begin{equation}\label{ppg:lead:eq:gapODE}
 d'=2Xd-d^2,\qquad v'=1-2Xv.
\end{equation}
Let $\delta=1-r$ and $\mathcal L(\delta)=1+\log(1/\delta)$. By \eqref{ppg:lead:eq:initialgap}, $c\tau\le v(\rho)\le C\tau$. With
\[
 P(r)=\exp\left(-2\int_\rho^rX(s)\,ds\right),
\]
variation of constants gives
\begin{equation}\label{ppg:lead:eq:vsolution}
 v(r)=P(r)\left\{v(\rho)+\int_\rho^rP(s)^{-1}\,ds\right\}.
\end{equation}
By \eqref{ppg:lead:eq:Xreal},
\[
 P(r)=\left(\frac\delta\tau\right)^2
          \exp\left(-2\int_\rho^rH(s)\,ds\right).
\]
The integral of $|H|$ is bounded uniformly over all these intervals by
$C\int_0^{t_0}(1+\log(1/u))\,du$. Hence
$c(\delta/\tau)^2\le P(r)\le C(\delta/\tau)^2$.
The positive bracket in \eqref{ppg:lead:eq:vsolution} is comparable to
\[
 \tau+\tau^2(1/\delta-1/\tau)=\tau^2/\delta,
\]
including $r$ arbitrarily close to $\rho$. Therefore
\begin{equation}\label{ppg:lead:eq:vcomparable}
 c\delta\le v(r)\le C\delta.
\end{equation}

A sharper difference estimate is essential. Put $k(r)=v(r)-\delta$. Since $\delta'=-1$, equation~\eqref{ppg:lead:eq:gapODE} gives
\[
 k'=-2k/\delta-2Hv,
\]
and thus
\[
 k(r)=\left(\frac\delta\tau\right)^2k(\rho)
       -2\delta^2\int_\rho^r\frac{H(s)v(s)}{(1-s)^2}\,ds.
\]
Use $|k(\rho)|\le C\tau$, \eqref{ppg:lead:eq:vcomparable}, and the radial bound on $H$ to obtain
\begin{equation}\label{ppg:lead:eq:vdifference}
 |v(r)-\delta|\le C\delta^2
 \left\{\tau^{-1}+\int_\delta^\tau\frac{1+\log(1/u)}u\,du\right\}
 \le C\delta^2\{\tau^{-1}+\mathcal L(\delta)^2\}.
\end{equation}
Finally,
\[
 h\le y=X-1/v=H+\frac{v-\delta}{v\delta}
 \le C/\tau+C\mathcal L(\delta)^2.
\]
No sign assumption on $H$ or $v-\delta$ is used. For small angles with $\delta\le|\theta|$, $1/|\theta|\le2/(\delta+|\theta|)$. If $\delta>|\theta|$, use instead $|X(re^{i\theta})|\le X(r)\le C/\delta$. For $|\theta|\ge t_0$, the constant $C/t_0$ is absorbed in the logarithmic term. At $\theta=0$ the radial bound suffices. This proves \eqref{ppg:lead:eq:modulus} with uniform constants.
\end{proof}

\begin{theorem}[Leading equivalent and consecutive ratio]\label{ppg:lead:thm:pointwise}
The normalized coefficients satisfy $x_n\to1$. Consequently,
\begin{equation}\label{ppg:lead:eq:pointwise}
 W_n\sim A\sqrt n\,r_*^n=r_*\sqrt{2\pi n}\,r_*^n,
 \qquad \frac{W_{n+1}}{W_n}\longrightarrow r_*.
\end{equation}
In particular $W_n$ is eventually strictly decreasing, although monotonicity of the consecutive ratios is not asserted.
\end{theorem}
\begin{proof}
The exact coefficient identity from \eqref{ppg:lead:eq:smallE} is
\[
 x_n=\frac1n[z^{n-1}]X(z)^2+\frac{e_{n-1}}n.
\]
The second term tends to zero. Compare the first with $M(z)^2=(1-z)^{-2}$, whose $z^{n-1}$ coefficient is exactly~$n$.

Set $r_n=1-1/n$ for $n\ge2$. Cauchy's coefficient formula gives
\begin{align}\label{ppg:lead:eq:cauchy}
 \frac1n[z^{n-1}](X^2-M^2)
 =\frac{r_n^{-(n-1)}}{2\pi n}\int_{-\pi}^{\pi}
 \{X(r_ne^{i\theta})^2-M(r_ne^{i\theta})^2\}
 e^{-i(n-1)\theta}\,d\theta.
\end{align}
The radial factor is bounded. For each fixed finite $T$, Lemma~\ref{ppg:lead:lem:cone} yields uniformly on $|t|\le T$,
\[
 \frac1nX(r_ne^{it/n})\longrightarrow\frac1{1-it},\qquad
 \frac1nM(r_ne^{it/n})\longrightarrow\frac1{1-it}.
\]
Therefore their normalized square difference tends uniformly to zero. Substitution $\theta=t/n$ shows that the central contribution $|\theta|\le T/n$ to \eqref{ppg:lead:eq:cauchy} tends to zero.

On the complement, \eqref{ppg:lead:eq:modulus} and $(a+b)^2\le2a^2+2b^2$ give
\[
 \frac1n\int_{T/n<|\theta|\le\pi}|X(r_ne^{i\theta})|^2\,d\theta
 \le\frac C{1+T}+\frac{C(1+\log n)^4}{n}.
\]
The reference function has the same tail bound without the logarithmic term: indeed
$|1-re^{i\theta}|^2=(1-r)^2+4r\sin^2(\theta/2)$ and
$\sin(|\theta|/2)\ge|\theta|/\pi$ on this interval give
$|1-r_ne^{i\theta}|\ge c(n^{-1}+|\theta|)$.
Take the limsup in \eqref{ppg:lead:eq:cauchy} as $n\to\infty$ at fixed $T$, then let $T\to\infty$. This proves
$n^{-1}[z^{n-1}]X^2\to1$, and hence $x_n\to1$.

Undo the normalization to obtain the equivalent. The exact ratio identity
\[
 \frac{W_{n+1}}{W_n}
 =r_*\sqrt{\frac{n+2}{n+1}}\frac{x_{n+1}}{x_n}
\]
gives the ratio limit. Since $r_*<1$, these ratios are eventually less than~1, proving eventual decrease of $W_n$ itself.
\end{proof}

The proof extracts coefficients of $X^2$ through the differential equation, with an integrable square-modulus tail. It neither transfers a radial logarithm directly to individual coefficients nor assumes analytic continuation across the unit circle. The decisive extra inputs beyond the averaged law are the nonnegative defect and its coefficientwise $o(n)$ bound.
\begin{writenote}[5 October 2026]
This theorem proves Serra's Conjecture~6.10 (stated in his Section~1.3, printed page~7, and posed as his Open Problem~4, printed page~69), with $r_\infty=r_*\in(0,\frac34)$ since $79/125<3/4$, and it determines the constant: $C=r_*\sqrt{2\pi}$, not his $\sqrt{2\pi}$ (Remark~\ref{ppg:lead:rem:serra}). Serra's Remark~8.5 makes his Open Problems~1--3 logical prerequisites of Open Problem~4; the proof here uses none of them, and Open Problems~2 and~3 are false as stated (Remark~\ref{ppg:all:rem:op2}, Proposition~\ref{ppg:lead:prop:op3}). The convergence $x_n\to1$ is quantified in Part~\ref{ppg:log:part}: $x_n=1-\frac1{4n}+o(n^{-1})$ (Proposition~\ref{ppg:log:prop:first}), which answers the first question of Section~\ref{ppg:lead:sec:questions} asymptotically.
\end{writenote}
\section{Necessary exponent and corrected amplitude}\label{ppg:lead:sec:conditional}
\begin{theorem}[Every exponential-power equivalent]\label{ppg:lead:thm:conditional}
Suppose for fixed $C,r>0$ and $\gamma\in\mathbb R$ that
$W_n\sim Cn^\gamma r^n$. Then
$\gamma=1/2$, $r=r_*$, and $C=r_*\sqrt{2\pi}$.
\end{theorem}
\begin{proof}
Taking nth roots already forces $r=r_*$. From \eqref{ppg:lead:eq:kbound} and \eqref{ppg:lead:eq:alpha}, uniformly for all $n,k$,
\begin{equation}\label{ppg:lead:eq:Qupper}
 Q_{n,k}\le C_0\sqrt{\frac n{k(n-k+1)}}
\end{equation}
with, for example, $C_0=\sqrt2$. Moreover, for each fixed $0<\delta<1/2$, \eqref{ppg:lead:eq:kerror} and \eqref{ppg:lead:eq:alpha} give uniformly on $\delta\le k/n\le1-\delta$,
\begin{equation}\label{ppg:lead:eq:Qbulk}
 Q_{n,k}=\frac{1+O_\delta(n^{-1})}{\sqrt{2\pi n(k/n)(1-k/n)}}.
\end{equation}
The assumed equivalent supplies a constant $D$ such that
$W_j\le D(j+1)^\gamma r^j$ for every $j\ge0$. Divide the recurrence by $n^\gamma r^n$ and use \eqref{ppg:lead:eq:Qupper}. Up to a constant independent of $n$, its right side is bounded by
\begin{equation}\label{ppg:lead:eq:alphasum}
 n^{-\gamma-1/2}\sum_{k=1}^n
 k^{\gamma-1/2}(n+1-k)^{\gamma-1/2}.
\end{equation}
For $\gamma<-1/2$, splitting the sum at its midpoint bounds the sum by $O(n^{\gamma-1/2})$, and \eqref{ppg:lead:eq:alphasum} is $O(n^{-1})$. For $\gamma=-1/2$, it is $O(n^{-1}\log n)$. For $-1/2<\gamma<1/2$, the sum is $O(n^{2\gamma})$, making \eqref{ppg:lead:eq:alphasum} $O(n^{\gamma-1/2})$. All three cases contradict convergence of the left side to $C>0$.

If $\gamma>1/2$, restrict to $\delta\le k/n\le1-\delta$. The assumed equivalent is uniform over both growing child indices there. Using \eqref{ppg:lead:eq:Qbulk}, the positive bulk in the normalized recurrence is at least a positive constant times $n^{\gamma-1/2}$, again contradicting convergence to $C$. Thus $\gamma=1/2$.

Now $W_j\le D\sqrt{j+1}r^j$. Each summand after division by $\sqrt n r^n$ is bounded by $C_0D^2/(rn)$ by \eqref{ppg:lead:eq:Qupper}. Hence the two endpoint regions contribute $O(\delta+1/n)$. In the compact bulk each normalized summand is, uniformly,
\[
 \frac{C^2}{r\sqrt{2\pi}\,n}(1+o(1)).
\]
Let $n\to\infty$ and then $\delta\downarrow0$. The recurrence gives
$C=C^2/(r\sqrt{2\pi})$, and positivity yields $C=r\sqrt{2\pi}$.
\end{proof}

\subsection{The source normalization}\label{ppg:lead:sec:sourcecorrection}
Equation~(140) of Serra~\cite{serra}, printed page~46, replaces the first child in the exact product $W_{k-1}W_{n-k}$ by $W_{xn}$ with $x=k/n$. Under an exponential-power equivalent, the exact exponential factor is $r^{n-1}$; the displayed unshifted integral instead produces $r^n$. The missing factor $r^{-1}$ changes the leading amplitude. Equation~(147) and Proposition~6.17, printed pages~47--48, consequently propose $\sqrt{2\pi}$ where the exact recurrence forces $r_*\sqrt{2\pi}$, now also established by the pointwise theorem. The same shift must be retained when formulating a rigorous integral approximation of the type requested in the source's Open Problem~3, equation~(200).

This discrepancy is an order-one normalization issue rather than a lower-order finite-size correction. The corrected amplitude agrees with the empirical amplitude recorded in OEIS when evaluated at its proposed numerical rate, but that numerical consistency is not used in any proof. We do not adopt the source's extrapolated digits as certified values. Theorem~\ref{ppg:lead:thm:conditional} gives an independent local balance check; Theorem~\ref{ppg:lead:thm:pointwise} supplies existence of the equivalent.

\begin{remark}[\wtag\ Serra's prefactor, 5 October 2026]\label{ppg:lead:rem:serra}
Serra's version~1 \cite{serra} asserts the prefactor $C=\sqrt{2\pi}\approx2.507$ in $W(n)\sim C\sqrt n\,r^n$ repeatedly: Proposition~6.17 with its display~(149) (printed page~48), the outcome of the formal computation of his Section~6.3 described in Section~\ref{ppg:lead:sec:sourcecorrection}; Remark~6.20(ii) (page~49) and Remark~8.4 (page~69), which restate it; and his abstract and Section~1.3 (page~7), which add that the prefactor is ``approached from below''. His Remark~6.19 (page~48) reconciles it with the regression value $C\approx1.46$ of his Section~6.2 through an effective prefactor $C_{\rm eff}(n)=C(1-b/\sqrt n+O(1/n))$ with $b>0$. For the game these statements are wrong:
\begin{enumerate}[label=(\alph*)]
\item $W_n/(\sqrt n\,r_*^n)\to A=r_*\sqrt{2\pi}$, and $1.32<A<1.59$; in particular the limit is not $\sqrt{2\pi}$.
\item $W_n/(\sqrt n\,r_*^n)=A\bigl(1+\frac1{4n}+O(n^{-2}\log n)\bigr)$. If $W_n/(\sqrt n\,r_*^n)=C(1-b/\sqrt n+O(1/n))$ for constants $C,b$, then $C=A$ and $b=0$; the limit is approached from above, not from below.
\item The exact balance is $C=C^2/(r\sqrt{2\pi})$, not Serra's~(147), $C=C^2/\sqrt{2\pi}$; the factor $r^{-1}$ is lost when the child sizes $k-1$ and $n-k$, whose sum is $n-1$, are replaced by $xn$ and $(1-x)n$, whose sum is $n$.
\end{enumerate}
\end{remark}
\begin{proof}
(a) The limit is Theorem~\ref{ppg:lead:thm:pointwise}. By \eqref{ppg:lead:eq:numeric} and $2.5066<\sqrt{2\pi}<2.5067$, $A>0.527\cdot2.5066>1.32$ and $A<0.632\cdot2.5067<1.59$, while $\sqrt{2\pi}>2.5$.

(b) Multiply \eqref{ppg:log:eq:main} of Theorem~\ref{ppg:log:thm:main} by $A$. Then $W_n/(\sqrt n\,r_*^n)-A=\frac A{4n}+O(n^{-2}\log n)$, which is positive for all large $n$. If also $W_n/(\sqrt n\,r_*^n)=C(1-b/\sqrt n+O(1/n))$, letting $n\to\infty$ gives $C=A$, and then $\sqrt n\,(W_n/(\sqrt n\,r_*^n)-A)=-Ab+O(n^{-1/2})$ while the left side is $O(n^{-1/2})$; so $Ab=0$ and $b=0$.

(c) Insert $W_j=C\sqrt j\,r^j(1+o(1))$ in the recurrence \eqref{ppg:lead:eq:recurrence}. In the bulk $\delta\le k/n\le1-\delta$, \eqref{ppg:lead:eq:Qbulk} gives the summand $\frac1nC^2\sqrt{(k-1)(n-k)}\,r^{n-1}/\sqrt{2\pi n(k/n)(1-k/n)}\,(1+o(1))$; dividing by $\sqrt n\,r^n$ and summing, as in the proof of Theorem~\ref{ppg:lead:thm:conditional}, gives $C=C^2/(r\sqrt{2\pi})$ in the limit. Serra's Steps~2--3 (his (141)--(147)) perform the same computation with $W(xn)W((1-x)n)$, whose exponential factor is $r^{xn}r^{(1-x)n}=r^n$ instead of $r^{(k-1)+(n-k)}=r^{n-1}$, and so obtain $C=C^2/\sqrt{2\pi}$.
\end{proof}
Numerically, Serra's regression value $1.46$ and Kot\v{e}\v{s}ovec's OEIS amplitude $1.4424953427$ are close to $A$, which is $1.44249534\ldots$ at the posted digits of $r_*$; no proof uses these values.

\begin{proposition}[\wtag\ Serra's integral equation~(200) read through the profile, 5 October 2026]\label{ppg:lead:prop:op3}
Let $\widetilde W:[0,\infty)\to(0,\infty)$ be measurable, with $\widetilde W(t)\le K\sqrt{t+1}\,r_*^t$ for all $t\ge0$ and some constant $K$, and $\widetilde W(t)\sim A\sqrt t\,r_*^t$ as $t\to\infty$. Examples are Serra's profile $\widetilde W(t)=A\sqrt t\,r_*^t$ of his Section~6.3, Step~2, and the interpolation $\widetilde W(t)=W_m\,r_*^{t-m}\sqrt{(t+1)/(m+1)}$, $m=\lfloor t\rfloor$, which equals $W_m$ at every integer $m$. Put
\[
 I_n=\int_0^1\frac{\widetilde W(xn)\,\widetilde W((1-x)n)}{\sqrt{2\pi nx(1-x)}}\,\dd x .
\]
Then $I_n/W_n\to r_*$. Hence if $W(xn)$ in Serra's equation~(200) (his Open Problem~3, printed page~68) is read as $\widetilde W(xn)$, its error term $E(n)=W_n-I_n$ satisfies $E(n)/W_n\to1-r_*$, and $\frac{46}{125}<1-r_*<\frac{473}{1000}$: the equation does not hold with $E(n)=o(W(n))$, let alone with Serra's $E(n)=O(W(n)/n)$.
\end{proposition}
\begin{proof}
Put $\varphi(t)=\widetilde W(t)/(\sqrt{t+1}\,r_*^t)$, so that $0<\varphi\le K$ and $\varphi(t)\to A$ as $t\to\infty$. (For the first example $\varphi(t)=A\sqrt{t/(t+1)}$. For the second, $\varphi(t)=W_mR^m/\sqrt{m+1}=Ax_m$ with $m=\lfloor t\rfloor$, which is bounded and tends to $A$ by Theorem~\ref{ppg:lead:thm:pointwise}; so both examples satisfy the hypotheses.) Since $r_*^{xn}r_*^{(1-x)n}=r_*^n$,
\[
 \frac{I_n}{\sqrt n\,r_*^n}
 =\frac1{\sqrt{2\pi}}\int_0^1\varphi(xn)\,\varphi((1-x)n)\,g_n(x)\,\dd x,
 \qquad g_n(x)=\sqrt{\frac{(x+1/n)(1-x+1/n)}{x(1-x)}} .
\]
For $n\ge1$ and $0<x<1$, $(x+1/n)(1-x+1/n)\le(x+1)(2-x)\le\frac94$, so the integrand is at most $\frac32K^2(x(1-x))^{-1/2}$, whose integral over $(0,1)$ is $\frac32K^2\pi$. For fixed $x\in(0,1)$, $xn\to\infty$ and $(1-x)n\to\infty$, so the integrand tends to $A^2$. By dominated convergence $I_n/(\sqrt n\,r_*^n)\to A^2/\sqrt{2\pi}=Ar_*$. By Theorem~\ref{ppg:lead:thm:pointwise}, $W_n/(\sqrt n\,r_*^n)\to A$, hence $I_n/W_n\to r_*$ and $E(n)/W_n\to1-r_*$. The bounds follow from \eqref{ppg:lead:eq:numeric}.
\end{proof}
The proposition concerns only readings of $W(xn)$ that follow the asymptotic profile at non-integer arguments; it says nothing about, for example, $W(\lfloor xn\rfloor)$, for which $\lfloor xn\rfloor+\lfloor(1-x)n\rfloor=n-1$ whenever $xn\notin\mathbb Z$, so that the shift of Section~\ref{ppg:lead:sec:sourcecorrection} is built in. As the delivered text says, a rigorous integral approximation must retain that shift.

\section{Two different inverse questions}\label{ppg:lead:sec:inverses}
\begin{theorem}[Continuous generating-function inverse]\label{ppg:lead:thm:GFinverse}
For every $Y\ge1$ there is a unique $z(Y)\in[0,R)$ with $F(z(Y))=Y$, and \eqref{ppg:lead:eq:summaryinverse} holds.
\end{theorem}
\begin{proof}
Positivity of the coefficients and $u_1>0$ imply that $F$ is continuous and strictly increasing on $[0,R)$. It starts at~1 and diverges at $R$, so the inverse is defined exactly. Put $\delta=R-z(Y)$. The leading pole gives $\delta\sim1/(aY)$, while the logarithmic expansion gives
\[
 aY\delta=1-\frac\delta{4R}\log(1/\delta)+O(\delta).
\]
Thus $\delta=1/(aY)+O(\log Y/Y^2)$ and
$\delta\log(1/\delta)=(\log Y)/(aY)+O(1/Y)$. Substituting and dividing by $aY$ gives
\[
 \delta=\frac1{aY}-\frac1{4a^2R}\frac{\log Y}{Y^2}+O(Y^{-2}).
\]
Subtract from $R$. Additive constants such as $\log a$ are absorbed by the $O(Y^{-2})$ term.
\end{proof}

\begin{theorem}[Ceiling safe probability threshold inverse]\label{ppg:lead:thm:threshold}
Let $L=\log(1/\varepsilon)$, $\lambda=\log(1/r_*)$, and
$N(\varepsilon)=\min\{n\ge0:W_n\le\varepsilon\}$. For sufficiently small $\varepsilon>0$, let $t_\varepsilon$ be the large real solution of
$A\sqrt t\,e^{-\lambda t}=\varepsilon$. Then
\begin{equation}\label{ppg:lead:eq:lambert}
 t_\varepsilon=-\frac{\operatorname W_{-1}(-2\lambda\varepsilon^2/A^2)}{2\lambda},
\end{equation}
where $\operatorname W_{-1}$ is the real branch taking values at most $-1$ of the inverse of $w\mapsto we^w$. There exists $\eta_\varepsilon\ge0$ with $\eta_\varepsilon\to0$ such that
\begin{equation}\label{ppg:lead:eq:thresholdceil}
 \lceil t_\varepsilon-\eta_\varepsilon\rceil
 \le N(\varepsilon)\le
 \lceil t_\varepsilon+\eta_\varepsilon\rceil.
\end{equation}
The same conclusion holds, with another vanishing nonnegative uncertainty, for
\begin{equation}\label{ppg:lead:eq:thresholdcenter}
 c_\varepsilon=
 \frac{L+\tfrac12\log L+\log A-\tfrac12\log\lambda}{\lambda}
\end{equation}
in place of $t_\varepsilon$. In particular
$N(\varepsilon)=c_\varepsilon+O(1)$.
\end{theorem}
\begin{proof}
First the nth-root limit alone gives the leading law
$N(\varepsilon)\sim L/\lambda$, without assuming monotonicity. Indeed, for any fixed $0<\xi<\lambda$, all sufficiently large $n$ satisfy
\[
 e^{-(\lambda+\xi)n}\le W_n\le e^{-(\lambda-\xi)n}.
\]
A fixed positive prefix contributes no crossing when $\varepsilon$ is small. The lower envelope excludes all indices below $L/(\lambda+\xi)$, whereas $\lceil L/(\lambda-\xi)\rceil$ is a crossing. Let $\xi\downarrow0$.

The model $A\sqrt t\,e^{-\lambda t}$ is strictly decreasing for $t>1/(2\lambda)$ and tends to zero. Its large root exists uniquely for small $\varepsilon$. Squaring the equation and putting $w=-2\lambda t$ gives $we^w=-2\lambda\varepsilon^2/A^2$; the branch $\operatorname W_{-1}$ selects the large root, proving \eqref{ppg:lead:eq:lambert}.

Put $f(t)=\lambda t-\tfrac12\log t-\log A$, so that
$f(t_\varepsilon)=L$ and $f'(t)\ge\lambda/2$ for large $t$.
Theorem~\ref{ppg:lead:thm:pointwise} says
\[
 \log W_n=-f(n)+d_n,\qquad d_n\longrightarrow0.
\]
Since $f(t)/t\to\lambda$, $t_\varepsilon\sim L/\lambda$; the leading first-passage law therefore gives $N(\varepsilon)>t_\varepsilon/2$ eventually, excluding every earlier index. Define
\[
 s_\varepsilon=\sup_{n\ge\lfloor t_\varepsilon/2\rfloor}|d_n|,
 \qquad \eta_\varepsilon=4s_\varepsilon/\lambda+1/t_\varepsilon.
\]
Then $s_\varepsilon\to0$ and $\eta_\varepsilon>0$ tends to zero. If
$t_\varepsilon/2\le n<t_\varepsilon-\eta_\varepsilon$, the mean value theorem gives
\[
 \log W_n>-L+(\lambda/2)\eta_\varepsilon-s_\varepsilon>-L,
\]
so this index is not a crossing. At
$n=\lceil t_\varepsilon+\eta_\varepsilon\rceil$ the analogous upper bound gives
\[
 \log W_n\le-L-(\lambda/2)\eta_\varepsilon+s_\varepsilon<-L,
\]
so it is a crossing. These strict comparisons prove both ceilings in \eqref{ppg:lead:eq:thresholdceil}, even when a bracket endpoint is an integer. The squeeze does not require monotonicity of $W_n$.

Finally,
\[
 \lambda t_\varepsilon=L+\tfrac12\log t_\varepsilon+\log A
 =L+\tfrac12\log L+\log A-\tfrac12\log\lambda+o(1).
\]
Thus $t_\varepsilon=c_\varepsilon+o(1)$. Increase the uncertainty by
$|t_\varepsilon-c_\varepsilon|$ to transfer the bracket to $c_\varepsilon$.
\end{proof}

The uncertainty vanishes \emph{before rounding}. One cannot generally replace the ceiling bracket by $N(\varepsilon)=c_\varepsilon+o(1)$, since the first passage is integer-valued. The proof gives no effective rate for the vanishing uncertainty. A finer expansion of the model root alone would not control the still unquantified forward error in the pointwise equivalent.

\begin{remark}[\wtag\ Repository context and later refinements, 5 October 2026]\label{ppg:lead:rem:inverse}
(i) \emph{The Lambert root.} Squaring $A\sqrt t\,e^{-\lambda t}=\varepsilon$ and taking logarithms gives $\lambda t-\frac12\log t=\log(A/\varepsilon)$, the equation $aX+b\log X=L$ of \texttt{p0:thm:lambert-core} in the transseries volume \cite{tsvol} with $a=\lambda>0$, $b=-\frac12<0$ and $L=\log(A/\varepsilon)$. For $b<0$ the branch rule there gives the large root $X=\frac ba\,W_{-1}\bigl(\frac ab\,e^{L/b}\bigr)=-\frac1{2\lambda}W_{-1}(-2\lambda\varepsilon^2/A^2)$, which is \eqref{ppg:lead:eq:lambert}, whenever $L>L_c=\frac12(1+\log(2\lambda))$, that is for all sufficiently small $\varepsilon$.

(ii) \emph{The first passage.} By Theorem~\ref{ppg:lead:thm:pointwise}, $W_n$ is strictly decreasing for $n\ge n_1$, so $\log(1/W_n)$ is strictly increasing there, and once $\varepsilon<\min_{n<n_1}W_n$ the first passage $N(\varepsilon)$ is the integer staircase of \texttt{p0:def:three-inverses} for that sequence at $y=\log(1/\varepsilon)$. A bracket \eqref{ppg:lead:eq:thresholdceil} whose two ceilings agree determines it, as in the separation condition of \texttt{p0:thm:staircase}(2), and the mean-value step of the proof is the estimate of \texttt{p0:cor:forward-to-inverse}. The proof here uses neither the monotonicity nor an admissible interpolation: it compares $\log W_n$ with the smooth model $f$ directly. These are pointers to known mechanics, not new inverse theory.

(iii) \emph{Later refinements.} Theorem~\ref{ppg:log:thm:inverse} sharpens the uncertainty to $o(x^{-2})$ and Theorem~\ref{ppg:all:thm:threshold} to $O((1+\log x)^M/x^{M+1})$ at every fixed order. Theorem~\ref{ppg:all:thm:threshold} with $M=1$ implies \eqref{ppg:lead:eq:thresholdceil}: its leading root $x$ in \eqref{ppg:all:eq:x0} is $t_\varepsilon$ of \eqref{ppg:lead:eq:lambert}, its centre is $z_1=t_\varepsilon+b/(\lambda t_\varepsilon)$ with $b=\frac14$ by \eqref{ppg:all:eq:R123}, and with $\eta'_\varepsilon=\frac1{4\lambda t_\varepsilon}+C_1E_1(t_\varepsilon)\to0$ one has $t_\varepsilon-\eta'_\varepsilon\le z_1-C_1E_1(t_\varepsilon)$ and $z_1+C_1E_1(t_\varepsilon)=t_\varepsilon+\eta'_\varepsilon$, so the monotonicity of the ceiling turns \eqref{ppg:all:eq:ceiling} into \eqref{ppg:lead:eq:thresholdceil} with $\eta_\varepsilon=\eta'_\varepsilon$. The proof of Theorem~\ref{ppg:all:thm:threshold} does not use Theorem~\ref{ppg:lead:thm:threshold}.
\end{remark}

\section{Questions beyond these results}\label{ppg:lead:sec:questions}
The results establish a computable exponential rate, the full leading square-root equivalent, and convergence of the consecutive ratios, while also retaining a refined real-axis logarithm. The coefficient-level bridge is proved here through kernel positivity, the Riccati comparison, and square-modulus Cauchy extraction. Several concrete questions remain.

\paragraph{Quantifying the coefficient bridge.}
The proof gives $x_n\to1$ without an explicit rate. Can the central and tail estimates be sharpened enough to transfer the real-axis logarithmic term coefficientwise? A quantitative modulus estimate, a controlled remainder in the cumulative law, or stronger analytic information could help, but each requires a new argument.

\paragraph{Monotonicity of the ratios.}
The convergence $W_{n+1}/W_n\to r_*<1$ implies eventual decrease of $W_n$. It does not prove that the sequence of ratios is itself monotone, either eventually or at every index. That stronger property is not resolved here.

\paragraph{Pointwise corrections and logarithms.}
Formally, transferring \eqref{ppg:lead:eq:xlog} coefficientwise suggests
$u_nR^n=A-A/(4n)+\cdots$. Multiplication by $\sqrt{n+1}$ would then give
$W_n\sim A\sqrt n\,r_*^n(1+1/(4n)+\cdots)$. These are formal implications, not proved coefficient estimates. Higher kernel orders interact with endpoint contributions and integrated logarithms, so a pure power series in $1/n$ should not be presumed. Whether terms such as $\log n/n^2$ occur, and what remainders hold, is not established here. No full transseries is obtained here.

\paragraph{Efficient certified digits.}
The convergent enclosures prove computability but provide no useful a priori complexity or convergence rate. Sharper split inequalities, stronger majorants, or a proved coefficient remainder could make rigorous high precision practical. The pointwise theorem now identifies $r_*$ with the consecutive-ratio limit described in OEIS. It does not supply a certified rate of numerical convergence or validate its posted digits.

\paragraph{Other policies and universality.}
The arguments concern the stated proportional policy. They do not prove policy optimality or classify other policies. It is natural to ask which other positive size-dependent convolution kernels admit the same balanced-merge rate argument and pole mechanism, and which changes in the reciprocal-index kernel correction alter the logarithmic term. Any such extension needs explicit hypotheses and a separate proof.
\begin{writenote}[5 October 2026]
Status of these questions after Parts~\ref{ppg:log:part} and~\ref{ppg:all:part}. \emph{Quantifying the coefficient bridge}: answered asymptotically. Part~\ref{ppg:log:part} proves $y_n=O(\log n/n)$ by a tail contraction (Lemma~\ref{ppg:log:lem:tail}) and then $y_n=-\frac1{4n}+o(n^{-1})$ (Proposition~\ref{ppg:log:prop:first}); Part~\ref{ppg:all:part} gives every fixed order. Effective constants remain open (Section~\ref{ppg:sec:further}, question~\ref{ppg:fq:effective}). \emph{Monotonicity of the ratios}: eventual strict decrease is proved (Corollary~\ref{ppg:log:cor:monotonicity}); strict decrease at every index, Serra's Open Problem~1, is open (question~\ref{ppg:fq:monotone}). \emph{Pointwise corrections and logarithms}: answered. The formal suggestion $W_n\sim A\sqrt n\,r_*^n(1+\frac1{4n}+\cdots)$ is Theorem~\ref{ppg:log:thm:main}; a term $\log n/n^2$ does occur, with coefficient $-\frac7{540}$; and the expansion to every fixed order is Theorem~\ref{ppg:all:thm:all}. No transseries is obtained (question~\ref{ppg:fq:growth}). \emph{Efficient certified digits}: open; Theorem~\ref{ppg:all:thm:rate} adds an explicit contraction but no complexity bound and no run (question~\ref{ppg:fq:numerics}). \emph{Other policies and universality}: open (question~\ref{ppg:fq:policies}).
\end{writenote}

\section{Reproducible exact computation}\label{ppg:lead:app:reproduce}
The accompanying archive contains this LaTeX source, its PDF, a deterministic build script, and a Python-standard-library companion with exact fixtures and tests. The companion evaluates \eqref{ppg:lead:eq:poly} and \eqref{ppg:lead:eq:tail} independently through $n=100$, compares every rational probability, checks the prefix \eqref{ppg:lead:eq:prefix}, verifies integer-square-root floors, and proves \eqref{ppg:lead:eq:finiteLower}--\eqref{ppg:lead:eq:finiteUpper} by strict rational comparisons. Its fixture validation uses explicit checks that remain active under Python optimization. The README gives the exact commands and the output-containment rules.

The finite computation verifies the stated finite inputs and certificate inequalities. It does not mechanize the analytic proofs, including the pointwise theorem, or certify extrapolated decimal digits. The theoretical inputs linking those inequalities to $r_*$ are Theorems~\ref{ppg:lead:thm:rate}, \ref{ppg:lead:thm:kernel}, and \ref{ppg:lead:thm:uppercert}. The source manuscript and other copyrighted references are cited by link and are not included in the archive.
\begin{writenote}[5 October 2026]
In this directory the archive's files are shipped under prefixed names: the verifier, its tests and the archive builder in \file{code/} (\file{144-leading-companion-verify_certificate.py}, \file{144-leading-companion-test_verifier.py}, \file{144-leading-bundle.py}), the fixture as \file{data/144-leading-companion-fixtures-rate_certificate.json}, and the companion's README, schema and provenance notes and the archive's \file{SOURCES.md} at the root with the prefix \file{144-leading-}. The delivered PDF, README and checksum manifest are not shipped; the archive survives in commit \texttt{60f54ea06}. The README of this report gives rerun routes that work with the shipped names.
\end{writenote}

\part{Pointwise corrections for the proportional placement game}\label{ppg:log:part}
\partsource{Bundle Report~145, archive \file{A398540_Placement_Game_Logarithmic_Corrections_and_Inverse_Source.zip}, delivered as \file{Report145.tex} (995 lines, 21 pages), dated 3~October 2026; title block ``Report 145''; no author line. Printed as delivered, with \textbf{[write]} notes. Delivered Section~$k$ is Section~$k+13$ here, and its statement $k.j$ is $(k+13).j$; equations are renumbered within sections. ``Report 144'' is Part~\ref{ppg:lead:part}, whose numbers are unchanged. Remark~\ref{ppg:log:rem:inverse} is added in the write.}
\begin{partabstract}
For the winning probabilities $W_n$ of the continuous proportional placement game associated with OEIS A398540, we prove
\[
 \frac{W_n}{A\sqrt n\,r_*^n}
 =1+\frac1{4n}-\frac7{540}\frac{\log n}{n^2}
   +\frac{B_2}{n^2}+o(n^{-2}),\qquad
 A=r_*\sqrt{2\pi}.
\]
Starting from the established leading equivalent, we construct an exact coefficientwise inverse of the critical linearized Riccati equation. A tail-supremum contraction supplies the initial decay rate without a regularity assumption. Exact endpoint profiles and a centered covariance expansion then isolate the logarithmic forcing, with a uniform mixed remainder whose absolute diagonal sum is $O(n^{-1})$. A boundary derivative and signed summable-residue convolutions further give an absolutely convergent-series formula for the constant $B_2$. The stronger remainder proves eventual strict decrease of consecutive ratios. Summing a hypothetical ratio expansion rules out a pure inverse-power expansion of consecutive ratios through order $n^{-3}$; no uncontrolled little-o term is differentiated. We also refine the Lambert-W threshold inverse with uncertainty $o(x^{-2})$ before integer rounding. A full expansion, all-index ratio monotonicity, effective correction errors, and certified digits of $B_2$ or additional digits of $r_*$ remain open here.
\end{partabstract}
{\small\emph{Readings in this Part} (\wtag). $x_n=W_nR^n/(A\sqrt{n+1})$, $y_n=x_n-1$; $\beta_{ij}$ with $m=i+j$ the \emph{child} sum (Part~\ref{ppg:lead:part}'s $\alpha_{i,j}/a$); $z_n=\sum y_iy_{n-i}$, not the variable $z$; $p=\frac14$, $q=-\frac5{12}$, $c=-\frac7{360}$ (the covariance coefficient, not Part~\ref{ppg:all:part}'s $c$), $s=-\frac{11}{180}$, $d=-\frac14$, $\kappa=-\frac7{540}$; $A_i$ endpoint factors, not the amplitude; $a_i=(b_i-1)x_i$ (Part~\ref{ppg:all:part}'s $\delta_i$); $\gamma$ Euler's constant; in Section~\ref{ppg:log:sec:inverse}, $x_0$ is the Lambert root, not $1/A$, and $b=\frac14$.\par}
\section{The model and the conclusions}\label{ppg:log:sec:model}
There are $n$ ordered slots and $n$ independent uniform values on $(0,1)$ arriving one at a time. Each arriving value must immediately occupy an empty slot, with all placed values increasing from left to right. The \emph{proportional policy} places a value $t$ in a gap $(L,U)$ with $k$ empty slots into the slot numbered
\[
 \left\lceil k\frac{t-L}{U-L}\right\rceil
\]
within that gap. It loses if there is no empty slot in the containing gap. Ties have probability zero. Let $W_n$ be the probability of winning under this policy, with $W_0=1$. We make no assertion that the policy is optimal.

The exact recurrence is
\begin{equation}\label{ppg:log:eq:recurrence}
 W_n=\frac1n\sum_{k=1}^n W_{k-1}W_{n-k}Q_{n,k},\qquad
 Q_{n,k}=n\int_{(k-1)/n}^{k/n}\binom{n-1}{k-1}
 t^{k-1}(1-t)^{n-k}\,\dd t.
\end{equation}
To see this, condition on the first value $t$ in the $k$th interval. Exactly $k-1$ remaining values must be below it. Their binomial probability supplies the integrand, and, conditional on the count, the two affinely rescaled subgames are independent. Integrating and summing proves \eqref{ppg:log:eq:recurrence}. In particular $W_1=1$, $W_2=3/4$, and every $W_n$ is a positive rational number.
\begin{writenote}[5 October 2026]
This is the recurrence of Part~\ref{ppg:lead:part}, Proposition~\ref{ppg:lead:prop:recurrence}, where it is derived in full; the sentences above are Report~145's own summary of that derivation.
\end{writenote}

Report 144~\cite{r144} proves from this recurrence that
\begin{equation}\label{ppg:log:eq:foundation}
 r_*:=\lim_{n\to\infty}W_n^{1/n}\in(0,1),\qquad
 W_n\sim A\sqrt n\,r_*^n,\qquad
 \frac{W_{n+1}}{W_n}\longrightarrow r_*,\qquad A=r_*\sqrt{2\pi}.
\end{equation}
Its proof and first uniform kernel bounds are the foundation used below. For clarity, the unchanged PDF and LaTeX source of that report accompany this article. The new pointwise corrections are proved here in full; they are not inferred solely from a real-axis generating-function expansion.
\begin{writenote}[5 October 2026]
Report~144 is Part~\ref{ppg:lead:part} of this report, with the same section numbers: \eqref{ppg:log:eq:foundation} is Theorems~\ref{ppg:lead:thm:rate} and~\ref{ppg:lead:thm:pointwise}, and \eqref{ppg:log:eq:firstkernel}--\eqref{ppg:log:eq:firstrho} are \eqref{ppg:lead:eq:kexpansion}--\eqref{ppg:lead:eq:kremainder} divided by $a$. The ``accompanying'' copy of Report~144 (the archive's \file{foundation/} directory) is byte-identical to Report~144's own delivery and is not shipped.
\end{writenote}

\begin{theorem}[Two pointwise corrections]\label{ppg:log:thm:main}
For the probabilities defined by \eqref{ppg:log:eq:recurrence},
\begin{equation}\label{ppg:log:eq:main}
 \frac{W_n}{A\sqrt n\,r_*^n}
 =1+\frac1{4n}-\frac7{540}\frac{\log n}{n^2}
   +o\!\left(\frac{\log n}{n^2}\right).
\end{equation}
In particular the first correction is $1/(4n)+o(n^{-1})$. There is no finite constant $c_2$ for which the right-hand side can instead be written as
$1+1/(4n)+c_2/n^2+o(n^{-2})$.
\end{theorem}

Theorem~\ref{ppg:log:thm:ratio} below rules out a pure inverse-power expansion of the consecutive ratio through order $n^{-3}$ with remainder $o(n^{-3})$. Theorem~\ref{ppg:log:thm:inverse} gives the refined ceiling-safe inverse. Theorem~\ref{ppg:log:thm:full} strengthens \eqref{ppg:log:eq:main} by identifying the constant second term through absolutely convergent series. Corollary~\ref{ppg:log:cor:monotonicity} then proves eventual strict decrease of the ratios. An all-orders expansion, all-index monotonicity, and effective onsets remain outside the conclusions.
\begin{writenote}[5 October 2026]
Theorems~\ref{ppg:log:thm:main} and~\ref{ppg:log:thm:full} follow from the case $M=2$ of Part~\ref{ppg:all:part}'s Theorem~\ref{ppg:all:thm:all}, which takes this Part's results as inputs (\eqref{ppg:all:eq:base}, \eqref{ppg:all:eq:oldsecond}) and improves the remainder of \eqref{ppg:log:eq:full} to $O((1+\log n)^2/n^3)$ (\eqref{ppg:all:eq:strongsecond}). The all-orders expansion that this paragraph leaves outside the conclusions is Theorem~\ref{ppg:all:thm:all}; all-index monotonicity and effective onsets remain open (Section~\ref{ppg:sec:further}).
\end{writenote}

\subsection{Source scope and notation}\label{ppg:log:sub:scope}
Serra's manuscript~\cite{serra}, version 1, gives the model and recurrence in Theorem 3.12, equations (47)--(48), printed pages 23--24. OEIS A398540~\cite{oeis} records digits associated with the ratio limit, rather than the rational probability sequence itself. We use $r_*$ as the proved limit in \eqref{ppg:log:eq:foundation}; no decimal approximation is needed anywhere in this article. The statements below concern this exact recurrence relative to the inspected sources, without a claim of universal priority.
\begin{writenote}[5 October 2026]
Checked by the write against Serra's version-1 PDF: accurate. OEIS data are licensed CC~BY-SA~4.0.
\end{writenote}

All logarithms are natural. All asymptotic statements about sequences are coefficientwise as their integer index tends to infinity. Uniformity in two kernel indices will be stated explicitly. Write
\begin{equation}\label{ppg:log:eq:normalization}
 R=r_*^{-1},\qquad a=(2\pi)^{-1/2},\qquad A=\frac1{aR},\qquad
 x_n=\frac{W_nR^n}{A\sqrt{n+1}},\qquad y_n=x_n-1.
\end{equation}
The leading result \eqref{ppg:log:eq:foundation} says $x_n\to1$. Thus $x$ and $y$ are bounded, even though their finitely many initial values need not be close to 1 or 0.

\section{The normalized kernel and critical equation}\label{ppg:log:sec:critical}
For child indices $i,j\ge0$, put $m=i+j$ and define
\begin{equation}\label{ppg:log:eq:beta}
 \beta_{ij}=\frac{Q_{m+1,i+1}}a
             \sqrt{\frac{(i+1)(j+1)}{m+2}}.
\end{equation}
The distinction between the child sum $m$ and parent size $m+1$ matters. Symmetry of the integral gives $\beta_{ij}=\beta_{ji}$. Since $aAR=1$, the recurrence becomes
\begin{equation}\label{ppg:log:eq:xrecurrence}
 (n+1)x_{n+1}=\sum_{i+j=n}\beta_{ij}x_ix_j\qquad(n\ge0).
\end{equation}
The first uniform kernel expansion from Report 144, Section 4, is
\begin{align}
 \beta_{ij}&=1+\frac14\left(\frac1{i+1}+\frac1{j+1}\right)
                    -\frac5{12(i+j+1)}+\rho_{ij},\label{ppg:log:eq:firstkernel}\\
 |\rho_{ij}|&\le10\left((i+1)^{-2}+(j+1)^{-2}\right).
 \label{ppg:log:eq:firstrho}
\end{align}
In particular $|\beta_{ij}-1|\le C((i+1)^{-1}+(j+1)^{-1})$. That expansion is uniform through arbitrarily unbalanced splits, including the endpoints. Its proof uses the exact beta-bin integral and Robbins' factorial inequalities~\cite{robbins}. Section~\ref{ppg:log:sec:kernel} supplies the stronger boundary-matched estimate needed at the next order.

Define, for $|z|<1$,
\begin{align*}
 X(z)&=\sum_{n\ge0}x_nz^n,&Y(z)&=X(z)-\frac1{1-z}=\sum_{n\ge0}y_nz^n,\\
 e_n&=\sum_{i+j=n}(\beta_{ij}-1)x_ix_j,&z_n&=\sum_{i+j=n}y_iy_j,\qquad h_n=e_n+z_n.
\end{align*}
The use of $z_n$ for a convolution coefficient is distinct from the generating-function variable $z$. Boundedness and the kernel error give
\begin{equation}\label{ppg:log:eq:coarseforcing}
 e_n=O(1+\log(n+1)),\qquad z_n=O(n+1),\qquad h_n=O(n+1).
\end{equation}
Absolute convergence inside the disk allows products and differentiation, so
\begin{equation}\label{ppg:log:eq:Riccati}
 X'=X^2+E,\qquad
 Y'-\frac{2Y}{1-z}=H,\qquad E(z)=\sum e_nz^n,\quad H(z)=\sum h_nz^n.
\end{equation}

\begin{lemma}[Exact critical inverse]\label{ppg:log:lem:inverseT}
For every $n\ge1$,
\begin{equation}\label{ppg:log:eq:T}
 y_n=(Th)_n:=\frac{h_{n-1}}{n+2}
       -2(n+1)\sum_{k\ge n}\frac{h_k}{(k+1)(k+2)(k+3)}.
\end{equation}
The tail is absolutely convergent. On the positive indices $n\ge1$, $T$ annihilates every constant sequence.
\end{lemma}
\begin{proof}
The integrating factor in \eqref{ppg:log:eq:Riccati} is $(1-z)^2$. Boundedness of $y$ alone implies $|Y(r)|\le C/(1-r)$, hence $(1-r)^2Y(r)\to0$ as $r\uparrow1$. Moreover,
\begin{equation}\label{ppg:log:eq:absoluteintegral}
 \sum_{k\ge0}|h_k|\int_0^1(1-t)^2t^k\,\dd t
 =\sum_{k\ge0}\frac{2|h_k|}{(k+1)(k+2)(k+3)}<\infty.
\end{equation}
Integrating to the boundary, with this absolute justification, gives
\begin{equation}\label{ppg:log:eq:boundary}
 (1-r)^2Y(r)=-\int_r^1(1-t)^2H(t)\,\dd t\qquad(0\le r<1).
\end{equation}
In particular the homogeneous solution $C/(1-z)^2$, which would contribute $C(n+1)$ to $y_n$, has been eliminated by the boundary condition, rather than dropped formally. The right-hand side is the difference of an absolutely convergent constant and an analytic primitive inside the disk. The identity on the real interval therefore justifies analytic coefficient extraction within the disk without continuation across its boundary.

Put $D_k=(k+1)(k+2)(k+3)$. For the basis forcing $H(z)=z^k$, the numerator in \eqref{ppg:log:eq:boundary}, before division by $(1-z)^2$, is
\[
 -\frac2{D_k}+\frac{z^{k+1}}{k+1}-\frac{2z^{k+2}}{k+2}
                   +\frac{z^{k+3}}{k+3}.
\]
Since $[z^j](1-z)^{-2}=j+1$, its contribution to coefficient $n$ is zero for $k\le n-2$, is
$1/n-2/[n(n+2)]=1/(n+2)$ for $k=n-1$, and is $-2(n+1)/D_k$ for $k\ge n$. This proves \eqref{ppg:log:eq:T}. Finally,
\begin{equation}\label{ppg:log:eq:telescopes}
 \sum_{k\ge n}\frac1{D_k}=\frac1{2(n+1)(n+2)},\qquad
 \sum_{k\ge n}\frac{k+1}{D_k}=\frac1{n+2},
\end{equation}
by telescoping. The first equality proves constant cancellation.
\end{proof}
\begin{remark}
The cancellation is stated only for $n\ge1$. The separately fixed initial coefficient is $y_0=-\sum_{k\ge0}2h_k/D_k$; for constant forcing $h_k=c$ it is $-c/2$. Finite changes to a forcing through index $K$ are invisible in \eqref{ppg:log:eq:T} for $n\ge K+2$.
\end{remark}

\section{A decay rate from a tail contraction}\label{ppg:log:sec:contraction}
The first correction requires an actual coefficientwise rate beyond $y_n\to0$. We derive it before using any sharper forcing expansion.

\begin{lemma}[Floor-safe tail bound]\label{ppg:log:lem:tail}
The normalized error satisfies
\begin{equation}\label{ppg:log:eq:coarserate}
 y_n=O\!\left(\frac{\log(n+2)}{n+1}\right).
\end{equation}
\end{lemma}
\begin{proof}
Let $M_N=\sup_{k\ge N}|y_k|$. This is finite, nonincreasing, and tends to zero. Since $y_i\to0$, the absolute Cesaro averages vanish, and
\[
 \sum_{i\le\lfloor k/2\rfloor}|y_i|=o(k+1).
\]
Symmetry of the convolution gives
\[
 |z_k|\le2M_{\lfloor k/2\rfloor}
                 \sum_{i\le\lfloor k/2\rfloor}|y_i|.
\]
For even $k$ the middle term is counted twice, which only enlarges the upper bound. Thus, for any $\sigma>0$, there is a fixed threshold beyond which
\begin{equation}\label{ppg:log:eq:zcontraction}
 |z_k|\le\sigma(k+1)M_{\lfloor k/2\rfloor}.
\end{equation}
If $n$ is past that threshold plus one, put $m=\lfloor(n-1)/2\rfloor$. Monotonicity of $M$ and \eqref{ppg:log:eq:telescopes} yield
\begin{equation}\label{ppg:log:eq:Tzcontraction}
 |(Tz)_n|\le\sigma M_m\left(\frac n{n+2}+\frac{2(n+1)}{n+2}\right)
 \le3\sigma M_m.
\end{equation}
The coarse forcing bound \eqref{ppg:log:eq:coarseforcing} directly gives
\[
 |(Te)_n|\le C f_n,\qquad f_n=\frac{1+\log(n+1)}{n+1},
\]
since the local term has this order and the tail is bounded by a constant times $(n+1)\int_n^\infty(1+\log(t+1))/(t+1)^3\,\dd t$.
The sequence $f_N$ is nonincreasing. Taking the supremum over every $n\ge N$ in $y=Te+Tz$, including the floor in \eqref{ppg:log:eq:Tzcontraction}, gives
\begin{equation}\label{ppg:log:eq:Mrecurrence}
 M_N\le qM_{\lfloor(N-1)/2\rfloor}+Cf_N
\end{equation}
for all $N\ge N_0$, where $q=3\sigma$ may be chosen arbitrarily small. Fix $q=1/6$ and $N_0\ge2$.

For $N\ge2$ and $m=\lfloor(N-1)/2\rfloor$,
\[
 \frac{f_m}{f_N}\le\frac{N+1}{m+1}\le3.
\]
Indeed, for $N=2\ell$ the last ratio is $(2\ell+1)/\ell\le3$, and for $N=2\ell+1$ it equals 2. Choose $B\ge2C$ large enough that $M_j\le Bf_j$ for every $0\le j<N_0$. Strong induction using $m<N$ in \eqref{ppg:log:eq:Mrecurrence} then gives
\[
 M_N\le qBf_m+Cf_N\le(B/2+C)f_N\le Bf_N.
\]
All startup indices, including $m=0$ when $N=2$, are covered. This proves \eqref{ppg:log:eq:coarserate} without presupposing a decay rate for $M_N$.
\end{proof}

\begin{lemma}[Quadratic error at the first scale]\label{ppg:log:lem:firstquadratic}
The convolution satisfies
\begin{equation}\label{ppg:log:eq:quadraticcoarse}
 z_n=O\!\left(\frac{\log^3(n+2)}{n+1}\right),\qquad
 (Tz)_n=O\!\left(\frac{\log^3(n+2)}{(n+1)^2}\right)=o(n^{-1}).
\end{equation}
\end{lemma}
\begin{proof}
On $i\le n/2$, the other index is at least $n/2$. Apply \eqref{ppg:log:eq:coarserate} to that index, then sum the remaining factor:
\[
 |z_n|\le C\frac{\log(n+2)}{n+1}
             \sum_{i\le n/2}\frac{\log(i+2)}{i+1}
 =O\!\left(\frac{\log^3(n+2)}{n+1}\right).
\]
The local term in \eqref{ppg:log:eq:T} has the claimed extra inverse power. The tail does also, by integrating $\log^3(t+2)/(t+1)^4$ and multiplying by $n+1$.
\end{proof}

\section{The first pointwise correction}\label{ppg:log:sec:first}
For a fixed endpoint distance $i\ge0$, define
\begin{equation}\label{ppg:log:eq:bi}
 b_i=\lim_{j\to\infty}\beta_{ij}
     =\frac{\sqrt{2\pi(i+1)}}{i!}\int_i^{i+1}s^ie^{-s}\,\dd s.
\end{equation}
Here $s^0=1$, including at zero, so $b_0=\sqrt{2\pi}(1-e^{-1})$. To verify the limit directly, substitute $s=(i+j+1)t$ in the beta-bin integral. The factor $\binom{i+j}{i}/(i+j+1)^i$ tends to $1/i!$, the square-root factor tends to $\sqrt{i+1}$, and $(1-s/(i+j+1))^j\to e^{-s}$ uniformly on the fixed interval $[i,i+1]$.

\begin{lemma}[Forcing to constant order]\label{ppg:log:lem:firstforcing}
There is a finite real constant $B$ such that
\begin{equation}\label{ppg:log:eq:firstforcing}
 e_n=\tfrac12\log(n+1)+B+o(1).
\end{equation}
\end{lemma}
\begin{proof}
By symmetry, \eqref{ppg:log:eq:firstkernel} gives exactly
\begin{equation}\label{ppg:log:eq:edecompose}
 e_n=\frac12\sum_{i=0}^n\frac{x_ix_{n-i}}{i+1}
    -\frac5{12(n+1)}\sum_{i=0}^nx_ix_{n-i}
    +\sum_{i=0}^n\rho_{i,n-i}x_ix_{n-i}.
\end{equation}
The bound \eqref{ppg:log:eq:coarserate} makes $K=\sum_{i\ge0}y_i/(i+1)$ absolutely convergent. Hence
\[
 \sum_{i=0}^n\frac{x_i}{i+1}=\log(n+1)+\gamma+K+o(1).
\]
Replacing the omitted factor 1 by $x_{n-i}$ changes this by
\[
 \sum_{i=0}^n\frac{x_iy_{n-i}}{i+1}
 =O\!\left(\frac{\log^2(n+2)}{n+1}\right)=o(1).
\]
For $i\le n/2$, use the bound on $y_{n-i}$ and the harmonic sum; for $i>n/2$, use $(i+1)^{-1}=O(n^{-1})$ and $\sum_{j\le n/2}|y_j|=O(\log^2(n+2))$.

Boundedness and $x_n\to1$ also imply
\[
 \frac1{n+1}\sum_{i=0}^nx_ix_{n-i}\longrightarrow1.
\]
For example expand the numerator as $n+1+2\sum_{i\le n}y_i+z_n$ and use the absolute Cesaro limit and \eqref{ppg:log:eq:quadraticcoarse}.

Set $\zeta_i=b_i-1-1/[4(i+1)]$. Taking a fixed-$i$ limit in \eqref{ppg:log:eq:firstkernel} gives $\rho_{ij}\to\zeta_i$ and $|\zeta_i|\le10/(i+1)^2$. On the half diagonal $i\le n/2$, boundedness of $x$ and \eqref{ppg:log:eq:firstrho} give the summable domination
\[
 |\rho_{i,n-i}x_ix_{n-i}|\le C/(i+1)^2.
\]
Each fixed summand tends to $\zeta_ix_i$. Dominated convergence on the strict half-sum, extended by zero, and symmetry therefore show
\[
 \sum_{i=0}^n\rho_{i,n-i}x_ix_{n-i}\longrightarrow2\sum_{i\ge0}\zeta_ix_i.
\]
The possible even-index center is $O((n+1)^{-2})$ and tends to zero. This treats both endpoints and both parities. Inserting the three limits into \eqref{ppg:log:eq:edecompose} gives \eqref{ppg:log:eq:firstforcing}, with
\begin{equation}\label{ppg:log:eq:B}
 B=\tfrac12(\gamma+K)-\tfrac5{12}+2\sum_{i\ge0}\zeta_ix_i.
\end{equation}
Its value is unnecessary because $T$ annihilates constants at positive indices.
\end{proof}

\begin{lemma}[The logarithm under the critical inverse]\label{ppg:log:lem:Tlog}
For $\ell_k=\log(k+1)$,
\begin{equation}\label{ppg:log:eq:Tlog}
 (T\ell)_n=-\frac1{2n}+O(n^{-2}).
\end{equation}
For an arbitrary sequence $v_k=o(1)$, $(Tv)_n=o(n^{-1})$.
\end{lemma}
\begin{proof}
Use $D_k^{-1}=\tfrac12([(k+1)(k+2)]^{-1}-[(k+2)(k+3)]^{-1})$ and summation by parts. The boundary term at infinity vanishes, giving
\begin{equation}\label{ppg:log:eq:Tlogexact}
 (T\ell)_n=\frac{\log n-\log(n+1)}{n+2}
 -(n+1)\sum_{k\ge n+1}\frac{\log(1+1/k)}{(k+1)(k+2)}.
\end{equation}
Now $\log(1+1/k)=1/k+O(k^{-2})$ and
\[
 \sum_{k\ge n+1}\frac1{k(k+1)(k+2)}
       =\frac1{2(n+1)(n+2)}.
\]
The isolated logarithm difference and the summed error in \eqref{ppg:log:eq:Tlogexact} are $O(n^{-2})$. Thus $(T\ell)_n=-1/[2(n+2)]+O(n^{-2})$, which is \eqref{ppg:log:eq:Tlog}. Finally, the positive absolute weights of \eqref{ppg:log:eq:T} give
\[
 |(Tv)_n|\le\frac2{n+2}\sup_{k\ge n-1}|v_k|=o(n^{-1}).
\]
No variation or smoothness assumption on $v$ is required.
\end{proof}

\begin{proposition}[First correction]\label{ppg:log:prop:first}
With the normalization \eqref{ppg:log:eq:normalization},
\begin{equation}\label{ppg:log:eq:firstpointwise}
 x_n=1-\frac1{4n}+o(n^{-1}),\qquad
 \frac{W_n}{A\sqrt n\,r_*^n}=1+\frac1{4n}+o(n^{-1}).
\end{equation}
In particular, for a fixed finite constant $C$, $|y_n|\le C/(n+1)$ for all $n\ge0$.
\end{proposition}
\begin{proof}
Apply $T$ to \eqref{ppg:log:eq:firstforcing} and use Lemma~\ref{ppg:log:lem:Tlog}, constant cancellation, and Lemma~\ref{ppg:log:lem:firstquadratic} in $y=Te+Tz$. This gives $y_n=-1/(4n)+o(n^{-1})$. The finitely many small indices are absorbed by increasing $C$. Finally,
\[
 \frac{W_n}{A\sqrt n\,r_*^n}=\sqrt{1+1/n}\,x_n,
 \qquad \sqrt{1+1/n}=1+\frac1{2n}+O(n^{-2}),
\]
which explains the change of sign in the first correction.
\end{proof}

\section{The boundary covariance and a stronger kernel estimate}\label{ppg:log:sec:kernel}
The summable endpoint error in \eqref{ppg:log:eq:firstrho} suffices to determine the constant-order forcing. It does not suffice to determine the next $\log n/n$ forcing coefficient: a bound by $i^{-2}+j^{-2}$ alone may sum to $O(1)$. We retain the exact endpoint profiles and derive a genuinely mixed remainder.

\begin{lemma}[Exact boundary factorization]\label{ppg:log:lem:factor}
For integers $i,j\ge1$, put $m=i+j$. Define the exact Stirling remainder $r_k$ by
\begin{equation}\label{ppg:log:eq:Stirling}
 k!=\sqrt{2\pi k}(k/e)^ke^{r_k},\qquad
 \frac1{12k+1}<r_k<\frac1{12k}\quad(k\ge1),
\end{equation}
and set
\begin{align}
 A_i&=\sqrt{1+1/i}\,e^{-r_i},&
 f_i(v)&=(1+v/i)^ie^{-v},& I_i&=\int_0^1f_i(v)\,\dd v,\label{ppg:log:eq:fi}\\
 H_m&=(1+2/m)^{-1/2}\exp\{r_m+1-m\log(1+1/m)\},\label{ppg:log:eq:Hm}\\
 C_{ij}&=\frac{\int_0^1f_i(v)f_j(1-v)\,\dd v}{I_iI_j}.
 \label{ppg:log:eq:Cij}
\end{align}
Then $b_i=A_iI_i$ and
\begin{equation}\label{ppg:log:eq:factor}
 \beta_{ij}=b_ib_jH_mC_{ij}.
\end{equation}
Here $A_i$ is an indexed endpoint factor, not the amplitude $A$ in \eqref{ppg:log:eq:normalization}.
\end{lemma}
\begin{writenote}[5 October 2026]
Part~\ref{ppg:all:part}'s Lemma~\ref{ppg:all:lem:factor} writes this factorization as $\beta_{ij}=H_n\int_0^1g_i(v)g_j(1-v)\,\dd v$, with $g_i=A_if_i$ in the present notation, and extends it to a zero child with $g_0(v)=\sqrt{2\pi}e^{-v}$; its $b_i=\int_0^1g_i$ is the $b_i$ of \eqref{ppg:log:eq:bi}.
\end{writenote}
\begin{proof}
Substitute $t=(i+v)/(m+1)$ in \eqref{ppg:log:eq:beta} and \eqref{ppg:log:eq:recurrence}. The factor $m+1$ in $Q$ cancels the Jacobian. Applying \eqref{ppg:log:eq:Stirling} to the binomial coefficient gives the exact identity
\begin{align}\label{ppg:log:eq:modalexact}
 \beta_{ij}
 ={}&\sqrt{\frac{m(i+1)(j+1)}{ij(m+2)}}\,
 e^{r_m-r_i-r_j}\left(\frac m{m+1}\right)^m\nonumber\\
 &\hspace{15mm}\cdot\int_0^1(1+v/i)^i(1+(1-v)/j)^j\,\dd v.
\end{align}
Removing the factors $A_iA_j$ leaves $(1+2/m)^{-1/2}e^{r_m}$. Replacing the two integrand factors by $f_i(v)f_j(1-v)$ introduces $e$, leaving exactly $H_m$. Finally, substitute $s=i+v$ in \eqref{ppg:log:eq:bi} and use \eqref{ppg:log:eq:Stirling}; the result is $b_i=A_iI_i$. This proves \eqref{ppg:log:eq:factor}.
\end{proof}

\begin{lemma}[Uniform mixed covariance]\label{ppg:log:lem:covariance}
Uniformly for all integers $i,j\ge1$, including arbitrarily unbalanced pairs,
\begin{equation}\label{ppg:log:eq:covariance}
 C_{ij}=1-\frac7{360ij}
       +O\!\left(i^{-2}j^{-1}+i^{-1}j^{-2}\right).
\end{equation}
\end{lemma}
\begin{proof}
For $0\le v\le1$, elementary logarithm bounds give
\[
 \log f_i(v)=-\frac{v^2}{2i}+\delta_i(v),\qquad
 0\le\delta_i(v)\le\frac{v^3}{3i^2},\qquad
 -\frac{v^2}{2i}\le\log f_i(v)\le0.
\]
For $u\le0$, $0\le e^u-1-u\le u^2/2$. Consequently
\begin{equation}\label{ppg:log:eq:uniformfi}
 f_i(v)=1-\frac{v^2}{2i}+\varepsilon_i(v),\qquad
 0\le\varepsilon_i(v)\le\frac{v^3/3+v^4/8}{i^2}
                         \le\frac{11}{24i^2}.
\end{equation}
Also $0\le1-f_i(v)\le v^2/(2i)$, so
\begin{equation}\label{ppg:log:eq:Ibounds}
 1-\frac1{6i}\le I_i\le1,\qquad I_i\ge\frac56.
\end{equation}
Let $\overline\varepsilon_i=\int_0^1\varepsilon_i(v)\,\dd v$ and $E_i(v)=\varepsilon_i(v)-\overline\varepsilon_i$. Because $\varepsilon_i$ is confined to an interval of length $11/(24i^2)$,
\begin{equation}\label{ppg:log:eq:centered}
 f_i(v)-I_i=-\frac{v^2-1/3}{2i}+E_i(v),\qquad
 \|E_i\|_\infty\le\frac{11}{24i^2}.
\end{equation}
Center \emph{before} multiplying. The numerator of $C_{ij}-1$ is
\begin{align*}
 &\int_0^1f_i(v)f_j(1-v)\,\dd v-I_iI_j\\
 &\hspace{12mm}=\int_0^1[f_i(v)-I_i][f_j(1-v)-I_j]\,\dd v.
\end{align*}
The leading term from \eqref{ppg:log:eq:centered} is
\begin{align}\label{ppg:log:eq:covintegral}
 \frac1{4ij}\int_0^1(v^2-1/3)((1-v)^2-1/3)\,\dd v
 &=\frac1{4ij}\left(\frac1{30}-\frac19\right)\nonumber\\
 &=-\frac7{360ij}.
\end{align}
The cross terms with $E_i$ and $E_j$ are respectively $O(i^{-2}j^{-1})$ and $O(i^{-1}j^{-2})$. Their product is $O(i^{-2}j^{-2})$, absorbed by the same mixed bound. This argument uses $|v^2-1/3|\le2/3$ and uniform suprema on $[0,1]$.

Finally \eqref{ppg:log:eq:Ibounds} gives
\[
 \frac1{I_iI_j}\le\frac{36}{25},\qquad
 \left|\frac1{I_iI_j}-1\right|
 \le\frac6{25}\left(\frac1i+\frac1j\right).
\]
Dividing the centered numerator by $I_iI_j$ therefore preserves the mixed error and gives \eqref{ppg:log:eq:covariance}. The cancellation of constant parts is the reason no unmixed $i^{-2}$ or $j^{-2}$ remainder remains at this step.
\end{proof}

\begin{proposition}[Boundary-matched kernel with an absolute summed error]\label{ppg:log:prop:matched}
Put
\begin{equation}\label{ppg:log:eq:pqs}
 p=\frac14,\qquad q=-\frac5{12},\qquad c=-\frac7{360},\qquad
 s=p^2+pq+c=-\frac{11}{180}.
\end{equation}
For $i,j\ge1$ and $m=i+j$,
\begin{equation}\label{ppg:log:eq:matched}
 \beta_{ij}=b_i+b_j-1+\frac q m+\frac s{ij}
       +O\!\left(m^{-2}+i^{-2}j^{-1}+i^{-1}j^{-2}\right),
\end{equation}
uniformly. Including the two endpoints by omitting $s/(ij)$ there, define $R_{ij}$ by
\begin{equation}\label{ppg:log:eq:matchedall}
 \beta_{ij}-1=(b_i-1)+(b_j-1)+\frac q m
       +\frac s{ij}\mathbf1_{\{i,j\ge1\}}+R_{ij}\qquad(m\ge2),
\end{equation}
where the indicated term is interpreted as zero at an endpoint. Then
\begin{equation}\label{ppg:log:eq:summedR}
 \sum_{i=0}^m|R_{i,m-i}|=O(m^{-1}).
\end{equation}
The same order holds after multiplying by any uniformly bounded weights $x_ix_{m-i}$.
\end{proposition}
\begin{proof}
Robbins' bounds imply $r_k=1/(12k)+O(k^{-2})$, uniformly for $k\ge1$. Equations \eqref{ppg:log:eq:fi}, \eqref{ppg:log:eq:uniformfi}, and \eqref{ppg:log:eq:Ibounds} give
\begin{equation}\label{ppg:log:eq:bHexp}
 A_i=1+\frac5{12i}+O(i^{-2}),\quad
 I_i=1-\frac1{6i}+O(i^{-2}),\quad
 b_i=1+\frac p i+O(i^{-2}).
\end{equation}
In particular $b_i$ is bounded. Similarly, for $m\ge2$,
\begin{align*}
 \log H_m
 &=-\tfrac12\log(1+2/m)+r_m+1-m\log(1+1/m)\\
 &=-\frac1m+\frac1{12m}+\frac1{2m}+O(m^{-2})
   =\frac q m+O(m^{-2}),
\end{align*}
so $H_m=1+q/m+O(m^{-2})$. These estimates are uniform down to the smallest stated indices: the arguments of each logarithm remain in a fixed bounded interval and the Taylor remainders are bounded there.

Write $u_i=b_i-1=p/i+O(i^{-2})$ and $M_{ij}=i^{-2}j^{-1}+i^{-1}j^{-2}$. Keeping the endpoint factors exact,
\[
 b_ib_j=b_i+b_j-1+\frac{p^2}{ij}+O(M_{ij}).
\]
The common first-order factor contributes
\[
 \frac q m b_ib_j=\frac q m+\frac{pq}{ij}+O(M_{ij}),
\]
because $(1/i+1/j)/m=1/(ij)$ exactly, while
\[
 \frac1m(i^{-2}+j^{-2})\le M_{ij},\qquad
 \frac1{mij}\le\frac14M_{ij}.
\]
The common $O(m^{-2})$ term stays of that order after multiplication by bounded factors. Since $b_ib_jH_m=1+O(1/i+1/j)$, the covariance in \eqref{ppg:log:eq:covariance} contributes $c/(ij)+O(M_{ij})$. Combining these terms proves \eqref{ppg:log:eq:matched} and the arithmetic in \eqref{ppg:log:eq:pqs}.

The exact endpoint integral is, for $m\ge1$,
\begin{equation}\label{ppg:log:eq:exactendpoint}
 \beta_{0m}=\beta_{m0}
 =\sqrt{2\pi}\left[1-\left(\frac m{m+1}\right)^{m+1}\right]
               \sqrt{\frac{m+1}{m+2}}=b_0+O(m^{-1}).
\end{equation}
For an explicit error argument, put $s_m=(m/(m+1))^{m+1}$. The elementary bounds $z/(1+z)\le\log(1+z)\le z$ give $-1-1/m\le\log s_m\le-1$, hence $|s_m-e^{-1}|\le e^{-1}/m$. The square-root factor differs from 1 by at most $1/(m+2)$. Since $b_m=1+O(m^{-1})$, each endpoint contributes $O(m^{-1})$ to the remainder in \eqref{ppg:log:eq:matchedall}.

In the interior, the common error sums to $O(m\cdot m^{-2})$. For the mixed errors, splitting at $m/2$ gives
\begin{equation}\label{ppg:log:eq:mixedsum}
 \sum_{i=1}^{m-1}M_{i,m-i}
 =m\sum_{i=1}^{m-1}\frac1{i^2(m-i)^2}
 \le\frac8m\sum_{i\ge1}\frac1{i^2}=O(m^{-1}).
\end{equation}
This proves \eqref{ppg:log:eq:summedR}. If $|x_i|\le L$, the absolute weighted remainder is at most $L^2$ times this sum.
\end{proof}

\section{The signed logarithmic forcing}\label{ppg:log:sec:forcing}
Write $d=-1/4$. Proposition~\ref{ppg:log:prop:first} gives
\begin{equation}\label{ppg:log:eq:yfirstinput}
 y_n=\frac d n+o(n^{-1}),\qquad |y_n|\le\frac C{n+1}.
\end{equation}
All formulas with $1/n$ in this section are used only for $n\ge1$; finitely many initial terms are handled separately.

\begin{lemma}[Signed harmonic convolution]\label{ppg:log:lem:harmonicconv}
Suppose $u_i=P/i+O(i^{-2})$ for $i\ge1$, with $u_0$ finite, and $v_j=D/j+o(j^{-1})$ for $j\ge1$, with $|v_j|\le C/(j+1)$ including $j=0$. Then
\begin{equation}\label{ppg:log:eq:convlemma}
 \sum_{i=0}^nu_iv_{n-i}
 =2PD\frac{\log n}{n}+o\!\left(\frac{\log n}{n}\right).
\end{equation}
The conclusion also holds when both factors have the second kind of expansion, with their leading coefficients in place of $P,D$.
\end{lemma}
\begin{proof}
The exact identity
\begin{equation}\label{ppg:log:eq:harmonicidentity}
 \sum_{i=1}^{n-1}\frac1{i(n-i)}=\frac{2\HH_{n-1}}n
          =2\frac{\log n}{n}+O(n^{-1}),\qquad
 \HH_N=\sum_{j=1}^N\frac1j,
\end{equation}
is obtained by partial fractions. The two endpoint terms are $O(n^{-1})$. Convolving an $O(i^{-2})$ remainder with $O((j+1)^{-1})$ gives $O(n^{-1})$: on one half use the summability of $i^{-2}$, and on the other use $i^{-2}=O(n^{-2})$ and a harmonic bound. For the little-o remainder in $v_j$, choose $J$ so that its absolute value is at most $\epsilon/j$ for $j\ge J$. These terms are bounded by a constant times $\epsilon\log n/n$ by \eqref{ppg:log:eq:harmonicidentity}; the finitely many $j<J$ contribute $O_J(n^{-1})$. Let $\epsilon\downarrow0$. If both factors have little-o remainders, apply this argument to each after writing their difference from the product of leading terms as the sum of two such convolutions. The uniform inverse-index bounds control the other factor. Signs of $P$ and $D$ are retained throughout.
\end{proof}

Set $a_i=(b_i-1)x_i$. From \eqref{ppg:log:eq:bHexp} and \eqref{ppg:log:eq:yfirstinput},
\begin{equation}\label{ppg:log:eq:ai}
 a_i=\frac p i+O(i^{-2})\quad(i\ge1),\qquad
 \sum_{i=0}^na_i=p\log n+C_a+O(n^{-1})
\end{equation}
for a finite constant $C_a$. Indeed the difference $a_i-p/i$ is absolutely summable with tail $O(n^{-1})$.
By Proposition~\ref{ppg:log:prop:matched}, the exact forcing has the representation
\begin{equation}\label{ppg:log:eq:eforcingmatched}
 e_n=2\sum_{i=0}^na_ix_{n-i}
       +\frac q n\sum_{i=0}^nx_ix_{n-i}
       +s\sum_{i=1}^{n-1}\frac{x_ix_{n-i}}{i(n-i)}+O(n^{-1}).
\end{equation}
This includes the endpoints with their verified error, and the total remainder is bounded absolutely.

\begin{proposition}[Logarithmic forcing coefficient]\label{ppg:log:prop:forcing}
For a finite constant $B$,
\begin{align}
 e_n&=\tfrac12\log n+B-\frac{59}{360}\frac{\log n}{n}
                  +o\!\left(\frac{\log n}{n}\right),\label{ppg:log:eq:esharp}\\
 z_n&=\frac18\frac{\log n}{n}
                  +o\!\left(\frac{\log n}{n}\right),\label{ppg:log:eq:zsharp}\\
 h_n&=\tfrac12\log n+B-\frac7{180}\frac{\log n}{n}
                  +o\!\left(\frac{\log n}{n}\right).\label{ppg:log:eq:hsharp}
\end{align}
The constant is the same $B$ as in \eqref{ppg:log:eq:firstforcing}.
\end{proposition}
\begin{proof}
Split the first term of \eqref{ppg:log:eq:eforcingmatched} as $2\sum a_i+2\sum a_iy_{n-i}$. Equations \eqref{ppg:log:eq:ai}, \eqref{ppg:log:eq:yfirstinput}, and Lemma~\ref{ppg:log:lem:harmonicconv} give
\begin{equation}\label{ppg:log:eq:forcing1}
 2\sum_{i=0}^na_ix_{n-i}
 =2p\log n+2C_a+4pd\frac{\log n}{n}
                         +o\!\left(\frac{\log n}{n}\right).
\end{equation}
Next $\sum_{i=0}^ny_i=d\log n+o(\log n)$; this follows by the same epsilon and finite-prefix argument applied to the harmonic sum. Also $z_n=O(\log n/n)$ by \eqref{ppg:log:eq:yfirstinput} and the half-split bound. Therefore
\begin{equation}\label{ppg:log:eq:forcing2}
 \sum_{i=0}^nx_ix_{n-i}=n+1+2\sum_{i=0}^ny_i+z_n
      =n+1+2d\log n+o(\log n).
\end{equation}
Finally $x_ix_{n-i}-1=O(1/i+1/(n-i))$ in the interior, so \eqref{ppg:log:eq:harmonicidentity} and \eqref{ppg:log:eq:mixedsum} give
\begin{equation}\label{ppg:log:eq:forcing3}
 \sum_{i=1}^{n-1}\frac{x_ix_{n-i}}{i(n-i)}
       =2\frac{\log n}{n}+O(n^{-1}).
\end{equation}
Substituting \eqref{ppg:log:eq:forcing1}--\eqref{ppg:log:eq:forcing3} into \eqref{ppg:log:eq:eforcingmatched} gives leading term $2p\log n=\tfrac12\log n$, constant $2C_a+q$, and coefficient
\begin{equation}\label{ppg:log:eq:efactor}
 4pd+2qd+2s=-\frac14+\frac5{24}-\frac{11}{90}
             =-\frac{59}{360}.
\end{equation}
Because \eqref{ppg:log:eq:esharp} also implies \eqref{ppg:log:eq:firstforcing}, the constants agree: $2C_a+q=B$.

Apply the final assertion of Lemma~\ref{ppg:log:lem:harmonicconv} to the two factors $y_i,y_{n-i}$. Their signed leading coefficients are both $d$, giving
\[
 z_n=2d^2\frac{\log n}{n}+o(\log n/n)
     =\frac18\frac{\log n}{n}+o(\log n/n).
\]
Adding it to $e_n$ produces $-59/360+1/8=-7/180$ in \eqref{ppg:log:eq:hsharp}.
\end{proof}

\section{The logarithmic pointwise correction}\label{ppg:log:sec:second}
\begin{lemma}[Critical inverse of a logarithm divided by its index]\label{ppg:log:lem:Tg}
For $g_k=\log(k+1)/(k+1)$,
\begin{equation}\label{ppg:log:eq:Tg}
 (Tg)_n=\frac13\frac{\log n}{n^2}+O(n^{-2}).
\end{equation}
For arbitrary sequences at the indicated scales,
\begin{align}
 v_k=o\!\left(\frac{\log(k+2)}{k+1}\right)
 &\quad\Longrightarrow\quad (Tv)_n=o\!\left(\frac{\log n}{n^2}\right),\label{ppg:log:eq:To}\\
 w_k=O((k+1)^{-1})
 &\quad\Longrightarrow\quad (Tw)_n=O(n^{-2}).\label{ppg:log:eq:TO}
\end{align}
\end{lemma}
\begin{proof}
The local term in \eqref{ppg:log:eq:T} for $g$ is $\log n/[n(n+2)]=\log n/n^2+O(\log n/n^3)$. In the tail,
\begin{align*}
 \sum_{k\ge n}\frac{g_k}{D_k}
 &=\sum_{k\ge n}\frac{\log(k+1)}{(k+1)^4}
                    +O(\log n/n^4)\\
 &=\int_n^\infty\frac{\log t}{t^4}\,\dd t+O(\log n/n^4)\\
 &=\frac{\log n}{3n^3}+\frac1{9n^3}+O(\log n/n^4).
\end{align*}
For large $n$, the summand is decreasing, so integral comparison bounds the sum-integral error by its first term. Multiplication by $2(n+1)$ in \eqref{ppg:log:eq:T} gives the coefficient $1-2/3=1/3$ in \eqref{ppg:log:eq:Tg}. The remaining terms are $O(n^{-2})$.

For \eqref{ppg:log:eq:To}, bound $|v_k|$ by $\epsilon\log(k+2)/(k+1)$ uniformly for every $k\ge N_\epsilon$. Once $n-1\ge N_\epsilon$, the local term and the entire tail in the absolute version of \eqref{ppg:log:eq:T} obey the same comparison. Their sum is at most $C\epsilon\log n/n^2$. Let $\epsilon\downarrow0$. Similarly, for \eqref{ppg:log:eq:TO}, the local term is $O(n^{-2})$ and the tail is bounded by $C(n+1)\sum_{k\ge n}(k+1)^{-4}=O(n^{-2})$. Neither conclusion requires smoothness or monotonicity of the input sequence.
\end{proof}

\begin{proof}[Proof of Theorem~\ref{ppg:log:thm:main}]
Replace $\log n$ and $\log n/n$ in \eqref{ppg:log:eq:hsharp} by the everywhere-defined sequences $\ell_n=\log(n+1)$ and $g_n=\log(n+1)/(n+1)$. The first change is $O(n^{-1})$, the second is $O(\log n/n^2)$; both are $o(\log n/n)$. Thus
\[
 h_n=\tfrac12\ell_n+B-\frac7{180}g_n+v_n,
 \qquad v_n=o(\log(n+2)/(n+1)).
\]
Initial-index changes do not affect the following asymptotics. Apply \eqref{ppg:log:eq:Tlog}, \eqref{ppg:log:eq:Tg}, \eqref{ppg:log:eq:To}, and constant cancellation to obtain
\begin{equation}\label{ppg:log:eq:xsharp}
 y_n=-\frac1{4n}-\frac7{540}\frac{\log n}{n^2}
                +o\!\left(\frac{\log n}{n^2}\right).
\end{equation}
The $O(n^{-2})$ error in \eqref{ppg:log:eq:Tlog} is crucially smaller than the displayed logarithmic scale. Equation \eqref{ppg:log:eq:TO} also shows directly that the absolutely summed $O(n^{-1})$ kernel forcing cannot hide another logarithmic coefficient.

Multiply $1+y_n$ by $\sqrt{1+1/n}=1+1/(2n)+O(n^{-2})$. The first correction becomes $+1/(4n)$; the induced constant-order $n^{-2}$ terms are $o(\log n/n^2)$ and do not alter $-7/540$. This proves \eqref{ppg:log:eq:main}. If a pure constant $c_2/n^2$ with remainder $o(n^{-2})$ replaced the displayed logarithmic term, subtracting $1+1/(4n)$ and dividing by $\log n/n^2$ would give limit 0 rather than $-7/540$, a contradiction.
\end{proof}

\section{The full second-order constant}\label{ppg:log:sec:constant}
The logarithmic result also supplies enough summability to resolve the next constant. In this section $d=-1/4$, $p=1/4$, $q=-5/12$, and $s=-11/180$ retain their meanings. We use $\kappa=-7/540$ for the sequence-level logarithmic coefficient and $L_h=-7/180$ for the forcing coefficient. In particular, \eqref{ppg:log:eq:xsharp} gives the already-proved bound
\begin{equation}\label{ppg:log:eq:residuebound}
 y_i-d/i=O\!\left(\frac{\log(i+2)}{(i+1)^2}\right)\qquad(i\ge1).
\end{equation}
This input, rather than any assumed constant second term, underlies the following residue sums.

\begin{lemma}[Second one-variable coefficients]\label{ppg:log:lem:secondfactors}
Uniformly over the positive indices in their domains,
\begin{align}
 H_m&=1+\frac q m+\frac{F_2}{m^2}+O(m^{-3}),& F_2&=\frac{217}{288},\label{ppg:log:eq:Hsecond}\\
 b_i&=1+\frac p i+\frac{b^{(2)}}{i^2}+O(i^{-3}),& b^{(2)}&=-\frac{179}{1440}.
 \label{ppg:log:eq:bsecond}
\end{align}
\end{lemma}
\begin{proof}
The exact Stirling remainder satisfies
\[
 r_m-r_{m+1}=(m+\tfrac12)\log(1+1/m)-1
 =\frac1{12m^2}-\frac1{12m^3}+O(m^{-4}).
\]
The same two terms occur in $1/(12m)-1/[12(m+1)]$. Subtract these identities and sum the $O(m^{-4})$ tail; Robbins' bounds give $r_m\to0$. Thus
\begin{equation}\label{ppg:log:eq:Stirlingrefined}
 r_m=\frac1{12m}+O(m^{-3}).
\end{equation}
This self-contained refinement is enough; no unproved higher Stirling term is needed. Taylor expansion in \eqref{ppg:log:eq:fi} and \eqref{ppg:log:eq:Hm} now gives
\begin{align*}
 f_i(v)&=1-\frac{v^2}{2i}+\frac{v^3/3+v^4/8}{i^2}+O(i^{-3}),\quad 0\le v\le1,\\
 I_i&=1-\frac1{6i}+\frac{13}{120i^2}+O(i^{-3}),\\
 A_i&=1+\frac5{12i}-\frac{47}{288i^2}+O(i^{-3}),\\
 \log H_m&=\frac q m+\frac2{3m^2}+O(m^{-3}).
\end{align*}
All remainders are uniform, since the Taylor variables lie in fixed bounded intervals and \eqref{ppg:log:eq:Stirlingrefined} is uniform for integers; enlarging constants covers the finitely many small indices. Multiplication gives $b^{(2)}=-47/288+13/120-5/72=-179/1440$, and exponentiation gives $F_2=2/3+q^2/2=217/288$.
\end{proof}

\begin{lemma}[The first boundary derivative]\label{ppg:log:lem:boundaryderivative}
For $i\ge1$, define
\begin{equation}\label{ppg:log:eq:Di}
 \mathcal D_i=\lim_{j\to\infty}j(\beta_{ij}-b_i)
       =-\frac{A_i}2\int_0^1f_i(v)(1-v)^2\,\dd v.
\end{equation}
At zero put $A_0=\sqrt{2\pi}$, $f_0(v)=e^{-v}$; the same integral gives
\begin{equation}\label{ppg:log:eq:D0}
 \mathcal D_0=\sqrt{2\pi}(e^{-1}-\tfrac12)
             =\lim_{j\to\infty}j(\beta_{0j}-b_0).
\end{equation}
Moreover
\begin{equation}\label{ppg:log:eq:Dexp}
 \mathcal D_i=-\frac16-\frac{11}{180i}+O(i^{-2})\qquad(i\ge1).
\end{equation}
\end{lemma}
\begin{proof}
For fixed $i\ge1$, expand the $j$ factors in the exact identity \eqref{ppg:log:eq:modalexact} or \eqref{ppg:log:eq:factor}:
\[
 A_j=1+\frac5{12j}+O(j^{-2}),\quad
 H_{i+j}=1-\frac5{12j}+O_i(j^{-2}),\quad
 f_j(1-v)=1-\frac{(1-v)^2}{2j}+O(j^{-2}).
\]
The first-order terms of $A_jH_{i+j}$ cancel. Integrating the remaining term proves \eqref{ppg:log:eq:Di}. At zero the exact endpoint expression \eqref{ppg:log:eq:exactendpoint} expands to \eqref{ppg:log:eq:D0}; alternatively its integral is $\int_0^1e^{-v}(1-v)^2\,\dd v=1-2/e$. These definitions at zero are needed only for the boundary limit.

For growing $i$, insert the first uniform expansion of $f_i$ and $A_i$ into \eqref{ppg:log:eq:Di}. Since $\int_0^1(1-v)^2\,\dd v=1/3$ and $\int_0^1v^2(1-v)^2\,\dd v=1/30$, the constant is $-1/6$ and the inverse-index coefficient is
\[
 -\frac12\left(\frac5{36}-\frac1{60}\right)=-\frac{11}{180}.
\]
The uniform remainder under the integral is $O(i^{-2})$.
\end{proof}

Retaining the $F_2/m^2$ term in the proof of Proposition~\ref{ppg:log:prop:matched} gives, for $i,j\ge1$,
\begin{equation}\label{ppg:log:eq:matchedsecond}
 \beta_{ij}=b_i+b_j-1+\frac q m+\frac s{ij}+\frac{F_2}{m^2}
                    +\widetilde R_{ij},\qquad
 |\widetilde R_{ij}|\le C M_{ij},\quad m=i+j.
\end{equation}
Here $M_{ij}=i^{-2}j^{-1}+i^{-1}j^{-2}$. In particular the previous unspecific $O(m^{-2})$ remainder has been removed. The newly omitted terms $m^{-3}$ and $m^{-2}(i^{-1}+j^{-1})$ are bounded respectively by $M_{ij}/16$ and $M_{ij}/4$; all earlier omitted terms already obeyed this mixed bound. Define $\widetilde R$ at the two endpoints by the same formula without $s/(ij)$.

Put
\begin{equation}\label{ppg:log:eq:theta}
 \theta_i=\mathcal D_i-p-q-s/i
          =\mathcal D_i+\frac16+\frac{11}{180i}\ (i\ge1),\qquad
 \theta_0=\mathcal D_0+\frac16.
\end{equation}
Then $\theta_i=O(i^{-2})$ and, for each fixed $i\ge0$,
\begin{equation}\label{ppg:log:eq:Rboundarylimit}
 n\widetilde R_{i,n-i}\longrightarrow\theta_i.
\end{equation}
Indeed Lemma~\ref{ppg:log:lem:boundaryderivative} handles $\beta_{i,n-i}-b_i$, while $b_{n-i}-1=p/(n-i)+O_i(n^{-2})$, and $n/(n-i)\to1$. The $F_2/n^2$ term vanishes after multiplication by $n$.

\begin{lemma}[Full diagonal remainder limit]\label{ppg:log:lem:diagonallimit}
The series $\mathcal Q=\sum_{i\ge0}\theta_ix_i$ converges absolutely, and
\begin{equation}\label{ppg:log:eq:diagonallimit}
 n\sum_{i=0}^n\widetilde R_{i,n-i}x_ix_{n-i}\longrightarrow2\mathcal Q.
\end{equation}
\end{lemma}
\begin{proof}
Boundedness of $x$ and \eqref{ppg:log:eq:theta} give absolute convergence. On $1\le i\le n/2$, \eqref{ppg:log:eq:matchedsecond} gives
\[
 n|\widetilde R_{i,n-i}|
 \le C\left(\frac n{i^2(n-i)}+\frac n{i(n-i)^2}\right)
 \le\frac{4C}{i^2}.
\]
This is summable independently of $n$. The zero endpoint is bounded after multiplication by $n$ by \eqref{ppg:log:eq:D0}. Use dominated convergence on the strict half-sum with \eqref{ppg:log:eq:Rboundarylimit} and $x_{n-i}\to1$, then symmetry. The even-index midpoint contributes only $O(n^{-2})$ after multiplication by $n$. This proves the limit including both boundary layers.
\end{proof}

Define the absolutely convergent sums
\begin{align}
 S_y&=y_0+\sum_{i\ge1}(y_i-d/i),&
 S_a&=a_0+\sum_{i\ge1}(a_i-p/i),\label{ppg:log:eq:residuesums}\\
 K_1&=\sum_{i\ge1}\frac{y_i}{i},&
 \mathcal Q&=\sum_{i\ge0}\theta_ix_i,
 \qquad a_i=(b_i-1)x_i.\nonumber
\end{align}
Here $K_1$ differs from the shifted sum $K$ used locally in Section~\ref{ppg:log:sec:first}. Absolute convergence follows from \eqref{ppg:log:eq:residuebound}, \eqref{ppg:log:eq:bsecond}, and \eqref{ppg:log:eq:theta}. More precisely,
\begin{equation}\label{ppg:log:eq:asecond}
 a_i=\frac p i+\frac{a^{(2)}}{i^2}
           +O\!\left(\frac{\log(i+2)}{i^3}\right),\qquad
 a^{(2)}=b^{(2)}+pd=-\frac{269}{1440}.
\end{equation}
Consequently the harmonic expansion and subtraction of the residue tail give
\begin{equation}\label{ppg:log:eq:apartialsecond}
 \sum_{i=0}^na_i=p\log n+p\gamma+S_a
                 +\frac{p/2-a^{(2)}}n+o(n^{-1}).
\end{equation}
The sign of the last coefficient follows from
$\sum_{i>n}(a_i-p/i)=a^{(2)}/n+o(n^{-1})$.

\begin{lemma}[Summable-residue convolution]\label{ppg:log:lem:residueconv}
Put $l_0=0$, $l_i=1/i$ for $i\ge1$. If $t_i$ is absolutely summable and $t_i=O(\log(i+2)/(i+1)^2)$, then
\[
 (t*l)_n=\frac{\sum_{i\ge0}t_i}{n}+o(n^{-1}).
\]
If $t,u$ both obey these conditions, $(t*u)_n=O(\log n/n^2)=o(n^{-1})$. Thus if $v=\alpha l+t$ and $w=\delta l+u$, then
\begin{equation}\label{ppg:log:eq:residueconv}
 (v*w)_n=\frac{2\alpha\delta\log n+2\alpha\delta\gamma
                    +\delta\sum t_i+\alpha\sum u_i}{n}+o(n^{-1}).
\end{equation}
\end{lemma}
\begin{proof}
In $n\sum_{i=0}^{n-1}t_i/(n-i)$, restrict first to $i\le n/2$. There $n/(n-i)\le2$, so dominated convergence gives $\sum t_i$. On the other half, $|t_i|\le C\log n/n^2$, so the absolute contribution is $O(\log n\,\HH_n/n)=o(1)$. For $t*u$, on each half bound the large-index factor by $C\log n/n^2$ and sum the absolute values of the other factor. Finally $(l*l)_n=2\HH_{n-1}/n=(2\log n+2\gamma)/n+O(n^{-2})$. Expanding the convolution proves \eqref{ppg:log:eq:residueconv}. No positivity of the residues is needed.
\end{proof}

Apply this lemma to $a=pl+(a-pl)$, $y=dl+(y-dl)$, and to the sequence with value $x_i/i$ at $i\ge1$ and value zero at $i=0$. The results are
\begin{align}
 \sum_{i=0}^na_iy_{n-i}
 &=\frac{2pd\log n+2pd\gamma+dS_a+pS_y}{n}+o(n^{-1}),\label{ppg:log:eq:convayconstant}\\
 \sum_{i=0}^nx_ix_{n-i}
 &=n+1+2d\log n+2d\gamma+2S_y+o(1),\label{ppg:log:eq:convxxconstant}\\
 \sum_{i=1}^{n-1}\frac{x_ix_{n-i}}{i(n-i)}
 &=\frac{2\log n+2\gamma+2K_1}{n}+o(n^{-1}),\label{ppg:log:eq:convxxweightedconstant}\\
 z_n&=\frac{2d^2\log n+2d^2\gamma+2dS_y}{n}+o(n^{-1}).\label{ppg:log:eq:convyconstant}
\end{align}
For \eqref{ppg:log:eq:convxxconstant}, expand $x=1+y$, use $\sum_{i\le n}y_i=d\log n+d\gamma+S_y+o(1)$, and note that $z_n=o(1)$.

The refined kernel and Lemma~\ref{ppg:log:lem:diagonallimit} improve \eqref{ppg:log:eq:eforcingmatched} to
\begin{align}\label{ppg:log:eq:eforcingconstant}
 e_n={}&2\sum_{i=0}^na_ix_{n-i}
       +\frac q n\sum_{i=0}^nx_ix_{n-i}
       +s\sum_{i=1}^{n-1}\frac{x_ix_{n-i}}{i(n-i)}\nonumber\\
       &+\frac{F_2+2\mathcal Q}{n}+o(n^{-1}).
\end{align}
Indeed the pure term $(F_2/n^2)\sum x_ix_{n-i}$ is $F_2/n+o(n^{-1})$. Substituting \eqref{ppg:log:eq:apartialsecond} and \eqref{ppg:log:eq:convayconstant}--\eqref{ppg:log:eq:convyconstant} yields
\begin{equation}\label{ppg:log:eq:hfull}
 h_n=\tfrac12\log n+B+L_h\frac{\log n}{n}+\frac{C_h}{n}+o(n^{-1}),
 \qquad L_h=-\frac7{180},
\end{equation}
where the complete constant assembly is
\begin{align}
 C_h={}&p-2a^{(2)}+q+F_2+L_h\gamma
       +2dS_a+(2p+2q+2d)S_y+2sK_1+2\mathcal Q\nonumber\\
 ={}&\frac{461}{480}-\frac{7\gamma}{180}
       -\frac{S_a}2-\frac{5S_y}6-\frac{11K_1}{90}+2\mathcal Q.
 \label{ppg:log:eq:Ch}
\end{align}
Every sum in this expression has been justified absolutely before it is used.

\begin{lemma}[Critical inverse at constant second order]\label{ppg:log:lem:Tsecond}
For $\ell_k=\log(k+1)$, $g_k=\log(k+1)/(k+1)$, and $w_k=1/(k+1)$,
\begin{align}
 (T\ell)_n&=-\frac1{2n}+\frac1{6n^2}+O(n^{-3}),\label{ppg:log:eq:Tlogsecond}\\
 (Tg)_n&=\frac{\log n}{3n^2}-\frac2{9n^2}+O(\log n/n^3),\label{ppg:log:eq:Tgsecond}\\
 (Tw)_n&=\frac1{3n^2}+O(n^{-3}).\label{ppg:log:eq:Twsecond}
\end{align}
Also $T(o((k+1)^{-1}))_n=o(n^{-2})$ for arbitrary coefficientwise remainders.
\end{lemma}
\begin{proof}
In \eqref{ppg:log:eq:Tlogexact}, expand $\log(1+1/k)=1/k-1/(2k^2)+O(k^{-3})$. The $1/k$ part telescopes exactly, and integral comparison gives
\[
 \sum_{k\ge n+1}\frac1{k^2(k+1)(k+2)}=\frac1{3n^3}+O(n^{-4}).
\]
Thus
\[
 (T\ell)_n=-\frac{\log(1+1/n)}{n+2}
       -\frac1{2(n+2)}+\frac{n+1}{6n^3}+O(n^{-3}),
\]
which proves the constant $+1/6$ in \eqref{ppg:log:eq:Tlogsecond}. The more precise tail integral already evaluated in Lemma~\ref{ppg:log:lem:Tg} gives \eqref{ppg:log:eq:Tgsecond}, retaining its constant $1/(9n^3)$ before multiplication by $-2(n+1)$. The same comparison with $t^{-4}$ instead of $\log t/t^4$ proves \eqref{ppg:log:eq:Twsecond}. For an arbitrary little-o input, bound its entire tail by $\epsilon/(k+1)$ and apply the absolute weights in \eqref{ppg:log:eq:T}; the result is $O(\epsilon/n^2)$, as in \eqref{ppg:log:eq:TO}. Let $\epsilon\downarrow0$.
\end{proof}

\begin{theorem}[Full second-order expansion]\label{ppg:log:thm:full}
There is a finite real constant $B_2$, explicitly
\begin{equation}\label{ppg:log:eq:B2}
 B_2=-\frac{59}{12960}-\frac{7\gamma}{540}-\frac{S_a}6
               -\frac{5S_y}{18}-\frac{11K_1}{270}+\frac23\mathcal Q,
\end{equation}
with the absolutely convergent definitions \eqref{ppg:log:eq:residuesums}, such that
\begin{equation}\label{ppg:log:eq:full}
 \frac{W_n}{A\sqrt n\,r_*^n}
 =1+\frac1{4n}-\frac7{540}\frac{\log n}{n^2}
                       +\frac{B_2}{n^2}+o(n^{-2}).
\end{equation}
\end{theorem}
\begin{proof}
When expressing \eqref{ppg:log:eq:hfull} in the basis of Lemma~\ref{ppg:log:lem:Tsecond}, retain the shift correction
\[
 \log k=\ell_k-w_k+O(k^{-2}),\quad
 \frac{\log k}{k}=g_k+o(k^{-1}),\quad
 \frac1k=w_k+O(k^{-2}).
\]
Finite initial changes are irrelevant to $T$ at large indices. Hence
\[
 h_k=\tfrac12\ell_k+B+L_hg_k+(C_h-\tfrac12)w_k+o(k^{-1}),
\]
and Lemma~\ref{ppg:log:lem:Tsecond} gives
\[
 y_n=-\frac1{4n}-\frac7{540}\frac{\log n}{n^2}
        +\frac{C_x}{n^2}+o(n^{-2}),\qquad
 C_x=\frac{C_h}3-\frac1{12}-\frac{2L_h}9.
\]
Multiplying by $\sqrt{1+1/n}=1+1/(2n)-1/(8n^2)+O(n^{-3})$ changes the constant by $d/2-1/8=-1/4$. Thus
$B_2=C_h/3-1/3-2L_h/9$, and substitution of \eqref{ppg:log:eq:Ch} gives \eqref{ppg:log:eq:B2}.
\end{proof}

\begin{corollary}[Eventual strict decrease of the ratios]\label{ppg:log:cor:monotonicity}
For $\rho_n=W_{n+1}/W_n$,
\begin{equation}\label{ppg:log:eq:ratiosecond}
 \rho_n=r_*\left(1+\frac1{2n}-\frac3{8n^2}+o(n^{-2})\right),
 \qquad \rho_{n+1}-\rho_n=-\frac{r_*}{2n^2}+o(n^{-2}).
\end{equation}
Consequently the ratios are strictly decreasing for all sufficiently large $n$, and $W_n$ is eventually strictly log-concave. No effective onset or all-index monotonicity is asserted.
\end{corollary}
\begin{proof}
Let $U_n$ denote the left-hand side of \eqref{ppg:log:eq:full}. The explicit logarithmic and constant second terms change by $O(\log n/n^3)$ between consecutive indices. If $\varepsilon_n=o(n^{-2})$ is its remainder, then
$|\varepsilon_{n+1}-\varepsilon_n|\le|\varepsilon_{n+1}|+|\varepsilon_n|=o(n^{-2})$.
It follows that $U_{n+1}/U_n=1-1/(4n^2)+o(n^{-2})$. Multiply by $r_*\sqrt{1+1/n}$ to obtain the first formula in \eqref{ppg:log:eq:ratiosecond}. Its new remainder again has consecutive difference $o(n^{-2})$ by the same absolute bound, whereas the difference of $1/(2n)$ is $-1/(2n^2)+O(n^{-3})$. This proves the second formula and its eventual negative sign. This use of differences is at the unchanged $n^{-2}$ error scale; it does not infer an $n^{-3}\log n$ ratio coefficient by differentiating an uncontrolled error.
\end{proof}
\begin{writenote}[5 October 2026]
Part~\ref{ppg:all:part}, Corollary~\ref{ppg:all:cor:ratio}, extends \eqref{ppg:log:eq:ratiosecond} by its cubic term $\frac{(7/270)\log n+32/135-2B_2}{n^3}$, which answers the first half of item~3 of Section~\ref{ppg:log:sec:open}. Strict decrease at every index remains open. The write's computation of $x_n$ in 60-digit fixed point (Section~\ref{ppg:sec:further}, question~\ref{ppg:fq:sa}) finds $\rho_0>\rho_1>\cdots>\rho_{3999}$, and the placement dossier found the same in float64 for $n\le4095$; this is evidence, not proof. The same computation gives $B_2\approx-0.0573920672$, uncertified.
\end{writenote}

\section{An obstruction to pure inverse-power ratio expansions}\label{ppg:log:sec:ratio}
\begin{theorem}\label{ppg:log:thm:ratio}
For $\rho_n=W_{n+1}/W_n$, there do not exist finite constants $a_1,a_2,a_3$ such that
\begin{equation}\label{ppg:log:eq:ratiosuppose}
 \rho_n=r_*+\frac{a_1}{n}+\frac{a_2}{n^2}+\frac{a_3}{n^3}+o(n^{-3}).
\end{equation}
The same obstruction holds for the shifted convention $W_n/W_{n-1}$.
\end{theorem}
\begin{proof}
Assume \eqref{ppg:log:eq:ratiosuppose}. Since $r_*>0$, Taylor expansion of the logarithm gives constants $b_1,b_2,b_3$ with
\[
 \log\rho_n=\log r_*+\frac{b_1}{n}+\frac{b_2}{n^2}
                    +\frac{b_3}{n^3}+t_n,\qquad t_n=o(n^{-3}).
\]
The remainder series converges absolutely and its tail is $o(n^{-2})$: for every $\epsilon>0$, bound $|t_k|\le\epsilon k^{-3}$ for all sufficiently large $k$. Sum from a fixed such index $N$ through $n-1$. The elementary harmonic expansions are
\begin{align*}
 \sum_{k=N}^{n-1}\frac1k&=\log n+C_N-\frac1{2n}-\frac1{12n^2}+O(n^{-4}),\\
 \sum_{k=N}^{n-1}\frac1{k^2}&=C'_N-\frac1n-\frac1{2n^2}+O(n^{-3}),\\
 \sum_{k=N}^{n-1}\frac1{k^3}&=C''_N-\frac1{2n^2}+O(n^{-3}).
\end{align*}
It follows that there are constants $C,c_1,c_2$ with
\begin{equation}\label{ppg:log:eq:logpure}
 \log W_n=n\log r_*+b_1\log n+C+\frac{c_1}{n}
                             +\frac{c_2}{n^2}+o(n^{-2}).
\end{equation}
But Theorem~\ref{ppg:log:thm:main}, followed by a logarithm, gives
\begin{equation}\label{ppg:log:eq:logsharp}
 \log W_n=n\log r_*+\tfrac12\log n+\log A+\frac1{4n}
  -\frac7{540}\frac{\log n}{n^2}+o(\log n/n^2).
\end{equation}
The quadratic term in that last logarithm is $O(n^{-2})$, absorbed by its stated remainder. Comparing \eqref{ppg:log:eq:logpure} and \eqref{ppg:log:eq:logsharp} successively forces $b_1=1/2$, $C=\log A$, and $c_1=1/4$. Subtract these matching terms and divide by $\log n/n^2$. Equation \eqref{ppg:log:eq:logpure} tends to 0, while \eqref{ppg:log:eq:logsharp} tends to $-7/540$, a contradiction. Shifting the ratio index by one preserves the hypothetical pure-power form and its remainder order.
\end{proof}

\subsection{What the obstruction says about Richardson extrapolation}
In the inspected version of Serra~\cite{serra}, Section 6.4, printed page 49 (PDF page 50), equation (151) assumes a pure inverse-power expansion of $r(n)=W(n+1)/W(n)$ for Richardson extrapolation. The accompanying discussion explicitly identifies formal proof of the assumption as an open problem. Theorem~\ref{ppg:log:thm:ratio} shows that interpreting that premise as a full Poincare expansion is incompatible with the nonzero sequence-level logarithm proved here. Thus that pure-power asymptotic premise cannot provide the proposed justification at all orders.

This is not a proof that the decimal digits obtained numerically are wrong. Numerical accuracy, asymptotic model validity, and certified error bounds are distinct issues. The obstruction permits a pure-power ratio approximation through order $n^{-2}$ with $o(n^{-2})$ remainder; Corollary~\ref{ppg:log:cor:monotonicity} in fact establishes that approximation. Nor does it identify the exact ratio-level logarithmic correction: differencing an uncontrolled coefficientwise little-o remainder would be invalid. The proof above instead sums a hypothetical \emph{stronger} ratio expansion, whose tail error is controlled. The source comparison is based on the inspected version-1 PDF; no claim of fresh live access to its hosting record is needed.
\begin{writenote}[5 October 2026]
Theorem~\ref{ppg:log:thm:ratio} also follows from Part~\ref{ppg:all:part}'s Corollary~\ref{ppg:all:cor:ratio}: subtracting \eqref{ppg:all:eq:ratio3}, multiplied by $r_*$, from \eqref{ppg:log:eq:ratiosuppose} forces $a_1=r_*/2$ and $a_2=-3r_*/8$ and then leaves $\frac{7r_*}{270}\frac{\log n}{n^3}=O(n^{-3})$, which is false. The proof above, which sums a hypothetical expansion instead, came first and is independent; it is kept. The ratio-level logarithm that this subsection leaves unidentified is $\frac{7r_*}{270}\frac{\log n}{n^3}$ (Corollary~\ref{ppg:all:cor:ratio}). The write checked (151) on Serra's printed page~49. Remark~\ref{ppg:all:rem:op2} records that his Open Problem~2, the expansion~(199), fails for every $p\ge2$.
\end{writenote}

\section{A refined ceiling-safe probability threshold}\label{ppg:log:sec:inverse}
Let
\[
 N(\varepsilon)=\min\{n\ge0:W_n\le\varepsilon\},\qquad
 L=\log(1/\varepsilon),\qquad\lambda=\log(1/r_*)>0.
\]
The set is nonempty for $\varepsilon>0$ because $r_*<1$. Write $W_{-1}$ for the real branch of the Lambert-W function with values at most $-1$ on $[-e^{-1},0)$, so that $W_{-1}(z)e^{W_{-1}(z)}=z$.

\begin{theorem}[Inverse with the full second-order refinement]\label{ppg:log:thm:inverse}
Put $b=1/4$, $\kappa=-7/540$, and $k=B_2-b^2/2=B_2-1/32$. For sufficiently small $\varepsilon>0$, define
\begin{align}
 x_0&=-\frac1{2\lambda}W_{-1}\!\left(-\frac{2\lambda\varepsilon^2}{A^2}\right),
 &D_0&=\lambda-\frac1{2x_0},\label{ppg:log:eq:x0}\\
 x_3&=x_0+\frac{b}{x_0D_0}
             +\frac{\kappa\log x_0+k}{x_0^2D_0}.
 \label{ppg:log:eq:x3}
\end{align}
Then $D_0>0$ and there is $\eta_\varepsilon\ge0$, tending to zero as $\varepsilon\downarrow0$, such that
\begin{equation}\label{ppg:log:eq:ceilinverse}
 \left\lceil x_3-\frac{\eta_\varepsilon}{x_0^2}\right\rceil
 \le N(\varepsilon)\le
 \left\lceil x_3+\frac{\eta_\varepsilon}{x_0^2}\right\rceil.
\end{equation}
The uncertainty is asserted before rounding; the theorem does not assert $N(\varepsilon)-x_3\to0$ or provide an effective bound for $\eta_\varepsilon$.
\end{theorem}
\begin{proof}
For large real $t$, put
\[
 f_0(t)=\lambda t-\tfrac12\log t-\log A,\qquad
 h(t)=\frac b t+\frac{\kappa\log t+k}{t^2},\qquad f(t)=f_0(t)-h(t).
\]
Theorem~\ref{ppg:log:thm:full} and Taylor's formula for the logarithm imply
\begin{equation}\label{ppg:log:eq:inverseerror}
 \log W_n=-f(n)+r_n,\qquad r_n=o(n^{-2}).
\end{equation}
The logarithm's quadratic contribution is exactly $-b^2/(2n^2)$, which explains the necessary change from $B_2$ to $k$. Terms of higher degree are $o(n^{-2})$.

By the definition of the large Lambert-W branch, $f_0(x_0)=L$, $x_0\to\infty$, and $f_0'(x_0)=D_0>0$ eventually. Put $\Delta=x_3-x_0=h(x_0)/D_0=O(x_0^{-1})$. Taylor expansion on this short interval gives
\begin{align*}
 f_0(x_3)&=L+D_0\Delta+O(x_0^{-2}\Delta^2)
          =L+h(x_0)+O(x_0^{-4}),\\
 h(x_3)&=h(x_0)+O(x_0^{-3}),
\end{align*}
because $h'(t)=O(x_0^{-2})$ there. Therefore
\begin{equation}\label{ppg:log:eq:inverseresidual}
 f(x_3)-L=O(x_0^{-3})=o(x_0^{-2}).
\end{equation}
Also $f'(t)\ge\lambda/2$ for all sufficiently large $t$. These derivatives are of explicit comparison functions, not of the unknown probability sequence or its remainder.

The tail supremum in \eqref{ppg:log:eq:inverseerror} gives $s_\varepsilon\to0$ such that
\begin{equation}\label{ppg:log:eq:uniforminverseerror}
 |r_n|\le s_\varepsilon/x_0^2\qquad\text{for every integer }n\ge x_0/2.
\end{equation}
For example, a constant multiple of $\sup_{n\ge\lfloor x_0/2\rfloor}(n+1)^2|r_n|$ suffices. Choose a strictly positive $\eta_\varepsilon\to0$ large enough that, with $\delta_\varepsilon=\eta_\varepsilon/x_0^2$,
\begin{equation}\label{ppg:log:eq:inversemargin}
 (\lambda/2)\delta_\varepsilon>|f(x_3)-L|+s_\varepsilon/x_0^2.
\end{equation}
Such a choice exists by \eqref{ppg:log:eq:inverseresidual}. For $x_0/2\le n<x_3-\delta_\varepsilon$, the derivative lower bound and \eqref{ppg:log:eq:uniforminverseerror} give
\[
 \log W_n\ge-L-|f(x_3)-L|+(\lambda/2)\delta_\varepsilon
                  -s_\varepsilon/x_0^2>-L.
\]
At $n=\lceil x_3+\delta_\varepsilon\rceil$, the reversed inequalities give $\log W_n<-L$.

It remains to exclude earlier indices $n<x_0/2$. The root limit alone implies $N(\varepsilon)\sim L/\lambda$: for any $0<\tau<\lambda$, eventually
$e^{-(\lambda+\tau)n}\le W_n\le e^{-(\lambda-\tau)n}$, and the finitely many preceding probabilities have a positive minimum. These estimates squeeze the first crossing between $L/(\lambda+\tau)$ and $L/(\lambda-\tau)$ up to bounded integer errors. Let $\tau\downarrow0$. Since $f_0(x_0)=L$ also gives $x_0\sim L/\lambda$, eventually $N(\varepsilon)>x_0/2$. The two preceding inequalities therefore prove \eqref{ppg:log:eq:ceilinverse}. No additional monotonicity hypothesis is required.
\end{proof}

\begin{remark}[The logarithmic refinement without the constant]
If one retains only Theorem~\ref{ppg:log:thm:main}, the center
\begin{equation}\label{ppg:log:eq:x2}
 x_2=x_0+\frac{1}{4x_0D_0}
              -\frac7{540}\frac{\log x_0}{x_0^2D_0}
\end{equation}
already gives
\[
 \left\lceil x_2-\eta_\varepsilon\frac{\log x_0}{x_0^2}\right\rceil
 \le N(\varepsilon)\le
 \left\lceil x_2+\eta_\varepsilon\frac{\log x_0}{x_0^2}\right\rceil,
 \qquad\eta_\varepsilon\to0.
\]
This follows from the same proof with logarithmic-scale residuals, or from \eqref{ppg:log:eq:ceilinverse}, since $x_3-x_2=O(x_0^{-2})=o(\log x_0/x_0^2)$. Both statements retain integer ceilings.
\end{remark}

\begin{remark}[\wtag\ Relation to Parts~I and~III, 5 October 2026]\label{ppg:log:rem:inverse}
The root $x_0$ of \eqref{ppg:log:eq:x0} is $t_\varepsilon$ of Part~\ref{ppg:lead:part}'s \eqref{ppg:lead:eq:lambert} and $x$ of Part~\ref{ppg:all:part}'s \eqref{ppg:all:eq:x0}, the large root that the branch rule of \texttt{p0:thm:lambert-core} \cite{tsvol} assigns to $W_{-1}$ (Remark~\ref{ppg:lead:rem:inverse}). Theorem~\ref{ppg:log:thm:inverse} follows from Part~\ref{ppg:all:part}'s Theorem~\ref{ppg:all:thm:threshold} with $M=2$, whose proof does not use it. Indeed, with $\ell=\log x_0$, the centre there is $z_2=x_0+\frac b{\lambda x_0}+\frac{(c\ell+k)/\lambda+b/(2\lambda^2)}{x_0^2}$ by \eqref{ppg:all:eq:R123}, where $c=\kappa$ and $k$ is the same constant. Since $\frac1{D_0}=\frac1\lambda\bigl(1-\frac1{2\lambda x_0}\bigr)^{-1}=\frac1\lambda+\frac1{2\lambda^2x_0}+O(x_0^{-2})$,
\[
 x_3-z_2=\frac b{x_0}\Bigl(\frac1{D_0}-\frac1\lambda-\frac1{2\lambda^2x_0}\Bigr)
  +\frac{\kappa\ell+k}{x_0^2}\Bigl(\frac1{D_0}-\frac1\lambda\Bigr)
  =O\Bigl(\frac{1+\ell}{x_0^3}\Bigr).
\]
Put $\eta_\varepsilon=x_0^2\bigl(|x_3-z_2|+C_2E_2(x_0)\bigr)=O((1+\ell)^2/x_0)\to0$. Then $x_3-\eta_\varepsilon/x_0^2\le z_2-C_2E_2(x_0)$ and $z_2+C_2E_2(x_0)\le x_3+\eta_\varepsilon/x_0^2$, and the monotonicity of the ceiling turns \eqref{ppg:all:eq:ceiling} into \eqref{ppg:log:eq:ceilinverse}. The bracket is of the kind used in the separation condition of \texttt{p0:thm:staircase}; this proof, like Part~\ref{ppg:all:part}'s, needs no monotonicity of $W_n$.
\end{remark}

\section{What remains beyond the proved terms}\label{ppg:log:sec:open}
The results now determine both the logarithmic and constant terms at second order and prove eventual ratio decrease. Several concrete questions remain.
\begin{enumerate}
\item \textbf{A full logarithmic expansion.} Can one prove a systematic expansion in inverse powers of $n$ with polynomials in $\log n$? The boundary-matched factorization suggests a route, but each new order needs uniform mixed errors and control of the corresponding nonlinear convolutions. An all-orders statement does not follow from the two orders established here.
\item \textbf{The second constant numerically.} Can the absolutely convergent expression \eqref{ppg:log:eq:B2} be turned into certified decimal enclosures for $B_2$? Absolute convergence proves existence but does not alone provide effective tail bounds for the terms involving $x_i$ and $y_i$.
\item \textbf{Finer ratio terms and finite-index behavior.} Can one establish a strong enough recurrence-based difference estimate to identify the first logarithmic term at order $n^{-3}$ in $W_{n+1}/W_n$? Can one prove strict decrease at every relevant finite index, or identify an explicit onset of the eventual decrease? The present $o(n^{-2})$ remainder answers neither question.
\item \textbf{Effective certification.} Can the tail contraction, residue sums, and endpoint estimates be made into computable correction-error bounds with explicit onsets? Such bounds could support rigorously controlled numerical extrapolation. The exact finite arithmetic below does not itself provide them or certify further digits of $r_*$.
\item \textbf{Effective inverse intervals.} Effective bounds for $B_2$ and the full remainder could make \eqref{ppg:log:eq:ceilinverse} computable with a guaranteed error. Rounding remains necessary, especially when the center is close to an integer.
\end{enumerate}
These are research questions rather than consequences silently included in the full second-order theorem.
\begin{writenote}[5 October 2026]
Status after Part~\ref{ppg:all:part}. Item~1 is answered: Theorem~\ref{ppg:all:thm:all} proves the expansion in inverse powers of $n$ with polynomials in $\log n$ to every fixed order (not its convergence). Item~2 is open (Section~\ref{ppg:sec:further}, question~\ref{ppg:fq:effective}); an uncertified value is $B_2\approx-0.0573920672$. Item~3 is answered in its first part: the first logarithmic term of $W_{n+1}/W_n$ at order $n^{-3}$ is $\frac{7r_*}{270}\log n$ (Corollary~\ref{ppg:all:cor:ratio}); strict decrease at every index and an explicit onset remain open (question~\ref{ppg:fq:monotone}). Items~4 and~5 are open (questions~\ref{ppg:fq:effective} and~\ref{ppg:fq:thresholds}).
\end{writenote}

\section{Reproduction and the boundary of computation}\label{ppg:log:sec:reproduction}
The accompanying archive contains this LaTeX source and PDF, a deterministic build script, a new exact-arithmetic companion, and the unchanged source and PDF of Report 144. It does not redistribute the primary manuscript PDF. The foundation report supplies the proofs of \eqref{ppg:log:eq:foundation} and \eqref{ppg:log:eq:firstkernel}--\eqref{ppg:log:eq:firstrho}; its earlier computational archive is not needed for the new companion and is not duplicated here.

The companion uses Python's standard library and rational arithmetic. It derives the polynomial covariance integral, signed coefficient arithmetic through the full constant term, rational identities for the critical inverse, the ratio coefficients, and formal algebraic residual identities for the refined inverse. Analytic asymptotic constants for the inverse operator are explicitly identified as proof inputs, rather than conclusions of finite testing. Typed, closed fixtures are validated against recomputed values; unexpected keys, invalid rational representations, duplicate JSON keys, and semantic changes are rejected. The test suite also checks normal and optimized Python behavior and guards against overwriting outputs. The README gives the exact commands and explains each check's scope.

Finite exact checks are useful safeguards for algebra and indexing. They are not substitutes for the proofs of uniform remainders, infinite-tail estimates, or coefficientwise little-o propagation in this article. In particular, no floating estimate of $r_*$ enters the verifier, and no finite run is described as a certificate for the asymptotic theorem or the conjectural decimal expansion. Rebuilding the PDF with the recorded toolchain and environment produces deterministic bytes; the archive builder fixes file order, timestamps, permissions, and compression settings.
\begin{writenote}[5 October 2026]
Shipped names: the checker and its tests are \file{code/145-logcorr-companion-checks.py} and \file{code/145-logcorr-companion-test_checks.py}, the build scripts \file{code/145-logcorr-build_pdf.py} and \file{code/145-logcorr-build_archive.py}, the fixture, receipt and test records \file{data/145-logcorr-companion-*.json}, and the companion README \file{145-logcorr-companion-README.md}. The archive's copy of Report~144 is printed as Part~\ref{ppg:lead:part} and not shipped; its \file{toolchain.json} is identical to Report~146's, shipped as \file{data/146-allorders-toolchain.json}; its PDF, README and checksum manifest are not shipped. On a copy of the shipped files, the checker run with \texttt{--fixture} on the shipped fixture reproduces \file{data/145-logcorr-companion-receipt.json} as JSON (rerun by the write, 5~October 2026); the README gives the route.
\end{writenote}

\part{All fixed orders for the proportional placement game}\label{ppg:all:part}
\partsource{Bundle Report~146, archive \file{A398540_All_Orders_Logarithmic_Asymptotics_and_Inverses_Source.zip}, delivered as \file{Report146.tex} (950 lines, 22 pages), dated 3~October 2026; title block ``Report 146''; no author line; the base of the merge, staged as \file{article.tex}. It is the version ``corrected in place'' before delivery; the earlier version was never received (provenance section). Printed as delivered, with \textbf{[write]} notes, except that its table of contents is replaced by the report's. Delivered Section~$k$ is Section~$k+25$ here, and its statement $k.j$ is $(k+25).j$; equations are renumbered within sections. ``Report 144'' is Part~\ref{ppg:lead:part} and ``Report 145'' is Part~\ref{ppg:log:part}. Remarks~\ref{ppg:all:rem:op2} and~\ref{ppg:all:rem:inverse} are added in the write.}
\begin{partabstract}
For the proportional placement probabilities associated with OEIS A398540, we prove a coefficientwise logarithmic Poincar\'e expansion at every fixed order. If $r_*\in(0,1)$ is the exponential rate and $A=r_*\sqrt{2\pi}$, then
\[
\frac{W_n}{A\sqrt n\,r_*^n}
 =1+\sum_{j=1}^M\frac{\widehat P_j(\log n)}{n^j}
   +O_M\!\left(\frac{(1+\log n)^M}{n^{M+1}}\right),
\qquad \deg\widehat P_j\le j-1.
\]
The proof bootstraps signed, continuous-function-valued convolutions using exact logarithmic coefficient models and absolutely convergent subtracted endpoint moments. Each step reduces to finite algebra and finitely many specified scalar series. The third-order logarithm is explicit and the possible squared logarithm at that order vanishes. We give all-fixed-order Lambert-W threshold inverses with ceiling-safe error brackets. Separately, a finite-coefficient interval algorithm has a proved quadratic-times-log radius-width bound; this is a computability theorem, not a complexity bound or a calculation of new digits. The expansion is not asserted to converge as an infinite series, and its remainder constants and inverse margins are not made effective here.
\end{partabstract}
{\small\emph{Readings in this Part} (\wtag). $x_n$, $y_n$, $\beta_{ij}$, $b_i$, $H_n$, $B$, $d$, $p=\frac14$, $q=-\frac5{12}$ as in Part~\ref{ppg:log:part}; $c=-\frac7{540}$ (Part~\ref{ppg:log:part}'s $\kappa$, not its $c$); $C=B_2+\frac14$; $f=h_2=\frac{217}{288}$; $u=1-z$ and $L=-\log(1-z)$ in Sections~\ref{ppg:all:sec:models}--\ref{ppg:all:sec:third}, where $p,q$ are also integer indices of $B_{p,q}=u^{p-1}L^q$; $\delta_i=(b_i-1)x_i$ (Part~\ref{ppg:log:part}'s $a_i$) and $a_i(v)=g_i(v)x_i-1$; $S_D$ (Part~\ref{ppg:log:part}'s $S_a$) and the new $S_A$; in Section~\ref{ppg:all:sec:inverse}, $x$ is the Lambert root and $b=\frac14$; in Section~\ref{ppg:all:sec:rate}, $u_n=W_n/\sqrt{n+1}$ and $F$ as in Part~\ref{ppg:lead:part}, and $A$, $B$, $P$, $L$, $U$ are interval endpoints.\par}
\section{The sequence and the main theorem}\label{ppg:all:sec:main}
There are $n$ ordered empty slots and $n$ independent uniform values on $(0,1)$, arriving sequentially. In a gap $(L,U)$ with $k$ available slots, the proportional policy places an arriving value $t$ in position $\lceil k(t-L)/(U-L)\rceil$ within that gap. It loses when the containing gap has no empty slot. Let $W_n$ be its winning probability, with $W_0=1$. This article makes no optimality assertion about the policy.

Conditioning on the first value and on the number of subsequent values below it gives the exact recurrence
\begin{align}
 W_n&=\frac1n\sum_{k=1}^n W_{k-1}W_{n-k}Q_{n,k},\label{ppg:all:eq:Wrec}\\
 Q_{n,k}&=n\int_{(k-1)/n}^{k/n}\binom{n-1}{k-1}
          t^{k-1}(1-t)^{n-k}\,\dd t.\label{ppg:all:eq:Q}
\end{align}
All $W_n$ are positive rationals; $W_1=1$ and $W_2=3/4$. The two subgames are independent after affine rescaling, which proves the product in \eqref{ppg:all:eq:Wrec}.
\begin{writenote}[5 October 2026]
The recurrence is that of Part~\ref{ppg:lead:part}, Proposition~\ref{ppg:lead:prop:recurrence}, where it is derived in full.
\end{writenote}

Report 144~\cite{r144} establishes
\begin{equation}\label{ppg:all:eq:leading}
 r_*:=\lim_{n\to\infty}W_n^{1/n}\in(0,1),\qquad
 W_n\sim A\sqrt n\,r_*^n,\qquad A=r_*\sqrt{2\pi}.
\end{equation}
Set
\begin{equation}\label{ppg:all:eq:normalization}
 R=r_*^{-1},\qquad a=(2\pi)^{-1/2},\qquad
 x_n=\frac{W_nR^n}{A\sqrt{n+1}},\qquad y_n=x_n-1.
\end{equation}
Report 145~\cite{r145} supplies the starting rate and the second constant:
\begin{align}
 x_n&=1+\frac d n+O\!\left(\frac{1+\log n}{n^2}\right),\quad d=-\frac14,
 \label{ppg:all:eq:base}\\
 x_n&=1+\frac d n+\frac{c\log n+C}{n^2}+o(n^{-2}),
 \quad c=-\frac7{540},\quad C=B_2+\frac14.\label{ppg:all:eq:oldsecond}
\end{align}
The rate \eqref{ppg:all:eq:base}, not merely the weaker first correction with an $o(n^{-1})$ remainder, is the induction input. The constant $B_2$ is defined by absolutely convergent series in Report 145. Its decimal value is not needed.
\begin{writenote}[5 October 2026]
Report~144 is Part~\ref{ppg:lead:part} and Report~145 is Part~\ref{ppg:log:part}. \eqref{ppg:all:eq:leading} is Theorems~\ref{ppg:lead:thm:rate} and~\ref{ppg:lead:thm:pointwise}. \eqref{ppg:all:eq:base} follows from Part~\ref{ppg:log:part}'s \eqref{ppg:log:eq:xsharp}, which gives $y_n+\frac1{4n}=O(\log n/n^2)$, the finitely many small $n$ being absorbed in the constant. \eqref{ppg:all:eq:oldsecond} is Theorem~\ref{ppg:log:thm:full}, with $C=B_2+\frac14$ the constant $C_x$ of its proof and $B_2$ given by \eqref{ppg:log:eq:B2}.
\end{writenote}

\begin{theorem}[Every fixed order]\label{ppg:all:thm:all}
There is a unique sequence of real polynomials $P_j$ such that, for each fixed integer $M\ge1$,
\begin{equation}\label{ppg:all:eq:all}
 x_n=1+\sum_{j=1}^M\frac{P_j(\log n)}{n^j}
       +O_M\!\left(\frac{(1+\log n)^M}{n^{M+1}}\right),
 \qquad \deg P_j\le j-1.
\end{equation}
Here $P_1=d$ and $P_2(t)=ct+C$. Each order is determined by the finite scalar construction in Section~\ref{ppg:all:sec:scalar}. Multiplication by $\sqrt{1+1/n}$ gives a consistent expansion
\begin{equation}\label{ppg:all:eq:allW}
 \frac{W_n}{A\sqrt n\,r_*^n}
 =1+\sum_{j=1}^M\frac{\widehat P_j(\log n)}{n^j}
       +O_M\!\left(\frac{(1+\log n)^M}{n^{M+1}}\right),
 \qquad \deg\widehat P_j\le j-1.
\end{equation}
In particular $\widehat P_1=1/4$ and $\widehat P_2(t)=ct+B_2$.
\end{theorem}

Theorem~\ref{ppg:all:thm:third} gives the explicit third-order term. Theorem~\ref{ppg:all:thm:threshold} inverts \eqref{ppg:all:eq:allW} at every fixed order. The separate computable enclosure theorem in Section~\ref{ppg:all:sec:rate} concerns the leading rate only. Every big-O constant may depend on the fixed order. No estimate is uniform as $M\to\infty$.
\begin{writenote}[5 October 2026]
This theorem answers item~1 of Part~\ref{ppg:log:part}'s Section~\ref{ppg:log:sec:open} and the question ``Pointwise corrections and logarithms'' of Part~\ref{ppg:lead:part}'s Section~\ref{ppg:lead:sec:questions}, in the fixed-order sense fixed in Section~\ref{ppg:all:sub:scope}. It supersedes Theorems~\ref{ppg:log:thm:main} and~\ref{ppg:log:thm:full} in statement but takes Part~\ref{ppg:log:part}'s results as inputs.
\end{writenote}

\subsection{Sources and the meaning of all orders}\label{ppg:all:sub:scope}
The model and recurrence are in Serra's manuscript, version 1, Theorem 3.12 and equations (47)--(48), printed pages 23--24~\cite{serra}. OEIS A398540 records digits associated with the limiting rate, rather than the probability sequence itself~\cite{oeis}. The present article is a continuation of Reports 144 and 145; their unchanged PDF and LaTeX files accompany it. We use their proved input statements, rather than reusing their computational archives. The primary manuscript PDF is not redistributed.
\begin{writenote}[5 October 2026]
The unchanged files of Reports~144 and~145 that accompanied this manuscript are printed here as Parts~\ref{ppg:lead:part} and~\ref{ppg:log:part} and are not shipped; they are byte-identical to those Reports' own deliveries. The citations of Serra's manuscript were checked by the write against its version-1 PDF.
\end{writenote}

All logarithms are natural. A full fixed-order expansion means the quantified statement \eqref{ppg:all:eq:all} for each finite $M$, with consistent coefficient polynomials. It does not mean a convergent infinite expansion, a complete exponentially small transseries, a proof of analytic continuation, or a universal novelty claim. The proofs below are coefficientwise and work without extending the actual generating function through its circle of convergence.

\section{The exact kernel and its uniform expansion}\label{ppg:all:sec:kernel}
For $i,j\ge0$, with $n=i+j$, define
\begin{equation}\label{ppg:all:eq:beta}
 \beta_{ij}=\frac{Q_{n+1,i+1}}a
             \sqrt{\frac{(i+1)(j+1)}{n+2}}.
\end{equation}
Since $aAR=1$, \eqref{ppg:all:eq:Wrec} becomes
\begin{equation}\label{ppg:all:eq:xrec}
 (n+1)x_{n+1}=\sum_{i+j=n}\beta_{ij}x_ix_j.
\end{equation}
For $i\ge1$ let $r_i$ be the exact Stirling remainder defined by
\[
 i!=\sqrt{2\pi i}(i/e)^i e^{r_i}.
\]
Put
\begin{align}
 g_i(v)&=\sqrt{1+1/i}\,e^{-r_i}(1+v/i)^i e^{-v},\quad i\ge1,
 \label{ppg:all:eq:g}\\
 g_0(v)&=\sqrt{2\pi}\,e^{-v},\qquad
 b_i=\int_0^1g_i(v)\,\dd v,\label{ppg:all:eq:gzero}\\
 H_n&=(1+2/n)^{-1/2}\exp\{r_n+1-n\log(1+1/n)\},\quad n\ge1.
 \label{ppg:all:eq:H}
\end{align}

\begin{lemma}[Separated kernel including endpoints]\label{ppg:all:lem:factor}
For $i,j\ge0$ with $i+j=n\ge1$,
\begin{equation}\label{ppg:all:eq:factor}
 \beta_{ij}=H_n\int_0^1g_i(v)g_j(1-v)\,\dd v.
\end{equation}
\end{lemma}
\begin{proof}
For positive child indices this is the exact factorization obtained from \eqref{ppg:all:eq:Q} by the change of variables $t=(i+v)/(n+1)$ and the factorial formula. For completeness at a zero child, define locally
$f_n(v)=(1+v/n)^n e^{-v}$. Then
\begin{align*}
 H_n\sqrt{1+1/n}\,e^{-r_n}
 &=\sqrt{\frac{n+1}{n+2}}e^{1-n\log(1+1/n)},\\
 \int_0^1e^{-v}f_n(1-v)\,\dd v
 &=e^{-1}\frac n{n+1}\{(1+1/n)^{n+1}-1\}.
\end{align*}
Multiplication by $\sqrt{2\pi}$ yields
\[
 \beta_{0,n}=\sqrt{2\pi}\sqrt{\frac{n+1}{n+2}}
      \left\{1-\left(\frac n{n+1}\right)^{n+1}\right\},
\]
which is also the value from the original integral. Symmetry treats the other endpoint. The isolated pair $(0,0)$ is never needed in \eqref{ppg:all:eq:factor}.
\end{proof}
\begin{writenote}[5 October 2026]
For positive children this is Part~\ref{ppg:log:part}'s Lemma~\ref{ppg:log:lem:factor} ($g_i=A_if_i$ there). The zero-child value $\beta_{0,n}$ is Part~\ref{ppg:log:part}'s \eqref{ppg:log:eq:exactendpoint}, and Part~\ref{ppg:lead:part}'s \eqref{ppg:lead:eq:endpoint} divided by $a$ at parent size $n+1$. The coefficients $p$, $b_2^*$, $q$, $f$ of \eqref{ppg:all:eq:b12}--\eqref{ppg:all:eq:H12} are Part~\ref{ppg:log:part}'s $p$, $b^{(2)}$, $q$, $F_2$ (\eqref{ppg:log:eq:bHexp}, \eqref{ppg:log:eq:Hsecond}--\eqref{ppg:log:eq:bsecond}).
\end{writenote}

\begin{lemma}[Uniform finite kernel expansions]\label{ppg:all:lem:kernel}
For every fixed $K\ge1$ there are explicitly computable polynomials $G_k(v)$ and numbers $h_k$ for which
\begin{align}
 g_n(v)&=1+\sum_{k=1}^K G_k(v)n^{-k}+O_K(n^{-K-1}),\label{ppg:all:eq:gexp}\\
 H_n&=1+\sum_{k=1}^K h_kn^{-k}+O_K(n^{-K-1}),\label{ppg:all:eq:Hexp}
\end{align}
uniformly for $v\in[0,1]$ and integers $n\ge1$.
\end{lemma}
\begin{proof}
The exact difference and limit
\[
 r_n-r_{n+1}=(n+\tfrac12)\log(1+1/n)-1,\qquad r_n\to0,
\]
determine a finite expansion $r_n=\sum_{k=1}^K r[k]n^{-k}+O(n^{-K-1})$. Match differences successively through $n^{-K-1}$: the new coefficient $r[k]$ has nonzero pivot $k n^{-k-1}$. The unmatched difference is $O(n^{-K-2})$, whose absolutely convergent tail is $O(n^{-K-1})$. This proves the error instead of presupposing a full expansion of $r_n$.

For $t=1/n$ and the just-constructed finite series $\mathcal R(t)=\sum r[k]t^k$, the required recipes are the finite Taylor expansions of
\begin{align}
 G(t,v)&=(1+t)^{1/2}\exp\left\{-\mathcal R(t)
        +\sum_{k\ge1}\frac{(-1)^kv^{k+1}t^k}{k+1}\right\},\label{ppg:all:eq:Grecipe}\\
 H(t)&=(1+2t)^{-1/2}\exp\left\{\mathcal R(t)
        +\sum_{k\ge1}\frac{(-1)^{k+1}t^k}{k+1}\right\}.\label{ppg:all:eq:Hrecipe}
\end{align}
Only finitely many terms are retained at the requested order. Taylor errors are uniform on the compact interval in $v$ for small $t$; a larger constant covers the finitely many remaining positive $n$. No derivatives in $v$ are used.
\end{proof}
The first values, used later, are
\begin{align}
 G_1(v)&=\frac5{12}-\frac{v^2}2,
 &G_2(v)&=-\frac{47}{288}-\frac{5v^2}{24}+\frac{v^3}3+\frac{v^4}8,
 \label{ppg:all:eq:G12}\\
 p:=\int_0^1G_1&=\frac14,
 &b_2^*:=\int_0^1G_2&=-\frac{179}{1440},\label{ppg:all:eq:b12}\\
 q:=h_1&=-\frac5{12},& f:=h_2&=\frac{217}{288}.
 \label{ppg:all:eq:H12}
\end{align}
The star distinguishes an expansion coefficient $b_2^*$ from the exact integral $b_2$ in \eqref{ppg:all:eq:gzero}.

\section{Exact coefficient models}\label{ppg:all:sec:models}
Write $u=1-z$ and $L=-\log(1-z)$. For positive integers $p,q$ define
\[
 B_{p,q}(z)=u^{p-1}L^q,\qquad b_{p,q}(n)=[z^n]B_{p,q}(z).
\]
These are analytic on $|z|<1$. Their use requires no boundary regularity for the actual sequence.

\begin{lemma}[Coefficient expansions and triangular matching]\label{ppg:all:lem:basis}
Every $b_{p,q}(n)$ has a fixed-order power-log expansion to arbitrary precision, with leading term
\begin{equation}\label{ppg:all:eq:blead}
 b_{p,q}(n)=\frac{(-1)^{p-1}q(p-1)!(\log n)^{q-1}
                 +\text{lower log powers}}{n^p}
       +O\!\left(\frac{(1+\log n)^{q-1}}{n^{p+1}}\right).
\end{equation}
Every later inverse-power coefficient has log degree at most $q-1$. Any specified expansion $\sum_{j=1}^M n^{-j}P_j(\log n)$ with $\deg P_j\le j-1$ therefore has a unique model
\begin{equation}\label{ppg:all:eq:canonical}
 S_M(z)=\sum_{p=1}^M\sum_{q=1}^p s_{p,q}u^{p-1}L^q
\end{equation}
matching it through order $n^{-M}$. The matching is finite and triangular, and preserves uniformity for continuous parameter coefficients on a compact interval.
\end{lemma}
\begin{proof}
For $n\ge1$ one has the exact identity
\begin{equation}\label{ppg:all:eq:Lcoeff}
 [z^n]L^q=\frac{q!}{n}\,
 e_{q-1}\left(1,\frac12,\ldots,\frac1{n-1}\right),
\end{equation}
where $e_j$ is the elementary symmetric polynomial, zero when its index exceeds the number of variables. Indeed, extract the coefficient of $t^q$ in
\[
 [z^n](1-z)^{-t}=\frac t n\prod_{k=1}^{n-1}(1+t/k).
\]
Newton's identities express $e_{q-1}$ using finitely many generalized harmonic numbers. Euler--Maclaurin for $t^{-r}$ gives their expansions to every fixed order with bounded remainders; for $r=1$ the constant is Euler's $\gamma$, and for $r\ge2$ it is $\zeta(r)$. This proves the power-log expansion of \eqref{ppg:all:eq:Lcoeff} and its log-degree bound.

Multiplication by $u^{p-1}$ takes an exact backward difference. Before differencing a remainder, expand \eqref{ppg:all:eq:Lcoeff} sufficiently many orders farther: an $O(n^{-N}\log^{q-1}n)$ error remains of that order after any fixed number of differences, and $N$ is freely chosen. The retained explicit terms can be differenced directly. Their leading coefficient is the nonzero pivot in \eqref{ppg:all:eq:blead}. Solve first in ascending $p$, then in descending log degree. The model error is $O(n^{-M-1}(1+\log n)^{M-1})$. All operations are finite scalar linear operations, giving the asserted uniformity.
\end{proof}

The algebra of models closes exactly:
\begin{equation}\label{ppg:all:eq:product}
 B_{p,q}B_{r,s}=B_{p+r-1,q+s}.
\end{equation}
Polynomials in $u$ will represent moments. Their coefficients have finite support. Two bounds essential below are
\begin{equation}\label{ppg:all:eq:diffbounds}
 |[z^n]S_M|\le\frac{C}{n+1},\qquad
 |[z^n]u^kS_M|\le\frac{C_k}{(n+1)^{k+1}}\quad(k\ge0\text{ fixed}).
\end{equation}
For a term with $p=1$, necessarily $q=1$ and no logarithm occurs. For $p\ge2$, the additional factor $n^{-(p-1)}$ absorbs its fixed log power. Applying \eqref{ppg:all:eq:blead} after multiplication by $u^k$ proves \eqref{ppg:all:eq:diffbounds}. These are estimates on exact models, never on differences of an unspecified remainder.

\section{Signed convolution and endpoint moments}\label{ppg:all:sec:convolution}
We use the commutative Banach algebra $C[0,1]$, with pointwise multiplication and the supremum norm. Scalar real or complex sequences are special cases. Its product satisfies $\norm{fg}\le\norm f\norm g$; reflection $v\mapsto1-v$ preserves the norm.

\begin{lemma}[Uniform signed convolution]\label{ppg:all:lem:conv}
Fix $M\ge1$. Suppose $a_n=s_n+r_n$, where $s_n=[z^n]S_M$ is a model of the form \eqref{ppg:all:eq:canonical}, and
\begin{equation}\label{ppg:all:eq:residue}
 \norm{r_n}\le C\frac{(1+\log(n+2))^M}{(n+1)^{M+1}}.
\end{equation}
Define
\begin{equation}\label{ppg:all:eq:moments}
 \mu_k=\sum_{n\ge0}\binom nk r_n\quad(0\le k<M),
 \qquad J_M(z)=\sum_{k=0}^{M-1}(-1)^k\mu_ku^k.
\end{equation}
These sums converge absolutely in norm. For a second sequence $\widetilde a$ with the same type of decomposition,
\begin{equation}\label{ppg:all:eq:conv}
 (a*\widetilde a)_n=[z^n](S_M+J_M)(\widetilde S_M+\widetilde J_M)
      +O\!\left(\frac{(1+\log n)^{M+1}}{n^{M+1}}\right).
\end{equation}
The polynomial $J_M\widetilde J_M$ has no coefficient once $n>2M-2$.
\end{lemma}
\begin{proof}
The moment of index $k$ has a tail bounded by a constant times
$\sum n^{k-M-1}(1+\log n)^M$, convergent because $k\le M-1$. This also proves uniform convergence and continuity when the coefficients are functions of $v$. Integration in $v$ may therefore be interchanged with these sums.

For $i\le n/2$, apply integer Newton expansion to $f(i)=s_{n-i}$. With $\nabla$ denoting backward difference,
\begin{equation}\label{ppg:all:eq:newtondiscrete}
 s_{n-i}=\sum_{k=0}^{M-1}(-1)^k\binom ik(\nabla^ks)_n+E_{n,i},
 \qquad \norm{E_{n,i}}\le C i^M n^{-M-1}.
\end{equation}
If $i<M$ this is exact. Otherwise the exact remainder is
\[
 \sum_{r=0}^{i-M}\binom{i-1-r}{M-1}\Delta^M f(r).
\]
Its nonnegative weights sum to $\binom iM$, and
$\Delta^Mf(r)=(-1)^M(\nabla^Ms)_{n-r}$. All indices appearing are at least $n-i\ge n/2$, so \eqref{ppg:all:eq:diffbounds} proves the bound.

In the left half of $r*s$, the resulting remainder is bounded by
\begin{equation}\label{ppg:all:eq:borderline}
 Cn^{-M-1}\sum_{i\le n/2}i^M\norm{r_i}
   =O\!\left(n^{-M-1}(1+\log n)^{M+1}\right).
\end{equation}
The moment at index $M$ is generally divergent. Only its truncated size is used here. Replacing each truncated moment of index $k<M$ by its full sum costs at most
\[
 C\sum_{k=0}^{M-1}n^{-k-1}\sum_{i>n/2}i^k\norm{r_i}
 =O\!\left(n^{-M-1}(1+\log n)^M\right).
\]
The other half has the bound
\[
 \sup_{i>n/2}\norm{r_i}\sum_{j<n/2}\norm{s_j}
 =O\!\left(n^{-M-1}(1+\log n)^{M+1}\right),
\]
since the absolute partial sums of $s$ are $O(1+\log n)$. Hence $r*s$ is approximated by $[z^n]J_MS_M$ with the stated error. The same argument applies to either cross product with distinct sequences.

Finally, splitting according to which index is at least $n/2$ gives
\[
 \norm{(r*\widetilde r)_n}
 \le\sup_{i\ge n/2}\norm{r_i}\sum_j\norm{\widetilde r_j}
   +\sup_{j\ge n/2}\norm{\widetilde r_j}\sum_i\norm{r_i}
 =O\!\left(n^{-M-1}(1+\log n)^M\right).
\]
Both residue sums converge. The model-model product is exact, and the remaining moment-moment polynomial has finite support. This proves \eqref{ppg:all:eq:conv} with arbitrary signs and without regularity assumptions on the residual sequence.
\end{proof}

\section{The forcing identity and the induction}\label{ppg:all:sec:induction}
Report 145 proves the exact critical inverse, with its boundary condition already enforced:
\begin{align}
 y_n&=(Th)_n=\frac{h_{n-1}}{n+2}
       -2(n+1)\sum_{k\ge n}\frac{h_k}{(k+1)(k+2)(k+3)},\quad n\ge1,
 \label{ppg:all:eq:T}\\
 h_n&=\sum_{i+j=n}(\beta_{ij}-1)x_ix_j+(y*y)_n.
 \label{ppg:all:eq:hdef}
\end{align}
The tail is absolutely convergent. For context, the generating-function equation is
$Y'-2Y/(1-z)=H$, and boundedness of $y_n$ removes the homogeneous solution proportional to $(1-z)^{-2}$. We do not choose this solution by a formal asymptotic guess.
\begin{writenote}[5 October 2026]
This is Part~\ref{ppg:log:part}'s Lemma~\ref{ppg:log:lem:inverseT}, with $h_n=e_n+z_n$ as there.
\end{writenote}

\subsection{An identity that keeps every residual decaying}
For all $i\ge0$, define
\begin{align}
 a_i(v)&=g_i(v)x_i-1,&\bar a_i&=\int_0^1a_i(v)\,\dd v=b_ix_i-1,
 \label{ppg:all:eq:an}\\
 \delta_i&=(b_i-1)x_i,& U(v)&=d+G_1(v)=\frac16-\frac{v^2}2.
 \label{ppg:all:eq:delta}
\end{align}
The starting rate \eqref{ppg:all:eq:base} and \eqref{ppg:all:eq:gexp} give uniformly
\begin{equation}\label{ppg:all:eq:U}
 a_i(v)=\frac{U(v)}i+O\!\left(\frac{1+\log i}{i^2}\right),\qquad
 \int_0^1U(v)\,\dd v=0.
\end{equation}
Consequently $\bar a_i=O((1+\log i)/i^2)$ and
$\delta_i=p/i+O(i^{-2})$. The fixed sums
\begin{equation}\label{ppg:all:eq:SA}
 S_A=\sum_{i\ge0}\bar a_i,\qquad
 S_D=\delta_0+\sum_{i\ge1}(\delta_i-p/i)
\end{equation}
converge absolutely before the induction begins.
\begin{writenote}[5 October 2026]
$S_D$ is Part~\ref{ppg:log:part}'s $S_a$ of \eqref{ppg:log:eq:residuesums}. The constant $S_A$ is new here; numerically it equals $\frac1{12}$ to 18 digits (Section~\ref{ppg:sec:further}, question~\ref{ppg:fq:sa}), which is not proved. Since $b_i=\sqrt{2\pi(i+1)}\,q_i$ with $q_i=\frac1{i!}\int_i^{i+1}s^ie^{-s}\,\dd s$ (\eqref{ppg:log:eq:bi}) and $A=\sqrt{2\pi}/R$, one has $b_ix_i=q_iW_iR^{i+1}$, so $S_A=\sum_{i\ge0}\bigl(q_iW_iR^{i+1}-1\bigr)$.
\end{writenote}

Expand the exact factorization \eqref{ppg:all:eq:factor} in \eqref{ppg:all:eq:hdef}. The identities
$\bar a_i-y_i=\delta_i$ and
$\sum x_ix_j=n+1+2\sum y_i+(y*y)_n$ cancel the latter quadratic convolution. Thus, for every $n\ge1$,
\begin{align}
 h_n={}&(H_n-1)(n+1)+2\sum_{i=0}^n\delta_i
       +2(H_n-1)\sum_{i=0}^n\bar a_i\notag\\
      &+H_n\int_0^1\sum_{i+j=n}a_i(v)a_j(1-v)\,\dd v.
 \label{ppg:all:eq:forcingidentity}
\end{align}
This is an exact identity including zero endpoints. A constant sequence has not been left convolved with a poorly controlled residual.

\subsection{One order of forcing}
Assume \eqref{ppg:all:eq:all} at order $M$, including its displayed error. Multiplication by the uniform kernel expansion yields an expansion for $a_n(v)$ through order $M$ with degree at most $j-1$ at order $j$ and uniform error $O(n^{-M-1}(1+\log n)^M)$. Form its canonical model $S_M(v;z)$ and the moments \eqref{ppg:all:eq:moments}. Lemma~\ref{ppg:all:lem:conv}, with reflection in $v$, gives the last convolution in \eqref{ppg:all:eq:forcingidentity} from
\begin{equation}\label{ppg:all:eq:integratedproduct}
 \int_0^1[S_M(v;z)+J_M(v;z)]
          [S_M(1-v;z)+J_M(1-v;z)]\,\dd v,
\end{equation}
with error $O(n^{-M-1}(1+\log n)^{M+1})$.

Here is the degree and truncation bookkeeping. A product of two model terms has the form $u^aL^q$ with $q\le a+2$; its first coefficient order is $n^{-a-1}$ with log degree at most $a+1$. A cross product with $J_M$ satisfies $q\le a+1$. Retain $a<M$. Every later inverse-power term from a retained monomial has the same or lower log degree and thus satisfies the required degree bound at its later order. When $a=M$, the largest discarded term is $O(n^{-M-1}(1+\log n)^{M+1})$. When $a>M$, the extra inverse power absorbs the extra fixed logarithms. The polynomial product has finite support. Multiplication by a sufficiently long $H_n$ expansion preserves this precision; independently, the convolution is $O((1+\log n)/n)$ by an elementary half split of $a_n=O(1/n)$.

For the remaining terms, $(b_i-1)=O(i^{-1})$ means the known $x_i$ expansion through order $M$ determines $\delta_i$ through order $M+1$ with error $O(i^{-M-2}(1+\log i)^M)$. Hence
\begin{align}
 \sum_{i=0}^n\delta_i
  &=p\HH_n+S_D-\sum_{i>n}(\delta_i-p/i),\label{ppg:all:eq:deltasum}\\
 \sum_{i=0}^n\bar a_i&=S_A-\sum_{i>n}\bar a_i,\label{ppg:all:eq:abarSum}
\end{align}
where $\HH_n=\sum_{i=1}^n1/i$. The first formula has an expansion through $n^{-M}$ with error $O(n^{-M-1}(1+\log n)^M)$. The second only needs to be taken through $n^{-(M-1)}$: its error $O(n^{-M}(1+\log n)^M)$ acquires another inverse power from $H_n-1$.

These assertions use ordinary absolutely convergent model tails. For $F(t)=t^{-p}(\log t)^q$, $p\ge2$, Euler--Maclaurin with the appropriate endpoint convention gives
\begin{align}
 \sum_{k>n}F(k)={}&\int_n^\infty F(t)\,\dd t-\frac{F(n)}2
  -\sum_{r=1}^J\frac{B_{2r}}{(2r)!}F^{(2r-1)}(n)+\mathcal E_J,
 \label{ppg:all:eq:EM}\\
 |\mathcal E_J|&\le C_J\int_n^\infty |F^{(2J)}(t)|\,\dd t.\notag
\end{align}
Choose $J$ sufficiently large. For example, the tail of $k^{-2}$ begins
$n^{-1}-(2n^2)^{-1}+(6n^3)^{-1}$. Apply \eqref{ppg:all:eq:EM} only to explicit model terms; bound the residual tails absolutely, without differentiating them. At $M=1$, \eqref{ppg:all:eq:abarSum} simply gives $S_A+O((1+\log n)/n)$, which suffices.

For the factor $(H_n-1)(n+1)$, take the kernel expansion through order $M+1$, so multiplication by $n+1$ leaves an error $O(n^{-M-1})$. Combining all terms proves
\begin{equation}\label{ppg:all:eq:hexp}
 h_n=\frac12\log n+B+\sum_{k=1}^M\frac{Q_k(\log n)}{n^k}
       +O\!\left(\frac{(1+\log n)^{M+1}}{n^{M+1}}\right),
 \qquad \deg Q_k\le k.
\end{equation}
The leading logarithm is $2p\log n$. The non-convolution tails have degree at most $k-1$ at order $k$, while the model-model products may have degree $k$.

\subsection{The exact inverse improves a full power}
For integers $a\ge-1$, $q\ge0$, the coefficientwise inverse \eqref{ppg:all:eq:T} applied to $F(z)=u^aL^q$ is represented, on positive coefficients, by
\begin{equation}\label{ppg:all:eq:Tmodel}
 T[F](z)=-u^{a+1}\sum_{j=0}^q
       \frac{q!}{(q-j)!}\frac{L^{q-j}}{(a+3)^{j+1}}.
\end{equation}
Indeed the critical generating function is
\[
 -u^{-2}\int_0^u t^2F(1-t)\,\dd t,
\]
and repeated integration by parts proves \eqref{ppg:all:eq:Tmodel}. The integral converges at zero. To check its coefficient interpretation directly, for the forcing $F(z)=z^m$ the coefficient at positive $n$ is
\[
 \begin{cases}
  -2(n+1)/((m+1)(m+2)(m+3)),&1\le n\le m,\\
  1/(m+3),&n=m+1,\\
  0,&n>m+1.
 \end{cases}
\]
These are exactly \eqref{ppg:all:eq:T}. Absolute weighted summability extends this verification to the logarithmic models. A polynomial forcing can leave a polynomial output of one higher degree; it is ignored only for sufficiently large $n$.

For arbitrary signed residuals, the following estimate requires no smoothness:
\begin{equation}\label{ppg:all:eq:Timprovement}
 v_k=O\!\left(\frac{(1+\log(k+2))^D}{(k+1)^s}\right)
 \quad\Longrightarrow\quad
 (Tv)_n=O\!\left(\frac{(1+\log n)^D}{n^{s+1}}\right)
 \quad(s,D\ge0\text{ fixed}).
\end{equation}
The local term has that size. The absolute tail is at most a constant times
$n\sum_{k\ge n}(1+\log(k+2))^D/(k+1)^{s+3}$, which has the same bound by integration. Finite initial changes affect only finitely many output indices.

Match \eqref{ppg:all:eq:hexp} to a finite coefficient model using $u^{-1}L$, $u^{-1}$ and $u^{k-1}L^q$ for $1\le k\le M$, $1\le q\le k+1$. The first term has coefficients $\HH_n$; subsequent terms are handled by the triangular matching of Lemma~\ref{ppg:all:lem:basis}. The model error is within the error already displayed in \eqref{ppg:all:eq:hexp}. Under \eqref{ppg:all:eq:Tmodel}, the $u^{-1}L$ term maps to $-L/2-1/4$, and the $u^{-1}$ term maps to the constant polynomial $-1/2$. A term $u^{k-1}L^q$ maps to terms $u^kL^{q-j}$, whose first nonpolynomial order is $n^{-k-1}$ and log degree at most $q-1\le k$. Their later coefficients keep that degree bound. The leading forcing $(1/2)\log n$ therefore produces $P_1=-1/4$.

Taking each exact model expansion sufficiently far, and applying \eqref{ppg:all:eq:Timprovement} to the forcing residual, gives
\[
 (Th)_n=\sum_{j=1}^{M+1}\frac{P_j^{\rm new}(\log n)}{n^j}
      +O\!\left(\frac{(1+\log n)^{M+1}}{n^{M+2}}\right),
 \qquad\deg P_j^{\rm new}\le j-1.
\]
Since $y=Th$, this advances the induction. Subtracting the old and new expansions shows agreement at each earlier order: the least differing inverse power, with its highest nonzero logarithm, could not fit a remainder one full inverse power smaller. Starting from \eqref{ppg:all:eq:base} proves \eqref{ppg:all:eq:all}. Equation \eqref{ppg:all:eq:oldsecond} identifies $P_2$ by the same uniqueness argument. Multiplication by the finite square-root expansion proves \eqref{ppg:all:eq:allW}.

\begin{remark}[Why the moments are not circular]
At $M=1$ only a zeroth residual moment is used, whose summands are $O((1+\log n)/n^2)$. This first step proves the stronger second-order remainder $O((1+\log n)^2/n^3)$. Only then, at $M=2$, is a first residual moment introduced; its summands are $O((1+\log n)^2/n^2)$. A mere $o(n^{-2})$ residual would not justify that moment.
\end{remark}

\section{A finite scalar recursion}\label{ppg:all:sec:scalar}
The function-valued notation proves all required estimates uniformly, but a calculation at a fixed order does not require an arbitrary function as an independent input. Every coefficient $s_{p,q}(v)$ of the canonical $S_M(v;z)$ is a polynomial in $v$: the kernel coefficients are polynomials, and canonical matching uses only finite scalar linear operations.

\begin{proposition}[Scalar connection constants]\label{ppg:all:prop:scalar}
At the step from $M$ to $M+1$, the convolution contribution through order $n^{-M}$ is determined by finite polynomial algebra and at most
\[
 \sum_{p=1}^M p(M-p+1)=\frac{M(M+1)(M+2)}6
\]
indexed scalar sums, in addition to the fixed sums $S_A,S_D$. Specifically, put
\begin{align}
 r_{M,n}(v)&=g_n(v)x_n-1-[z^n]S_M(v;z),\label{ppg:all:eq:rMn}\\
 m_{k,p,q}&=\sum_{n\ge0}\binom nk\int_0^1s_{p,q}(v)r_{M,n}(1-v)\,\dd v,
 \quad 1\le p\le M,\quad1\le q\le p,\quad0\le k\le M-p.
 \label{ppg:all:eq:scalarseries}
\end{align}
Every series in \eqref{ppg:all:eq:scalarseries} is absolutely convergent.
\end{proposition}
\begin{proof}
Uniform norm convergence allows the sum and integral to be interchanged, giving
$m_{k,p,q}=\int_0^1s_{p,q}(v)\mu_k(1-v)\,\dd v$. Reflecting one cross term in \eqref{ppg:all:eq:integratedproduct} shows that the relevant nonpolynomial expression is
\begin{align}
 &\int_0^1 S_M(v;z)S_M(1-v;z)\,\dd v\notag\\
 &\qquad+2\sum_{p=1}^M\sum_{q=1}^p\sum_{k=0}^{M-p}
       (-1)^km_{k,p,q}u^{k+p-1}L^q.\label{ppg:all:eq:scalarproduct}
\end{align}
The retention condition is exactly $k+p-1<M$. Its summands in \eqref{ppg:all:eq:scalarseries} are $O((1+\log n)^M/n^{M+1-k})$ and the exponent is at least 2. The pure moment product is a polynomial of degree at most $2M-2$; even after applying $T$ it remains finitely supported. No unknown algebraic tail is lost by discarding it. Counting the indexed moments gives the bound.
\end{proof}

Here is a complete order-advancement procedure.
\begin{enumerate}
\item From the already proved $P_1,\ldots,P_M$ and \eqref{ppg:all:eq:Grecipe}--\eqref{ppg:all:eq:Hrecipe}, form the asymptotic coefficients of $a_n(v)$ through order $M$.
\item Use \eqref{ppg:all:eq:Lcoeff} and triangular matching to construct $S_M(v;z)$.
\item Define the finitely many scalar sums \eqref{ppg:all:eq:scalarseries}. Their convergence follows from the already proved order-$M$ remainder, before $P_{M+1}$ is known.
\item Form \eqref{ppg:all:eq:scalarproduct}; retain the coefficient expansion through $n^{-M}$ with the explicit discarded-term bounds of Section~\ref{ppg:all:sec:induction}.
\item Combine it with \eqref{ppg:all:eq:deltasum}--\eqref{ppg:all:eq:EM}, the finite $H_n$ expansion, and \eqref{ppg:all:eq:forcingidentity}, obtaining \eqref{ppg:all:eq:hexp}.
\item Match to the finite forcing basis, apply \eqref{ppg:all:eq:Tmodel}, and keep the new polynomial $P_{M+1}$.
\end{enumerate}
The exact recurrence and exact $R$ already define $x_n$. A series involving the full exact sequence is a definition of a connection constant, not a circular assumption about the next asymptotic coefficient. ``Finite'' here means finitely many algebraic operations and finitely many specified convergent series at each fixed order. It does not mean that an infinite sum is evaluated exactly by finitely many rational operations. The count is an upper bound on indexed constants, not a minimality or independence assertion.

\section{The explicit third order}\label{ppg:all:sec:third}
The first induction step strengthens \eqref{ppg:all:eq:oldsecond} to
\begin{equation}\label{ppg:all:eq:strongsecond}
 x_n=1+\frac d n+\frac{c\log n+C}{n^2}
       +O\!\left(\frac{(1+\log n)^2}{n^3}\right).
\end{equation}
All first moments used below are justified by this strengthened estimate.

\begin{theorem}[Third order and squared logarithm cancellation]\label{ppg:all:thm:third}
With $S_A$ defined by \eqref{ppg:all:eq:SA},
\begin{align}
 x_n={}&1-\frac1{4n}+\frac{c\log n+C}{n^2}
  +\frac{(323/22680+cS_A)\log n+K_x}{n^3}
  +O\!\left(\frac{(1+\log n)^3}{n^4}\right),\label{ppg:all:eq:thirdx}\\
 \frac{W_n}{A\sqrt n\,r_*^n}
 ={}&1+\frac1{4n}+\frac{c\log n+B_2}{n^2}
  +\frac{(22/2835+cS_A)\log n+B_3}{n^3}
  +O\!\left(\frac{(1+\log n)^3}{n^4}\right),\label{ppg:all:eq:thirdW}\\
 B_3={}&K_x+\frac C2+\frac3{32}.\label{ppg:all:eq:B3}
\end{align}
The constants $K_x,B_3$ are specified by ordinary convergent series and polynomial integrals below. The coefficient of $(\log n)^2/n^3$ is zero in both normalizations.
\end{theorem}
\begin{writenote}[5 October 2026]
A numerical check by the write (question~\ref{ppg:fq:sa} of Section~\ref{ppg:sec:further}; uncertified) fits $\log x_n$ for $300\le n\le4000$ and finds the coefficient of $(\log n)^2/n^3$ to be about $2\times10^{-12}$ and that of $\log n/n^3$ to be $0.009920635$ to within about $10^{-10}$. With $S_A=\frac1{12}$ the theorem predicts $0$ and $\frac{323}{22680}+\frac c{12}+\frac c4=\frac5{504}=0.0099206349\ldots$, where $\frac c4=-dc$ comes from taking the logarithm. If $S_A=\frac1{12}$, the logarithmic coefficients in \eqref{ppg:all:eq:thirdx} and \eqref{ppg:all:eq:thirdW} are $\frac{199}{15120}$ and $\frac{101}{15120}$.
\end{writenote}

\subsection{The model and convergent constants}
Further kernel coefficients obtained from \eqref{ppg:all:eq:Grecipe}--\eqref{ppg:all:eq:Hrecipe} are
\begin{equation}\label{ppg:all:eq:thirdkernel}
 b_3^*:=\int_0^1G_3(v)\,\dd v=\frac{25009}{362880},
 \qquad h_3=-\frac{71329}{51840}.
\end{equation}
These use $r_n=1/(12n)-1/(360n^3)+O(n^{-4})$. Set
\begin{equation}\label{ppg:all:eq:VF}
 V(v)=C+dG_1(v)+G_2(v),\qquad
 F(v)=c(\gamma-1)-V(v),\qquad K=-c/2.
\end{equation}
By \eqref{ppg:all:eq:strongsecond},
$a_n(v)=U(v)/n+[c\log n+V(v)]/n^2+O((1+\log n)^2/n^3)$ uniformly. Its exact canonical model is
\begin{equation}\label{ppg:all:eq:S2}
 S_2(v;z)=U(v)L+u\{KL^2+F(v)L\}.
\end{equation}
Indeed $[z^n]uL=-n^{-2}+O(n^{-3})$ and
$[z^n]uL^2=(-2\log n+2-2\gamma)n^{-2}+O((1+\log n)n^{-3})$.
Define, including all initial indices,
\begin{equation}\label{ppg:all:eq:muThird}
 r_n(v)=a_n(v)-[z^n]S_2(v;z),\qquad
 \mu_0(v)=\sum_{n\ge0}r_n(v),\qquad
 \mu_1(v)=\sum_{n\ge0}nr_n(v).
\end{equation}
Both converge absolutely uniformly. The coefficients of $uL$ and $uL^2$ are absolutely summable with sum zero, since their functions tend to zero as $z\uparrow1$. Thus
\begin{equation}\label{ppg:all:eq:muZeroSimple}
 \mu_0(v)=a_0(v)+\sum_{n\ge1}\{a_n(v)-U(v)/n\},\qquad
 \int_0^1\mu_0(v)\,\dd v=S_A.
\end{equation}
There is no analogous permission to omit the full model from the first moment.

Use the symmetric bilinear form
$\ip FG=\int_0^1F(v)G(1-v)\,\dd v$, and define
\begin{align}
 S_0&=\ip UU=-\frac7{360},& A_0&=2\ip U{\mu_0},\notag\\
 A_2&=2\ip UF+2KS_A,& A_1&=2\{\ip F{\mu_0}-\ip U{\mu_1}\}.
 \label{ppg:all:eq:Aconstants}
\end{align}
Every integral here is of a uniformly convergent series against a polynomial.

\subsection{The cancellation and the forcing coefficients}
At precision $O((1+\log n)^3/n^3)$, Lemma~\ref{ppg:all:lem:conv} says the integrated convolution is the coefficient of
\begin{equation}\label{ppg:all:eq:thirdconvmodel}
 S_0L^2+A_0L+u\{A_2L^2+A_1L\}.
\end{equation}
The candidate term $uL^3$ has coefficient
$K\int_0^1\{U(v)+U(1-v)\}\,\dd v=0$. This exact zero mean, rather than numerical fitting, removes the squared logarithm in the final third order.
Since $[z^n]L^2=2(\log n+\gamma)/n-1/n^2+O(n^{-3})$, the convolution equals
\begin{equation}\label{ppg:all:eq:thirdconv}
 \frac{2S_0\log n+c_{10}}n+\frac{-2A_2\log n+c_{20}}{n^2}
       +O\!\left(\frac{(1+\log n)^3}{n^3}\right),
\end{equation}
where
\begin{equation}\label{ppg:all:eq:c1020}
 c_{10}=2\gamma S_0+A_0,\qquad
 c_{20}=-S_0+(2-2\gamma)A_2-A_1.
\end{equation}
Put
\begin{equation}\label{ppg:all:eq:a23}
 a_2=b_2^*+pd=-\frac{269}{1440},\quad
 a_3=b_3^*+b_2^*d+pC,\quad\bar V=\int_0^1V=C+a_2.
\end{equation}
The two partial sums, computed from \eqref{ppg:all:eq:strongsecond} by convergent tails, are
\begin{align}
 \sum_{i=0}^n\delta_i={}&p\log n+p\gamma+S_D+\frac{p/2-a_2}n\notag\\
 &+\frac{-p/12+a_2/2-(pc/2)\log n-pc/4-a_3/2}{n^2}
       +O\!\left(\frac{(1+\log n)^2}{n^3}\right),\label{ppg:all:eq:thirdDsum}\\
 \sum_{i=0}^n\bar a_i={}&S_A-\frac{c(\log n+1)+\bar V}n
       +O\!\left(\frac{(1+\log n)^2}{n^2}\right).\label{ppg:all:eq:thirdAsum}
\end{align}
For example, the tails of $(\log i)/i^3$ and $(\log i)/i^2$ begin respectively
$(\log n/2+1/4)/n^2$ and $(\log n+1)/n$. Substitution into \eqref{ppg:all:eq:forcingidentity} gives
\begin{equation}\label{ppg:all:eq:thirdh}
 h_n=\tfrac12\log n+B+\frac{L_1\log n+C_1}n
         +\frac{L_2\log n+D}{n^2}
         +O\!\left(\frac{(1+\log n)^3}{n^3}\right),
\end{equation}
with the full constant definitions
\begin{align}
 L_1&=2S_0=-\frac7{180},\notag\\
 C_1&=q+f+p-2a_2+2qS_A+c_{10},\notag\\
 L_2&=-(p+2q)c-2A_2+2qS_0,\label{ppg:all:eq:forcingThirdConstants}\\
 D&=f+h_3-p/6+a_2-pc/2-a_3
        +2fS_A-2q(c+\bar V)+c_{20}+qc_{10}.\notag
\end{align}
Direct polynomial integration gives
\begin{equation}\label{ppg:all:eq:L2eval}
 \ip UV=\frac{173}{15120},\qquad
 A_2=-2\ip UV-cS_A,\qquad
 L_2=\frac{617}{11340}+2cS_A.
\end{equation}
There is no dependence on the unknown $C$ in $\ip UV$, because $\int U=0$.

\subsection{Applying the inverse}
Rule \eqref{ppg:all:eq:Tmodel} or the tail formula \eqref{ppg:all:eq:T} now yields
\begin{equation}\label{ppg:all:eq:Kx}
 K_x=C-\frac{L_1}4-\frac{L_2}8+\frac D2+\frac1{48},
 \qquad P_3(t)=(c+L_2/2)t+K_x.
\end{equation}
To make the index algebra explicit, if the leading forcing in \eqref{ppg:all:eq:thirdh} is $a_0\log n$ in place of $(1/2)\log n$, then
\begin{align}
 (Th)_n={}&-\frac{a_0}{2n}
 +\frac{C_1/3+(L_1/3)\log n-2L_1/9-a_0/6}{n^2}\notag\\
 &+\frac{C_1/3+D/2+(L_1/3+L_2/2)\log n
                   -17L_1/36-L_2/8-a_0/8}{n^3}
 +O\!\left(\frac{(1+\log n)^3}{n^4}\right).\label{ppg:all:eq:inverse3}
\end{align}
Taking $a_0=1/2$ gives $c=L_1/3$ and
$C=C_1/3-2L_1/9-1/12$, consistent with \eqref{ppg:all:eq:oldsecond}. It reduces the cubic constant to \eqref{ppg:all:eq:Kx}. Combining \eqref{ppg:all:eq:L2eval} with \eqref{ppg:all:eq:Kx} proves \eqref{ppg:all:eq:thirdx}. Finally, multiply by
$\sqrt{1+1/n}=1+1/(2n)-1/(8n^2)+1/(16n^3)+O(n^{-4})$
to obtain \eqref{ppg:all:eq:thirdW} and \eqref{ppg:all:eq:B3}. This proves Theorem~\ref{ppg:all:thm:third}.

\section{The first logarithm in the consecutive ratio}\label{ppg:all:sec:ratio}
\begin{corollary}\label{ppg:all:cor:ratio}
For $\rho_n=W_{n+1}/W_n$,
\begin{equation}\label{ppg:all:eq:ratio3}
 \frac{\rho_n}{r_*}
 =1+\frac1{2n}-\frac3{8n^2}
   +\frac{(7/270)\log n+32/135-2B_2}{n^3}
   +O\!\left(\frac{(1+\log n)^3}{n^4}\right).
\end{equation}
More generally, the ratio has a consistent expansion at every fixed inverse-power order, obtained from sufficiently many forward terms.
\end{corollary}
\begin{proof}
Let $F_n=W_n/(A\sqrt n\,r_*^n)$ and $b=1/4$. Taking a finite logarithm of \eqref{ppg:all:eq:thirdW} gives
\[
 \log F_n=\frac b n+\frac{c\log n+k}{n^2}
       +\frac{\text{an explicit linear log polynomial}}{n^3}
       +O\!\left(\frac{(1+\log n)^3}{n^4}\right),\qquad
 k=B_2-b^2/2.
\]
The consecutive difference of the explicit cubic term is $O((1+\log n)/n^4)$. The consecutive difference of the arbitrary remainder is bounded only at its original order $O((1+\log n)^3/n^4)$, which is sufficient. Expanding the two explicit lower-order terms and $\frac12\log(1+1/n)$ therefore gives
\[
 \log(\rho_n/r_*)
 =\frac1{2n}-\frac1{2n^2}
   +\frac{1/6+b+c-2c\log n-2k}{n^3}
   +O\!\left(\frac{(1+\log n)^3}{n^4}\right).
\]
A finite exponential yields \eqref{ppg:all:eq:ratio3}; its cubic constant is
$-1/16+b+c+b^2-2B_2=32/135-2B_2$.
For any requested finite order, first use a sufficiently long instance of \eqref{ppg:all:eq:allW}, and then take differences only of its explicit smooth terms. The untouched residual is bounded at its original order. The same finite algebra gives consistency by uniqueness. No improved bound for a difference of an unspecified residual is assumed.
\end{proof}
The logarithm in \eqref{ppg:all:eq:ratio3} resolves the cubic ratio coefficient left undetermined in Report 145. It does not prove all-index ratio monotonicity or an effective onset of the already established eventual monotonicity.
\begin{writenote}[5 October 2026]
This answers the first half of item~3 of Part~\ref{ppg:log:part}'s Section~\ref{ppg:log:sec:open} and implies Theorem~\ref{ppg:log:thm:ratio} (see the note at the end of Section~\ref{ppg:log:sec:ratio}). In float64, $n^3(\rho_n/r_*-1-\frac1{2n}+\frac3{8n^2})$ follows $\frac7{270}\log n+\frac{32}{135}-2B_2$ to within about $0.003$ for $250\le n\le4000$ (placement dossier).
\end{writenote}

\begin{corollary}[Logarithm-aware dyadic extrapolation]\label{ppg:all:cor:extrapolation}
The following asymptotic acceleration is valid:
\begin{equation}\label{ppg:all:eq:fivepoint}
 \frac{\rho_n-22\rho_{2n}+168\rho_{4n}-512\rho_{8n}+512\rho_{16n}}{147}
 =r_*+O\!\left(\frac{(1+\log n)^3}{n^4}\right).
\end{equation}
More generally, write $E f(n)=f(2n)$. For fixed $J\ge1$ the finite operator
\begin{equation}\label{ppg:all:eq:generalextrap}
 \mathcal Q_J(E)=\prod_{j=1}^J
     \left(\frac{E-2^{-j}}{1-2^{-j}}\right)^{\max(1,j-1)}
\end{equation}
satisfies
$\mathcal Q_J(E)\rho_n=r_*+O((1+\log n)^J/n^{J+1})$.
\end{corollary}
\begin{proof}
Using the forward expansion at order $J$ in the preceding proof gives ratio coefficients of degree zero at order one and degree at most $j-2$ at order $j\ge2$, with remainder $O((1+\log n)^J/n^{J+1})$. The degree drops because the consecutive difference raises the inverse-power order of each explicit forward term by at least one; nonlinear products do not increase this bound. The arbitrary forward residual retains its original bound.

On $n^{-j}P(\log n)$, the factor $E-2^{-j}$ replaces $P(t)$ by
$2^{-j}\{P(t+\log2)-P(t)\}$, decreasing a nonconstant polynomial's degree by one and annihilating a constant. Hence the multiplicities in \eqref{ppg:all:eq:generalextrap} kill all retained terms, and the denominators preserve constant sequences. A finite number of dyadic shifts preserves the displayed error. Expanding the product at $J=3$ gives \eqref{ppg:all:eq:fivepoint}; the repeated factor at $2^{-3}$ is essential for the cubic logarithm.
\end{proof}
The corollary is an asymptotic cancellation statement. It supplies neither an effective error constant nor a certificate for decimal digits obtained by numerical extrapolation.

\begin{remark}[\wtag\ Serra's Open Problem~2, 5 October 2026]\label{ppg:all:rem:op2}
Serra's Open Problem~2 (printed page~68) asks for a proof that, for every $p\ge1$, $r(n)=W(n+1)/W(n)$ satisfies his~(199),
\[
 r(n)=r_\infty+\frac{a_1}n+\frac{a_2}{n^2}+\cdots+\frac{a_p}{n^p}+O\Bigl(\frac1{n^{p+1}}\Bigr),
\]
with $a_1\approx0.288$; his Remark~8.2 says that the Richardson estimates of his Section~6.4 have a certified rate only if~(199) holds. For $p=1$, (199) holds with $a_1=r_*/2\in(0.263,0.316)$ (Corollary~\ref{ppg:log:cor:monotonicity}; the interval comes from \eqref{ppg:lead:eq:numeric}). For every $p\ge2$, (199) is false.
\end{remark}
\begin{proof}
If (199) holds for some $p\ge2$, it holds for $p=2$, because $a_3n^{-3}+\cdots+a_pn^{-p}+O(n^{-p-1})=O(n^{-3})$. By \eqref{ppg:all:eq:ratio3},
\[
 \rho_n=r_*+\frac{r_*}{2n}-\frac{3r_*}{8n^2}+\frac{r_*\bigl(\frac7{270}\log n+\frac{32}{135}-2B_2\bigr)}{n^3}+O\Bigl(\frac{(1+\log n)^3}{n^4}\Bigr).
\]
Subtract~(199) with $p=2$. Multiplying the difference by $n$ and letting $n\to\infty$ gives $a_1=r_*/2$; then multiplying by $n^2$ gives $a_2=-3r_*/8$. What remains is $\frac{7r_*}{270}\frac{\log n}{n^3}=O(n^{-3})$, impossible since $r_*>0$.
\end{proof}
Serra's Section~6.4 assumes the all-orders form~(151); Theorem~\ref{ppg:log:thm:ratio} already excludes it through order $n^{-3}$. His Proposition~6.21, which derives the rate of the Richardson table from~(199), is a conditional statement whose hypothesis fails for $p\ge2$; nothing here decides whether his numerical digits are correct. Corollary~\ref{ppg:all:cor:extrapolation} is a dyadic scheme whose cancellation is proved with the logarithms present.

\section{All fixed orders of the threshold inverse}\label{ppg:all:sec:inverse}
Let $\lambda=\log(1/r_*)>0$ and define
\begin{equation}\label{ppg:all:eq:N}
 N(\varepsilon)=\min\{n\ge0:W_n\le\varepsilon\},\qquad
 L_\varepsilon=\log(1/\varepsilon).
\end{equation}
Every admissible $W_n$ is strictly positive and $A>0$. These hypotheses matter for a first crossing: asymptotics alone do not exclude an initial zero. The set in \eqref{ppg:all:eq:N} is nonempty for $\varepsilon>0$ because $W_n\to0$.

The logarithm of \eqref{ppg:all:eq:allW}, at any fixed $M\ge1$, gives
\begin{equation}\label{ppg:all:eq:logW}
 \log W_n=\log A+\tfrac12\log n-\lambda n
       +\sum_{j=1}^M\frac{\mathcal L_j(\log n)}{n^j}+d_{M,n},
 \quad d_{M,n}=O_M\!\left(\frac{(1+\log n)^M}{n^{M+1}}\right),
\end{equation}
with $\deg\mathcal L_j\le j-1$. This follows by finite Taylor expansion of $\log(1+w)$: a product of $r$ forward terms of total order $j$ has degree at most $j-r$. In particular,
\begin{equation}\label{ppg:all:eq:logfirst}
 \mathcal L_1=b=\frac14,\qquad
 \mathcal L_2(t)=ct+k,\qquad k=B_2-\frac1{32}.
\end{equation}
Define the explicit smooth models
\begin{equation}\label{ppg:all:eq:fM}
 f_0(t)=\lambda t-\tfrac12\log t-\log A,
 \qquad f_M(t)=f_0(t)-\sum_{j=1}^M t^{-j}\mathcal L_j(\log t).
\end{equation}
The exact leading large root is
\begin{equation}\label{ppg:all:eq:x0}
 x=x_0(\varepsilon)
 =-\frac1{2\lambda}W_{-1}\!\left(-\frac{2\lambda\varepsilon^2}{A^2}\right),
 \qquad f_0(x)=L_\varepsilon,
\end{equation}
where $W_{-1}$ denotes the lower real Lambert-W branch. For sufficiently small $\varepsilon$ its argument belongs to $(-1/e,0)$ and $x\to\infty$, with $x\sim L_\varepsilon/\lambda$. The principal branch would give the wrong, small root.

\begin{theorem}[Ceiling-safe all-order inversion]\label{ppg:all:thm:threshold}
Fix $M\ge1$ and put
\[
 E_M(x)=\frac{(1+\log x)^M}{x^{M+1}}.
\]
There are uniquely determined polynomials $\mathcal R_j$ of degree at most $j-1$ such that the finite center
\begin{equation}\label{ppg:all:eq:formalcenter}
 z_M=x+\sum_{j=1}^M x^{-j}\mathcal R_j(\log x)
\end{equation}
satisfies, for some finite constant $C_M>0$ and all sufficiently small $\varepsilon>0$,
\begin{equation}\label{ppg:all:eq:ceiling}
 \left\lceil z_M-C_ME_M(x)\right\rceil
 \le N(\varepsilon)\le
 \left\lceil z_M+C_ME_M(x)\right\rceil.
\end{equation}
The same bracket holds, after adjusting $C_M$, for the Newton center $z^{(k)}$ defined by
\begin{equation}\label{ppg:all:eq:newton}
 z^{(0)}=x,\qquad
 z^{(k+1)}=z^{(k)}-
     \frac{f_M(z^{(k)})-L_\varepsilon}{f'_M(z^{(k)})},
\end{equation}
whenever $k$ is fixed and $3\cdot2^k-2\ge M+1$.
The constants may depend on $M$ and the fixed forward data. They are not effective error certificates in this article.
\end{theorem}

\subsection{The finite inverse recursion}
Write $\ell=\log x$, $\delta=\sum_{j=1}^M x^{-j}\mathcal R_j(\ell)$, and $w=\delta/x$. The exact expression to expand is
\begin{align}
 f_M(x+\delta)-f_0(x)
 =\lambda\delta-\tfrac12\log(1+w)
  -\sum_{j=1}^M x^{-j}(1+w)^{-j}
       \mathcal L_j\bigl(\ell+\log(1+w)\bigr).
 \label{ppg:all:eq:inversesub}
\end{align}
At inverse order $m$, the new unknown first appears as $\lambda\mathcal R_m$. Every other appearance is through $w$ and thus at a strictly higher inverse order. Since $\lambda>0$, consecutive coefficient cancellation determines the polynomials uniquely.

The degree bound can be checked without assuming a formal infinite solution. Suppose the earlier $\mathcal R_i$ have degree at most $i-1$. A product of $r$ terms of $w$ with indices $i_1,\ldots,i_r$ has inverse order $\sum(i_a+1)$ and degree at most $\sum(i_a-1)$. When multiplied by the $j$th forward term, the total order is $m=j+\sum(i_a+1)$ and its degree is at most $m-1-2r\le m-1$. Derivatives of $\mathcal L_j$ only reduce that degree. Terms from $\log(1+w)$ alone have degree at most $m-2r\le m-2$. This proves $\deg\mathcal R_m\le m-1$.

The coefficients of $\mathcal R_j$ lie in the polynomial ring over the rationals generated by $\lambda^{-1}$ and the forward coefficients through the same order $j$. There is no additional dependence on $A$ after choosing the exact leading root as coordinate. For example, if
$\mathcal L_3(\ell)=a_2\ell^2+d_1\ell+e_0$, then
\begin{align}
 \mathcal R_1(\ell)&=b/\lambda,\notag\\
 \mathcal R_2(\ell)&=(c\ell+k)/\lambda+b/(2\lambda^2),\label{ppg:all:eq:R123}\\
 \mathcal R_3(\ell)&=(a_2\ell^2+d_1\ell+e_0)/\lambda
 +(c\ell+k)/(2\lambda^2)-b^2/\lambda^2+b/(4\lambda^3).\notag
\end{align}
For the placement sequence the third logarithmic polynomial is linear, so $a_2=0$. These formulas concern explicit smooth models, not a derivative of the sequence remainder.

For each fixed $M$, $\delta=O(1/x)$ and $w=O(x^{-2})$. Finite Taylor expansions of the smooth factors in \eqref{ppg:all:eq:inversesub}, with bounded errors for $|w|\le1/2$, justify every cancellation and yield
\begin{equation}\label{ppg:all:eq:formalresidual}
 f_M(z_M)-L_\varepsilon=O(E_M(x)).
\end{equation}
At the first uncanceled inverse order the log degree is at most $M$; higher fixed log powers are absorbed by their additional inverse powers after a sufficiently deep finite Taylor expansion.

\subsection{Newton errors for the smooth model}
At large $t$,
\[
 f'_M(t)=\lambda-\frac1{2t}+O(t^{-2}),\qquad
 f''_M(t)=\frac1{2t^2}+O(t^{-3}).
\]
Thus $f'_M\ge\lambda/2$ for $t\ge x/2$ once $x$ is large. Since $f_M(x)-L_\varepsilon=O(x^{-1})$, the values at $x\pm B/x$, with sufficiently large fixed $B$, bracket a unique root $\alpha_M=x+O(x^{-1})$. Taylor's theorem on this neighborhood proves
\[
 |z^{(k+1)}-\alpha_M|\le Cx^{-2}|z^{(k)}-\alpha_M|^2.
\]
For each fixed $k$ the iterates remain there. The exponents satisfy
$p_0=1$, $p_{k+1}=2p_k+2$, hence $p_k=3\cdot2^k-2$, namely $1,4,10,22,\ldots$. If $p_k\ge M+1$, the error and smooth-model residual are $O(x^{-M-1})$, within $O(E_M(x))$. This proves a statement about exact-real Newton iteration, not unqualified floating-point certification.

\subsection{The first crossing and strict integer brackets}
The forward equivalent gives $\log W_n/n\to-\lambda$. For any $0<\xi<\lambda$, all sufficiently large indices satisfy
\[
 e^{-(\lambda+\xi)n}\le W_n\le e^{-(\lambda-\xi)n}.
\]
The remaining finite prefix has a positive minimum. For $\varepsilon$ below that minimum, the two exponential envelopes place the first crossing between $L_\varepsilon/(\lambda+\xi)$ and $\lceil L_\varepsilon/(\lambda-\xi)\rceil$, up to bounded integer adjustments. Letting $\xi\downarrow0$ proves
\begin{equation}\label{ppg:all:eq:firstcrossing}
 N(\varepsilon)\sim L_\varepsilon/\lambda\sim x,
 \quad\text{so eventually }N(\varepsilon)>x/2.
\end{equation}
No monotonicity of the actual sequence was required.

The function $g_M(t)=(1+\log t)^M/t^{M+1}$ is decreasing for $t\ge1$, since
\[
 \frac{g'_M(t)}{g_M(t)}=
 \frac1t\left\{\frac M{1+\log t}-(M+1)\right\}<0.
\]
Consequently \eqref{ppg:all:eq:logW} implies a uniform bound
$|d_{M,n}|\le K'_ME_M(x)$ for every integer $n\ge x/2$ once $x$ is large. This uniformity holds over the entire unbounded tail without interpolating the remainder.

For either center $z$ above, choose $D_M$ so that
$|f_M(z)-L_\varepsilon|\le D_ME_M(x)$, and choose a strict margin
$C_M>2(D_M+K'_M)/\lambda$. Put $\Delta=C_ME_M(x)$. For $x/2\le n<z-\Delta$, the derivative lower bound implies
\[
 \log W_n\ge-L_\varepsilon+
   \{\lambda C_M/2-D_M-K'_M\}E_M(x)>-L_\varepsilon.
\]
Earlier crossings are excluded by \eqref{ppg:all:eq:firstcrossing}. At
$n_+=\lceil z+\Delta\rceil$ the reversed inequalities give
$\log W_{n_+}<-L_\varepsilon$. Therefore
$\lceil z-\Delta\rceil\le N(\varepsilon)\le\lceil z+\Delta\rceil$, proving Theorem~\ref{ppg:all:thm:threshold}, including integer endpoint coincidences.

\begin{remark}[The rounding boundary]
The uncertainty tends to zero before rounding. It does not imply
$N=z+o(1)$ or $N=\lceil z\rceil$ for every sufficiently small $\varepsilon$. Near an integer, an arbitrarily small unknown signed error can change a ceiling. Effective forward constants and remainder bounds would be needed to turn the existence bracket into a numerical threshold certificate.
\end{remark}

\begin{remark}[\wtag\ Repository context, 5 October 2026]\label{ppg:all:rem:inverse}
The leading root $x$ of \eqref{ppg:all:eq:x0} is the large root of $aX+b\log X=L$ with $a=\lambda$, $b=-\frac12$ and $L=L_\varepsilon+\log A$, which the branch rule of \texttt{p0:thm:lambert-core} \cite{tsvol} assigns to $W_{-1}$ (Remark~\ref{ppg:lead:rem:inverse}); the bracket \eqref{ppg:all:eq:ceiling} determines $N(\varepsilon)$ whenever its two ceilings agree, as in the separation condition of \texttt{p0:thm:staircase}; and the comparison with $f_M$ is the mean-value step of \texttt{p0:cor:forward-to-inverse}. The centres $z_M$ with $M\ge2$ carry polynomials $\mathcal R_j(\log x)$ because the forward coefficients $\mathcal L_j$ are polynomials in $\log n$; the volume's reversion theorems \texttt{p0:thm:lambert-centered} and \texttt{p0:thm:flattening} treat correction series with constant coefficients, so the recursion of Section~\ref{ppg:all:sec:inverse} is an extension of them, not an instance. With $M=1$ the theorem implies Part~\ref{ppg:lead:part}'s Theorem~\ref{ppg:lead:thm:threshold} (Remark~\ref{ppg:lead:rem:inverse}), and with $M=2$ Part~\ref{ppg:log:part}'s Theorem~\ref{ppg:log:thm:inverse} (Remark~\ref{ppg:log:rem:inverse}).
\end{remark}

\section{A computable refinement of the exponential rate}\label{ppg:all:sec:rate}
This section is independent of the all-orders bootstrap. It uses the following inputs proved in Report 144. Let
\[
 u_n=\frac{W_n}{\sqrt{n+1}},\qquad F(z)=\sum_{n\ge0}u_nz^n.
\]
Then $u_n>0$, $u_0=1$, the radius is $R=1/r_*$, and $F(z)\to\infty$ as $z\uparrow R$. There is an exact normalized kernel with
\begin{equation}\label{ppg:all:eq:alphainput}
 nu_n=\sum_{i+j=n-1}\alpha_{ij}u_iu_j,\qquad
 a\le\alpha_{ij}\le1,\qquad
 0\le\alpha_{ij}-a\le\frac1{i+1}+\frac1{j+1}.
\end{equation}
The lower bound used is the established $a=(2\pi)^{-1/2}$, not a weaker preliminary bound. One may start from the already certified rational interval
\begin{equation}\label{ppg:all:eq:initialR}
 [L_0,U_0]=[125/79,1000/527]
\end{equation}
for $R$. This initial certificate is an input and is not recomputed here.

\subsection{Exact real-axis enclosures}
Set $v=1/F$ and $H=F'/F^2$. Absolute interior convergence and \eqref{ppg:all:eq:alphainput} give $a\le H\le1$ and $v'=-H$. Boundary divergence gives $v(R-)=0$. Hence
\begin{equation}\label{ppg:all:eq:reciprocalbounds}
 a(R-z)\le v(z)=\int_z^R H(t)\,\dd t\le R-z,
 \qquad z+v(z)\le R\le z+v(z)/a.
\end{equation}
For $t>0$, put
\[
 J(t)=\sum_{n\ge0}\frac{u_nt^n}{n+1}=\frac1t\int_0^tF(s)\,\dd s.
\]
The inverse-index kernel bound implies
\[
 0\le H(t)-a\le2J(t)/F(t)
 \le\frac{2(R-t)}{at}\log\frac R{R-t}.
\]
Indeed \eqref{ppg:all:eq:reciprocalbounds} gives $F(s)\le1/(a(R-s))$ and $v(t)\le R-t$. For $0<z<R$, integration with $\delta=R-z$ yields
\begin{equation}\label{ppg:all:eq:defect}
 0\le v(z)-a\delta
 \le\frac{\delta^2}{az}\left(\log\frac R\delta+\frac12\right).
\end{equation}
Given $U\ge R$ and $D=U-z$, first replace $R$ by $U$ in the logarithm and then replace $\delta$ by $D$. The second replacement is valid because
$\frac{\dd}{\dd t}\{t^2(\log(U/t)+1/2)\}=2t\log(U/t)>0$ for $0<t<U$. Thus, with
\begin{equation}\label{ppg:all:eq:Crate}
 C(z,U)=\frac{D^2}z\left(\log\frac UD+\frac12\right),
\end{equation}
we obtain the refined exact enclosure
\begin{equation}\label{ppg:all:eq:refinedR}
 z+\frac{v(z)}a-\frac{C(z,U)}{a^2}\le R\le z+\frac{v(z)}a.
\end{equation}
Its width is $C/a^2$. Its possibly negative lower endpoint causes no problem after intersection with the old positive bracket.

\subsection{A terminating evaluator using finite coefficients}
Fix rational $0<z<t<R$ with a known strict-interiority witness. Since each $W_n$ is a computable rational, each $u_n$ can be enclosed by rational square-root bounds. For $N\ge1$, let $P_N(t)=\sum_{n<N}u_nt^n$. Sequentially compute a rational
\[
 P_N(t)\le\overline P_N\le P_N(t)+1,
\]
and test $4t\overline P_N<N$. If this rational strict test fails, including equality, advance immediately to $N+1$. Each finite computation terminates, and the test succeeds for every $N>4t(F(t)+1)$, since $P_N(t)\le F(t)<\infty$. No equality test on an unknown real is left unresolved.

For the accepted $N$, the nonnegative formal solution
$B=P_N+(x/N)B^2$, with $P_N(x)=\sum_{n<N}u_nx^n$, dominates $F$ coefficientwise. Below degree $N$ this follows from the polynomial term. At degree $n\ge N$, \eqref{ppg:all:eq:alphainput} gives
$u_n\le N^{-1}\sum_{i+j=n-1}u_iu_j$, and the factor $x$ makes coefficient induction valid. The Catalan expansion converges at $t$ and selects the smaller quadratic root
\[
 B_N(t)=\frac{2P_N(t)}{1+\sqrt{1-4tP_N(t)/N}}
        \le2P_N(t)\le2\overline P_N=:b.
\]
Both numerical roots are positive; selecting this formal branch is essential. Positivity now gives the tail bound
\begin{equation}\label{ppg:all:eq:tailF}
 \sum_{k\ge M}u_kz^k\le b(z/t)^M.
\end{equation}
Given positive rational tolerance $\epsilon$, find $M$ by rational arithmetic so that the last quantity is at most $\epsilon/2$, and enclose the finite prefix to width at most $\epsilon/2$. Adding the tail only to the upper endpoint gives a rational interval of width at most $\epsilon$ for $F(z)$. Its lower endpoint may be replaced by its maximum with 1. This proves arbitrary-accuracy evaluation with finitely many coefficients, but makes no practical coefficient-count or runtime claim.

If $1\le f_-\le F(z)\le f_+$, then
\begin{equation}\label{ppg:all:eq:vinterval}
 v_-=1/f_+\le v(z)\le1/f_-=:v_+,
 \qquad v_+-v_-\le f_+-f_-.
\end{equation}
The reciprocal endpoints reverse order. The elementary constants also have computable rational enclosures: alternating arctangent sums enclose $\pi$, positive square roots admit rational bisection, and for rational $x>1$, $y=(x-1)/(x+1)$ gives
\begin{equation}\label{ppg:all:eq:loginterval}
 \log x=2\sum_{j=0}^{m-1}\frac{y^{2j+1}}{2j+1}+T_m,
 \qquad 0\le T_m\le\frac{2y^{2m+1}}{(2m+1)(1-y^2)}.
\end{equation}
These examples establish computability, not computational efficiency.

\subsection{Outward rounding and interval contraction}
For rational enclosures $0<a_-\le a\le a_+$, $0<v_-\le v\le v_+$, and $0<C_-\le C\le C_+$, define endpoint intervals for
$A=z+v$, $B=z+v/a$, and $P=z+v/a-C/a^2$ by
\begin{align}
 A_-&=z+v_-,&A_+&=z+v_+,\notag\\
 B_-&=z+v_-/a_+,&B_+&=z+v_+/a_-,\label{ppg:all:eq:roundedendpoints}\\
 P_-&=z+v_-/a_+-C_+/a_-^2,
 &P_+&=z+v_+/a_--C_-/a_+^2.\notag
\end{align}
All three enclosures are valid, and their widths tend to zero as inputs are refined. This subtraction form avoids the invalid rule of dividing a possibly negative lower numerator by an upper positive denominator.

\begin{theorem}[Computable nested rate enclosures]\label{ppg:all:thm:rate}
Start from any known rational $0<L_0\le R\le U_0$. There is a terminating finite-coefficient algorithm producing nested rational intervals $[L_k,U_k]$ for $R$, to any prescribed positive rational width tolerance. If $w_k=U_k-L_k>0$, each update satisfies
\begin{align}
 w_{k+1}&\le\frac{313}{400}w_k<\frac45w_k,\label{ppg:all:eq:coarsecontract}\\
 w_{k+1}&\le\frac{25}{8a^2L_0}w_k^2
       \left(\log\frac{U_0}{w_k}+\frac12\right)+\frac{w_k^2}{50}.
 \label{ppg:all:eq:supercontract}
\end{align}
The corresponding rate intervals are $[1/U_k,1/L_k]$, of width at most $w_k/L_0^2$. The second estimate is a quadratic-times-log width bound. Neither estimate is a bit-complexity bound or a claim that new digits have been computed.
\end{theorem}
\begin{proof}
For the current bracket $[L,U]$, stop if $w=U-L=0$. Otherwise choose rational
\begin{equation}\label{ppg:all:eq:interior}
 d=\min(w/4,L/2),\qquad z=L-d,\qquad t=L-d/2,
 \qquad\eta=\min(w,w^2)/100.
\end{equation}
Then $0<z<t<L\le R$, $z\ge L/2$, and $0<d/w\le1/4$. The strict witness for the evaluator is available even when $L=R$. Compute and refine the input intervals until each endpoint interval in \eqref{ppg:all:eq:roundedendpoints} has width at most $\eta$. The evaluator and elementary interval operations terminate; all endpoint widths tend to zero, so this refinement also terminates. Define
\begin{equation}\label{ppg:all:eq:newLU}
 L_{\rm new}=\max(L,A_-,P_-),\qquad
 U_{\rm new}=\min(U,B_+).
\end{equation}
The maximum consists of valid lower bounds, and the minimum of valid upper bounds. The new interval is nonempty, rational, nested, and contains $R$.

To bound its width, set $s=d/w$ and $T=v(z)/w$. In units of the old width, the exact coarse intersection in \eqref{ppg:all:eq:reciprocalbounds} has endpoints
$\max(0,T-s)$ and $\min(1,T/a-s)$. If $T\le a(1+s)$, its width is at most
$T(1/a-1)\le(1+s)(1-a)$. If $T\ge a(1+s)$, its width is at most
$1+s-T\le(1+s)(1-a)$. Thus its width is at most $q_0w$, where
$q_0=(5/4)(1-a)$. Outward errors add at most $2\eta$; the refined lower bound can only decrease the interval. The elementary bound $\pi<22/7$ implies $a>39/100$, so
\[
 w_{\rm new}\le q_0w+2\eta
 \le\left(\frac{61}{80}+\frac1{50}\right)w
 =\frac{313}{400}w<\frac45w.
\]
For the refined intersection, its exact width is $C/a^2$, hence
$w_{\rm new}\le C/a^2+2\eta$. Throughout iteration,
$z\ge L_0/2$, $U\le U_0$, and $w\le D=U-z=w+d\le5w/4$. Therefore
\[
 \frac C{a^2}\le\frac{25}{8a^2L_0}w^2
       \left(\log\frac{U_0}w+\frac12\right),\qquad
 2\eta\le w^2/50,
\]
which proves \eqref{ppg:all:eq:supercontract}. All factors are positive since $w<U\le U_0$. Its right side divided by $w$ tends to zero.

For a radius-width tolerance $\epsilon>0$, find an integer $K$ by rational multiplication until $(4/5)^Kw_0\le\epsilon$. Each of the $K$ updates terminates and \eqref{ppg:all:eq:coarsecontract} proves the stopping guarantee. Finally,
\[
 \frac1{L_k}-\frac1{U_k}=\frac{w_k}{L_kU_k}\le\frac{w_k}{L_0^2}.
\]
For rate-width tolerance $\epsilon$, use $(4/5)^Kw_0\le\epsilon L_0^2$, or stop sooner when the actual rational width is sufficient.
\end{proof}
For a wholly rational coefficient in \eqref{ppg:all:eq:supercontract}, one may replace $25/(8a^2L_0)$ by the larger $31250/(1521L_0)$. No run of this interval-refinement algorithm, beyond its symbolic finite checks, is claimed in the present report.
\begin{writenote}[5 October 2026]
The inputs \eqref{ppg:all:eq:alphainput} are Part~\ref{ppg:lead:part}'s \eqref{ppg:lead:eq:normalized}, \eqref{ppg:lead:eq:kbound}, \eqref{ppg:lead:eq:kerror} and Theorem~\ref{ppg:lead:thm:stronglower}; the divergence of $F$ at $R$ is Lemma~\ref{ppg:lead:lem:divergence}; the Catalan majorant is that of Theorem~\ref{ppg:lead:thm:uppercert}; and \eqref{ppg:all:eq:initialR} is \eqref{ppg:lead:eq:numeric} written for $R$. The theorem quantifies Theorem~\ref{ppg:lead:thm:uppercert}, which has no rate.
\end{writenote}

\section{Limits and further questions}\label{ppg:all:sec:limits}
The main advance over the foundation reports is an all-fixed-order coefficient theorem with specified logarithmic degree and remainder, obtained without circular endpoint moments. The explicit third order also shows that the generic degree bound need not be attained: the squared logarithm cancels. Several questions remain outside the proved conclusions.
\begin{enumerate}
\item \textbf{Effective correction constants.} The endpoint-series formulas prove existence and identify coefficients, but do not provide effective tail bounds or onset indices for $B_2$, $B_3$, or the higher correction constants. The independent rate algorithm does not automatically supply these bounds.
\item \textbf{Higher cancellations and coefficient structure.} The upper bound $\deg P_j\le j-1$ allows further cancellations. A systematic description of their exact degrees, signs, and dependence on the scalar moments would sharpen the formal structure.
\item \textbf{Growth with the order.} No estimate on the size of $P_j$ as $j\to\infty$ is proved. Convergence, divergence rates, optimal truncation, exponentially small sectors, and summability of an infinite expansion all require additional arguments.
\item \textbf{Practical certified numerics.} The rate-enclosure theorem has a terminating evaluator and explicit outer width contraction, but no useful bound on total arithmetic work or the number of probability coefficients. Implementing and benchmarking a certified enclosure computation is a separate task; no new rate digits are certified here.
\item \textbf{Finite-index behavior and numerical thresholds.} All-index ratio monotonicity, effective onsets of eventual conclusions, and numerical first-crossing certificates remain separate questions. The ceiling brackets cannot be replaced by an unrounded integer asymptotic equality.
\end{enumerate}
These limitations are compatible: computability of the leading rate by an independent method does not make the connection constants or the all-order remainders effective.
\begin{writenote}[5 October 2026]
All five items remain open. Section~\ref{ppg:sec:further} collects them with the questions of Parts~\ref{ppg:lead:part} and~\ref{ppg:log:part}: item~1 is its question~\ref{ppg:fq:effective}, item~2 question~\ref{ppg:fq:degrees}, item~3 question~\ref{ppg:fq:growth}, item~4 question~\ref{ppg:fq:numerics}, and item~5 questions~\ref{ppg:fq:monotone} and~\ref{ppg:fq:thresholds}. It adds the value of $S_A$ (question~\ref{ppg:fq:sa}).
\end{writenote}

\section{Exact checks and reproduction}\label{ppg:all:sec:reproduction}
The accompanying archive contains this PDF and editable LaTeX source, a new standard-library exact-arithmetic companion, deterministic build scripts, a source and toolchain record, and unchanged source/PDF copies of Reports 144 and 145. It contains no primary copyrighted manuscript PDF, prior computational archive, or development intermediates.

The companion recomputes finite rational identities for logarithmic coefficient models, signed discrete Taylor expansion, the critical inverse, kernel and third-order constants, and finite inverse algebra. Its closed fixtures are checked semantically, rather than accepted because a checksum matches. Normal and optimized Python tests exercise exact algebra and structural rejection behavior. The companion documentation states the precise finite scope and separates analytic proof inputs from recomputed rational identities.

These tests are safeguards for algebra, signs, indexing, and reproducibility. They do not prove the uniform convolution bounds, infinite moment convergence, the all-orders quantifier, or analytic error propagation; those are established in the preceding proofs. They do not evaluate the connection constants, execute a new rate-enclosure iteration, certify decimal digits, or certify an integer threshold. The interval-theorem checks concern finite identities, not computational-complexity claims.

The README gives exact commands for verification and rebuilding. The PDF build uses a clean local TeX configuration, fixed timestamp and locale, no shell escape, and stabilized references; it rejects warnings and box defects. With the recorded TeX/font toolchain it is byte-reproducible. The archive uses an explicit fixed allowlist, fixed order, timestamps, and permissions, and a generated hash manifest. Output guards reject overwrites, symlinks, and parent-traversal paths. The foundation documents are archival dependencies and are not rebuilt or changed by the new PDF build.
\begin{writenote}[5 October 2026]
Shipped names: the checker and its tests are \file{code/146-allorders-companion-checks.py} and \file{code/146-allorders-companion-test_checks.py}, the build scripts \file{code/146-allorders-build_pdf.py} and \file{code/146-allorders-build_archive.py}, the records \file{data/146-allorders-companion-*.json} and \file{data/146-allorders-toolchain.json}, and \file{146-allorders-SOURCES.md} and \file{146-allorders-companion-README.md} at the root. The delivered build script compiles \file{Report146.tex}, not this merged \file{article.tex}, which is built as the README of this report says. The foundation copies, the PDF, the delivered README (replaced by the report's) and \file{SHA256SUMS} are not shipped. The checker refuses to run without POSIX no-follow file operations; the README gives a rerun route.
\end{writenote}

\section{Further questions and research}\label{ppg:sec:further}
\wtag\ (5 October 2026.) This section collects, under Vladimir's standing rule of 4~October 2026, every question that the three Parts leave open, deduplicated, with its source, its status, the evidence available and what is missing. Questions that a later Part answers are not repeated as open: Part~\ref{ppg:lead:part}'s ``pointwise corrections and logarithms'' and Part~\ref{ppg:log:part}'s ``full logarithmic expansion'' are answered by Theorem~\ref{ppg:all:thm:all}, Part~\ref{ppg:log:part}'s first ``finer ratio term'' by Corollary~\ref{ppg:all:cor:ratio}, and the eventual monotonicity of the ratios by Corollary~\ref{ppg:log:cor:monotonicity}; the dated notes at the questions say so. No claim of the three manuscripts was found to be wrong. The wrong statements on record are Serra's, each with its proof: his prefactor $\sqrt{2\pi}$ and its approach from below (Remark~\ref{ppg:lead:rem:serra}), his integral equation~(200) read through the profile (Proposition~\ref{ppg:lead:prop:op3}), his expansion~(199) for $p\ge2$ (Remark~\ref{ppg:all:rem:op2}) and the slip in question~\ref{ppg:fq:monotone} below.
\begin{enumerate}[label=\arabic*.,ref=\arabic*]
\item\label{ppg:fq:effective} \emph{Effective constants and remainders.} Sources: Part~\ref{ppg:lead:part}, Section~\ref{ppg:lead:sec:questions} (``Quantifying the coefficient bridge'', answered only asymptotically by Parts~\ref{ppg:log:part} and~\ref{ppg:all:part}); Part~\ref{ppg:log:part}, Section~\ref{ppg:log:sec:open}, items~2 and~4; Part~\ref{ppg:all:part}, Section~\ref{ppg:all:sec:limits}, item~1. The constants $B_2$, $B_3$, $K_x$ and the higher coefficients are given by absolutely convergent series in the exact $x_i$ and $y_i$ (\eqref{ppg:log:eq:B2}, Section~\ref{ppg:all:sec:third}); the remainders of Theorems~\ref{ppg:log:thm:full} and~\ref{ppg:all:thm:all} carry no explicit constants or onsets. Missing: computable tail bounds for those series and explicit constants in the uniform kernel, convolution and inverse estimates. Evidence, uncertified: the write's fit of $\log x_n$ (question~\ref{ppg:fq:sa}) gives $B_2=-0.0573920672\ldots$, agreeing with the placement dossier's float64 evaluation $-0.0573921$ of Part~\ref{ppg:log:part}'s series (its direct estimate at $n=4096$, $-0.0573727$, was still drifting).
\item\label{ppg:fq:monotone} \emph{Strict decrease of $W_{n+1}/W_n$ at every index, and an onset.} Sources: Serra's Open Problem~1; Part~\ref{ppg:lead:part}, Section~\ref{ppg:lead:sec:questions}; Part~\ref{ppg:log:part}, Section~\ref{ppg:log:sec:open}, item~3; Part~\ref{ppg:all:part}, Section~\ref{ppg:all:sec:limits}, item~5. Proved: eventual strict decrease and eventual strict log-concavity (Corollary~\ref{ppg:log:cor:monotonicity}). Evidence: $\rho_0>\rho_1>\cdots>\rho_{3999}$ in the write's 60-digit computation, the same for $n\le4095$ in the dossier's float64 run, and for $n\le1276$ according to Serra (his Remark~8.1). Missing: an explicit onset $n_0$ for Corollary~\ref{ppg:log:cor:monotonicity} together with a certified finite check below it, or a structural argument such as total positivity. \emph{A slip in the source:} Serra's suggested approach (printed page~68) calls his inequality~(198), $W(n+2)/W(n+1)<W(n+1)/W(n)$, equivalent to log-convexity, $W(n+1)^2<W(n)W(n+2)$. Since every $W(n)$ is positive, multiplying~(198) by $W(n+1)W(n)$ shows that it is equivalent to $W(n)W(n+2)<W(n+1)^2$, strict log-concavity; log-convexity would make the ratios increase.
\item\label{ppg:fq:degrees} \emph{Exact logarithmic degrees.} Source: Part~\ref{ppg:all:part}, Section~\ref{ppg:all:sec:limits}, item~2. Theorem~\ref{ppg:all:thm:all} proves $\deg P_j\le j-1$; the bound is attained for $j=1,2$ and not for $j=3$, where the squared logarithm cancels (Theorem~\ref{ppg:all:thm:third}). Missing: the exact degrees, signs and dependence on the scalar moments for $j\ge4$.
\item\label{ppg:fq:growth} \emph{Growth with the order, and a transseries.} Sources: Part~\ref{ppg:lead:part}, Section~\ref{ppg:lead:sec:questions} (``No full transseries is obtained''); Part~\ref{ppg:all:part}, Section~\ref{ppg:all:sec:limits}, item~3. Nothing is known about the size of $P_j$ as $j\to\infty$, the convergence or divergence of the expansion, optimal truncation, exponentially small terms, Borel summability, or the continuation of $X(z)$ across $|z|=1$. All three Parts prove their expansions without continuation.
\item\label{ppg:fq:numerics} \emph{Practical certified digits of $r_*$.} Sources: Part~\ref{ppg:lead:part}, Section~\ref{ppg:lead:sec:questions} (``Efficient certified digits''); Part~\ref{ppg:all:part}, Section~\ref{ppg:all:sec:limits}, item~4. Proved: computability (Theorem~\ref{ppg:lead:thm:uppercert}) and a width contraction by $313/400$ per update (Theorem~\ref{ppg:all:thm:rate}), with no bound on the work. Missing: an implementation, a work bound and a run. The only certified digits remain those of $527/1000<r_*<79/125$; Serra's eight digits rest, by his own Remark~6.25, on the unproved expansion~(151), and OEIS's 18 digits are not certified. Evidence: in the write's computation, a free parameter for the relative error of the 18-digit value of $R$ is fitted as $1.146\times10^{-19}$ both from $x_n$ and from the partial sums of $S_A$, consistent with truncation in the last digit.
\item\label{ppg:fq:thresholds} \emph{Numerical threshold certificates.} Sources: Part~\ref{ppg:log:part}, Section~\ref{ppg:log:sec:open}, item~5; Part~\ref{ppg:all:part}, Section~\ref{ppg:all:sec:limits}, item~5, and Remark following Theorem~\ref{ppg:all:thm:threshold}. The brackets of Theorems~\ref{ppg:lead:thm:threshold}, \ref{ppg:log:thm:inverse} and~\ref{ppg:all:thm:threshold} are existence statements. Missing: effective forward constants and remainders (question~\ref{ppg:fq:effective}); even then, rounding near an integer needs a separation margin.
\item\label{ppg:fq:policies} \emph{Other policies, optimality, universality.} Source: Part~\ref{ppg:lead:part}, Section~\ref{ppg:lead:sec:questions}. Which positive size-dependent convolution kernels admit the balanced-merging rate argument, the pole mechanism and the Riccati comparison, and how changes in the reciprocal-index correction change the logarithmic terms. Serra's Open Problems~5--7 (the optimal strategy) and his questions on the arithmetic nature of $r_\infty$ (his Section~7: transcendence, non-D-finiteness of $\sum W(n)z^n$) are not addressed by any Part.
\item\label{ppg:fq:sa} \emph{Is $S_A=\frac1{12}$?} Source: placement dossier of batch~105 (5~October 2026; float64, about 8 digits); checked further by the write. Here
\[
 S_A=\sum_{i\ge0}\bigl(b_ix_i-1\bigr)=\sum_{i\ge0}\bigl(q_iW_iR^{i+1}-1\bigr),\qquad
 q_i=\frac1{i!}\int_i^{i+1}s^ie^{-s}\,\dd s,
\]
the constant \eqref{ppg:all:eq:SA} of Part~\ref{ppg:all:part} (the second form is the identity in the note there). Neither Part~\ref{ppg:all:part} nor any other Part states or proves a value. \emph{Evidence (uncertified).} The write computed $x_0,\ldots,x_{4000}$ from the exact factorization of Lemma~\ref{ppg:all:lem:factor}, evaluating each kernel integral by 24-point Gauss--Legendre quadrature (the integrand is a polynomial times $e^{-1}$, and the quadrature error is far below $10^{-40}$) and running $(n+1)x_{n+1}=\sum\beta_{ij}x_ix_j$ in fixed-point arithmetic with resolution $2^{-200}$, from $x_0=1/A$ with the 18-digit OEIS value of $r_*$. The computation reproduces the exact $W_2,\ldots,W_5$ of \eqref{ppg:lead:eq:prefix} to $10^{-59}$. A least-squares fit of the partial sums $\sum_{i\le N}\bar a_i$, $300\le N\le4000$, by $S_A$ plus the tail terms $(\log N)^k/N^{j-1}$, $0\le k\le j-1$, $2\le j\le5$, plus the term by which an error in $r_*$ enters, gives $S_A-\frac1{12}=-1.9\times10^{-19}$ for the window $600\le N\le4000$ and at most $3.8\times10^{-18}$ in absolute value for the windows $[300,4000]$ and $[300,2000]$ (with only four tail orders, at most $4\times10^{-14}$); the same fit returns the coefficient $\frac7{540}$ of $\log N/N$ to twelve digits. The fit of $\log x_n$ reported after Theorem~\ref{ppg:all:thm:third} is consistent with the third-order logarithmic coefficient that $S_A=\frac1{12}$ predicts. \emph{Consequence if true:} the third-order logarithmic coefficients are rational, $\frac{199}{15120}$ for $x_n$ and $\frac{101}{15120}$ for $W_n/(A\sqrt n\,r_*^n)$. \emph{Missing:} a proof, presumably an exact identity coming from the boundary structure of the recurrence (each $q_i$ is the limit of $Q_{n+1,i+1}$ as $n\to\infty$ with $i$ fixed), and an explanation of why $S_A$, unlike $B_2$, appears to be rational.
\end{enumerate}

\begin{thebibliography}{9}
\bibitem{serra}
Alessandro Serra, \emph{The Number Challenge: Asymptotic Analysis and Hidden Constant of a Recursive Placement Game}, version~1, manuscript dated July~22, 2026, published August~15, 2026. Zenodo record 21946055.
\href{https://doi.org/10.5281/zenodo.21946055}{doi:10.5281/zenodo.21946055}. Inspected version-1 PDF (92 pages), with printed page numbering one less than its PDF page numbering. Model and recurrence: Theorem~3.12, equations (47)--(48), printed pages~23--24, as documented in Parts~I and~II. \wtag\ Licensed CC~BY~4.0 according to the Zenodo record; read by the write on 5~October 2026, which checked the page and equation numbers cited in all three Parts.

\bibitem{oeis}
The OEIS Foundation Inc., \emph{The On-Line Encyclopedia of Integer Sequences}, entry A398540, inspected October~3, 2026. The association and source scope are documented in Parts~I and~II.
\url{https://oeis.org/A398540}. \wtag\ Fetched again on 5~October 2026; data licensed CC~BY-SA~4.0.

\bibitem{robbins}
Herbert Robbins, A remark on Stirling's formula, \emph{The American Mathematical Monthly} \textbf{62} (1955), no.~1, 26--29.
\href{https://doi.org/10.2307/2308012}{doi:10.2307/2308012}.

\bibitem{r144}
\emph{Asymptotics and a computable exponential rate for a proportional placement game}, Report 144, October 3, 2026. Unchanged LaTeX source and PDF included in the \texttt{foundation} directory of the archives of Reports~145 and~146. See Section~4 for uniform kernel bounds, Section~9 for the leading equivalent and ratio limit, and Section~11 for the earlier threshold inverse. Supplies the leading equivalent, kernel bounds, radius properties, initial certified rate bracket, and positive-prefix inverse argument. \wtag\ Printed as Part~\ref{ppg:lead:part} of this report (cited as ``foundation'' by Part~\ref{ppg:log:part}).

\bibitem{r145}
\emph{Pointwise corrections for the proportional placement game}, Report 145, October 3, 2026. Unchanged source and PDF included in the archive of Report~146. Supplies the starting rate, exact critical inverse with its boundary condition, and the absolutely convergent-series definition of $B_2$. \wtag\ Printed as Part~\ref{ppg:log:part} of this report.

\bibitem{tsvol}
\raggedright\wtag\ ProveIt, \emph{Transseries and inversion}, the transseries-and-inversion volume (\texttt{p0:thm:lambert-core}, \texttt{p0:def:three-inverses}, \texttt{p0:thm:staircase}, \texttt{p0:cor:forward-to-inverse}, \texttt{p0:thm:lambert-centered}, \texttt{p0:thm:flattening}), \file{Analysis/Transseries/docs/series-and-transseries/Transseries_And_Inversion/transseries_and_inversion.tex}.
\end{thebibliography}
\end{document}
